\section*{Anhang: Medizinisches Mappings-Register}
\addcontentsline{toc}{section}{Anhang: Medizinisches Mappings-Register}
Das folgende Mappings-Register zeigt die Normalisierungsregeln von Freitext-Angaben, Experten-Ground-Truths und Modell-Vorhersagen.\\[1em]
\begin{longtable}{|p{3.5cm}|p{5cm}|p{6.5cm}|}
\hline
\textbf{Kategorie} & \textbf{Eingabe (Freitext/Synonym)} & \textbf{Zielbegriff (Normalisiert)} \\ \hline
\endhead
ANTIMIKROBIELL\_GT\_MAPPING & \texttt{Silber (Ag⁺)} & Silber (Ag+) \\ \hline
ANTIMIKROBIELL\_GT\_MAPPING & \texttt{Silber (Ag+)} & Silber (Ag+) \\ \hline
ANTIMIKROBIELL\_GT\_MAPPING & \texttt{PHMB} & PHMB \\ \hline
ANTIMIKROBIELL\_GT\_MAPPING & \texttt{Octenidin} & Octenidin \\ \hline
ANTIMIKROBIELL\_GT\_MAPPING & \texttt{Cadexomer-Iod} & Cadexomer-Iod \\ \hline
ANTIMIKROBIELL\_GT\_MAPPING & \texttt{Honig (Medihoney)} & Honig (Medihoney) \\ \hline
DEBRIDEMENT\_GT\_MAPPING & \texttt{Autolytisch (Hydrogele, Hydrokolloide, Folienverbände)} & Autolytisch (Hydrogele, Hydrokolloide, Folienverbände) \\ \hline
DEBRIDEMENT\_GT\_MAPPING & \texttt{Mechanisch (Monofilament-Pad, feuchte Kompressen, Wundspülung)} & Mechanisch (Monofilament-Pad, feuchte Kompressen, Wundspülung) \\ \hline
DEBRIDEMENT\_GT\_MAPPING & \texttt{Chirurgisch/Scharf (Skalpell, Kürette)} & Chirurgisch/Scharf (Skalpell, Kürette) \\ \hline
DEBRIDEMENT\_GT\_MAPPING & \texttt{Enzymatisch (Kollagenase)} & Enzymatisch (Kollagenase) \\ \hline
DEBRIDEMENT\_GT\_MAPPING & \texttt{Kein Débridement erforderlich} & Kein Débridement erforderlich \\ \hline
DEBRIDEMENT\_GT\_MAPPING & \texttt{Autolytisch (Hydrogele, Hydrokolloide, Folienverbände), vorsichtige mechanische Reinigung lockerer Beläge} & Autolytisch (Hydrogele, Hydrokolloide, Folienverbände), Mechanisch (Monofilament-Pad, feuchte Kompressen, Wundspülung) \\ \hline
DEBRIDEMENT\_GT\_MAPPING & \texttt{Autolytisch (Hydrogele, Hydrokolloide, Folienverbände), vorsichtig mechanisches Debridement} & Autolytisch (Hydrogele, Hydrokolloide, Folienverbände), Mechanisch (Monofilament-Pad, feuchte Kompressen, Wundspülung) \\ \hline
DEBRIDEMENT\_GT\_MAPPING & \texttt{vorsichtig mechanisches Debridement} & Mechanisch (Monofilament-Pad, feuchte Kompressen, Wundspülung) \\ \hline
DEBRIDEMENT\_GT\_MAPPING & \texttt{vorsichtige mechanische Reinigung lockerer Beläge} & Mechanisch (Monofilament-Pad, feuchte Kompressen, Wundspülung) \\ \hline
DEBRIDEMENT\_GT\_MAPPING & \texttt{nicht aggressiv entfernen, solange keine Revaskularisation erfolgt ist, Chirurgisch/Scharf (Skalpell, Kürette)} & Chirurgisch/Scharf (Skalpell, Kürette) \\ \hline
DEBRIDEMENT\_GT\_MAPPING & \texttt{Autolytisch (Hydrogele, Hydrokolloide, Folienverbände), chirurgisches Debridement erst nach Perfusionsverbesserung, sofern möglich} & Autolytisch (Hydrogele, Hydrokolloide, Folienverbände), Chirurgisch/Scharf (Skalpell, Kürette) \\ \hline
DEBRIDEMENT\_GT\_MAPPING & \texttt{chirurgisches Debridement erst nach Perfusionsverbesserung, sofern möglich} & Chirurgisch/Scharf (Skalpell, Kürette) \\ \hline
DEBRIDEMENT\_GT\_MAPPING & \texttt{Chirurgisch/Scharf (Skalpell, Kürette), autolytisches Debridement, enzymatisches Debridement - ergänzend} & Autolytisch (Hydrogele, Hydrokolloide, Folienverbände), Chirurgisch/Scharf (Skalpell, Kürette), Enzymatisch (Kollagenase) \\ \hline
DEBRIDEMENT\_GT\_MAPPING & \texttt{autolytisches Debridement, enzymatisches Debridement - ergänzend} & Autolytisch (Hydrogele, Hydrokolloide, Folienverbände), Enzymatisch (Kollagenase) \\ \hline
DEBRIDEMENT\_GT\_MAPPING & \texttt{konservatives Abwarten bei unklarer Demarkation, selektives chirurgisches Débridement, enzymatisch/autolytisch ergänzend} & Autolytisch (Hydrogele, Hydrokolloide, Folienverbände), Chirurgisch/Scharf (Skalpell, Kürette), Enzymatisch (Kollagenase) \\ \hline
DEBRIDEMENT\_GT\_MAPPING & \texttt{kein aggressives Debridement ohne Gefäßstatus, trockene stabile Nekrosen ggf. belassen, Bei infizierten/fibrinösen Belägen: zurückhaltendes scharfes Debridement, autolytische Verfahren möglich} & Autolytisch (Hydrogele, Hydrokolloide, Folienverbände), Chirurgisch/Scharf (Skalpell, Kürette) \\ \hline
DEBRIDEMENT\_GT\_MAPPING & \texttt{Autolytisch (Hydrogele, Folienverbände), Hydrokolloide} & Autolytisch (Hydrogele, Hydrokolloide, Folienverbände) \\ \hline
DEBRIDEMENT\_GT\_MAPPING & \texttt{Autolytisch (Hydrogele, Folienverbände), Hydrokolloide, Mechanisch (Monofilament-Pad, Wundspülung), feuchte Kompressen} & Autolytisch (Hydrogele, Hydrokolloide, Folienverbände), Mechanisch (Monofilament-Pad, feuchte Kompressen, Wundspülung) \\ \hline
DEBRIDEMENT\_GT\_MAPPING & \texttt{Autolytisch (Hydrogele, Chirurgisch/Scharf (Skalpell, Folienverbände), Hydrokolloide, Kürette), Mechanisch (Monofilament-Pad, Wundspülung), feuchte Kompressen} & Autolytisch (Hydrogele, Hydrokolloide, Folienverbände), Mechanisch (Monofilament-Pad, feuchte Kompressen, Wundspülung), Chirurgisch/Scharf (Skalpell, Kürette) \\ \hline
DEBRIDEMENT\_GT\_MAPPING & \texttt{Chirurgisch/Scharf (Skalpell, Kürette), Mechanisch (Monofilament-Pad, Wundspülung), feuchte Kompressen} & Mechanisch (Monofilament-Pad, feuchte Kompressen, Wundspülung), Chirurgisch/Scharf (Skalpell, Kürette) \\ \hline
DEBRIDEMENT\_GT\_MAPPING & \texttt{Autolytisch (Hydrogele, Hydrokolloide, Folienverbände), Chirurgisch/Scharf (Skalpell, Kürette), Enzymatisch} & Autolytisch (Hydrogele, Hydrokolloide, Folienverbände), Chirurgisch/Scharf (Skalpell, Kürette), Enzymatisch (Kollagenase) \\ \hline
DEBRIDEMENT\_GT\_MAPPING & \texttt{Mechanisch (Monofilament-Pad, Wundspülung), feuchte Kompressen} & Mechanisch (Monofilament-Pad, feuchte Kompressen, Wundspülung) \\ \hline
DEBRIDEMENT\_GT\_MAPPING & \texttt{Autolytisch (Hydrogele, Chirurgisch/Scharf (Skalpell, Folienverbände), Hydrokolloide, Kürette)} & Autolytisch (Hydrogele, Hydrokolloide, Folienverbände), Chirurgisch/Scharf (Skalpell, Kürette) \\ \hline
EXSUDAT\_GT\_MAPPING & \texttt{Annahme stark} & Stark \\ \hline
EXSUDAT\_GT\_MAPPING & \texttt{Große Wunde: stark, da Mazerationen am Wundrand
Kleine Nekrose: trocken} & Stark \\ \hline
EXSUDAT\_GT\_MAPPING & \texttt{Keine} & Keine \\ \hline
EXSUDAT\_GT\_MAPPING & \texttt{Leicht} & Leicht \\ \hline
EXSUDAT\_GT\_MAPPING & \texttt{Mässig} & Mäßig \\ \hline
EXSUDAT\_GT\_MAPPING & \texttt{Mäßig} & Mäßig \\ \hline
EXSUDAT\_GT\_MAPPING & \texttt{Nekrosen: kein Exsudat. Unterschenkel mäßig} & Mäßig \\ \hline
EXSUDAT\_GT\_MAPPING & \texttt{Stark} & Stark \\ \hline
EXSUDAT\_GT\_MAPPING & \texttt{blutig bis serös teilweise vorhanden dann eher mäßig} & Mäßig \\ \hline
EXSUDAT\_GT\_MAPPING & \texttt{eher schwach b is mäßig} & Leicht, Mäßig \\ \hline
EXSUDAT\_GT\_MAPPING & \texttt{eher wenig} & Leicht \\ \hline
EXSUDAT\_GT\_MAPPING & \texttt{eher weniger} & Leicht \\ \hline
EXSUDAT\_GT\_MAPPING & \texttt{gering bis gar nicht} & Keine, Leicht \\ \hline
EXSUDAT\_GT\_MAPPING & \texttt{kein} & Keine \\ \hline
EXSUDAT\_GT\_MAPPING & \texttt{keine} & Keine \\ \hline
EXSUDAT\_GT\_MAPPING & \texttt{keine Angabe möglich} & Enthaltung / keine Angabe \\ \hline
EXSUDAT\_GT\_MAPPING & \texttt{keine Angabe möglich / eher wenig} & Leicht \\ \hline
EXSUDAT\_GT\_MAPPING & \texttt{keine Angabe möglich / vermutlich mässig} & Mäßig \\ \hline
EXSUDAT\_GT\_MAPPING & \texttt{keine Angabe möglich / vermutlich wenig} & Leicht \\ \hline
EXSUDAT\_GT\_MAPPING & \texttt{keine Angabe möglich, vermutlich gering} & Leicht \\ \hline
EXSUDAT\_GT\_MAPPING & \texttt{keine Angabe möglich, vermutlich trocken} & Keine \\ \hline
EXSUDAT\_GT\_MAPPING & \texttt{keine Einschätzung möglich} & Enthaltung / keine Angabe \\ \hline
EXSUDAT\_GT\_MAPPING & \texttt{keine genaue Angabe möglich, sieht eher trocken bis leicht exsudierend aus} & Keine, Leicht \\ \hline
EXSUDAT\_GT\_MAPPING & \texttt{könnte eher schwach sein} & Leicht \\ \hline
EXSUDAT\_GT\_MAPPING & \texttt{leicht} & Leicht \\ \hline
EXSUDAT\_GT\_MAPPING & \texttt{leicht bis mittel} & Leicht, Mäßig \\ \hline
EXSUDAT\_GT\_MAPPING & \texttt{leicht bis mäßig} & Leicht, Mäßig \\ \hline
EXSUDAT\_GT\_MAPPING & \texttt{leichte bis mäßige exsudation} & Leicht, Mäßig \\ \hline
EXSUDAT\_GT\_MAPPING & \texttt{mittel} & Mäßig \\ \hline
EXSUDAT\_GT\_MAPPING & \texttt{mittel bis stark} & Mäßig, Stark \\ \hline
EXSUDAT\_GT\_MAPPING & \texttt{mittelstark} & Mäßig \\ \hline
EXSUDAT\_GT\_MAPPING & \texttt{mäfig} & Mäßig \\ \hline
EXSUDAT\_GT\_MAPPING & \texttt{mäßig} & Mäßig \\ \hline
EXSUDAT\_GT\_MAPPING & \texttt{mäßig bis stark} & Mäßig, Stark \\ \hline
EXSUDAT\_GT\_MAPPING & \texttt{mäßig bis stark, sicher ist eine Geruchsbildung süßlich aromatisch teilweise beschrieben als traubenartig} & Mäßig, Stark \\ \hline
EXSUDAT\_GT\_MAPPING & \texttt{nicht beurteilbar} & Enthaltung / keine Angabe \\ \hline
EXSUDAT\_GT\_MAPPING & \texttt{nicht beurteilbar, vermutlich stark} & Stark \\ \hline
EXSUDAT\_GT\_MAPPING & \texttt{nicht genau definierbar, aus der Erfahrung heraus mäßige Exsudation} & Mäßig \\ \hline
EXSUDAT\_GT\_MAPPING & \texttt{nicht klar, eher trübe Wunde könnte Geruch abgeben} & Enthaltung / keine Angabe \\ \hline
EXSUDAT\_GT\_MAPPING & \texttt{nicht zu beschreiben} & Enthaltung / keine Angabe \\ \hline
EXSUDAT\_GT\_MAPPING & \texttt{offenen stelle mäßig, geschlossene stelle kein bis wenig exsudat} & Keine, Leicht \\ \hline
EXSUDAT\_GT\_MAPPING & \texttt{scheint mäßig zu sein} & Mäßig \\ \hline
EXSUDAT\_GT\_MAPPING & \texttt{schwach} & Leicht \\ \hline
EXSUDAT\_GT\_MAPPING & \texttt{schwach bin mäßig} & Leicht, Mäßig \\ \hline
EXSUDAT\_GT\_MAPPING & \texttt{schwach bis mäßig} & Leicht, Mäßig \\ \hline
EXSUDAT\_GT\_MAPPING & \texttt{schwach bis nicht vorhanden} & Keine, Leicht \\ \hline
EXSUDAT\_GT\_MAPPING & \texttt{sehr gering} & Leicht \\ \hline
EXSUDAT\_GT\_MAPPING & \texttt{stark} & Stark \\ \hline
EXSUDAT\_GT\_MAPPING & \texttt{vermutlich gering} & Leicht \\ \hline
EXSUDAT\_GT\_MAPPING & \texttt{vermutlich hoch} & Stark \\ \hline
EXSUDAT\_GT\_MAPPING & \texttt{vermutlich hohe Exsudation} & Stark \\ \hline
EXSUDAT\_GT\_MAPPING & \texttt{vermutlich kaum} & Keine \\ \hline
EXSUDAT\_GT\_MAPPING & \texttt{vermutlich mittel} & Mäßig \\ \hline
EXSUDAT\_GT\_MAPPING & \texttt{vermutlich mittelmässig} & Mäßig \\ \hline
EXSUDAT\_GT\_MAPPING & \texttt{vermutlich mittelmäßig} & Mäßig \\ \hline
EXSUDAT\_GT\_MAPPING & \texttt{vermutlich mässig} & Mäßig \\ \hline
EXSUDAT\_GT\_MAPPING & \texttt{vermutlich mäßig} & Mäßig \\ \hline
EXSUDAT\_GT\_MAPPING & \texttt{vermutlich starke Exsudation} & Stark \\ \hline
EXSUDAT\_GT\_MAPPING & \texttt{vermutlich vorhanden, in welcher Menge ist nicht zu beurteilen} & Enthaltung / keine Angabe \\ \hline
EXSUDAT\_GT\_MAPPING & \texttt{vermutlich wenig} & Leicht \\ \hline
EXSUDAT\_GT\_MAPPING & \texttt{wahrscheinlich mittelmäßig vorhanden} & Mäßig \\ \hline
EXSUDAT\_GT\_MAPPING & \texttt{warsch hoch bis sehr hoch,  klare exsudation} & Stark \\ \hline
HAUTSCHUTZ\_GT\_MAPPING & \texttt{Hautschutzfilm / Barrierespray} & Hautschutzfilm / Barrierespray \\ \hline
HAUTSCHUTZ\_GT\_MAPPING & \texttt{Zinksalbe / Zinkpaste} & Zinksalbe / Zinkpaste \\ \hline
HAUTSCHUTZ\_GT\_MAPPING & \texttt{Wundrandschutzpaste} & Wundrandschutzpaste \\ \hline
HAUTSCHUTZ\_GT\_MAPPING & \texttt{Nicht erforderlich} & Kein Hautschutz erforderlich \\ \hline
KOMPRESSION\_GT\_MAPPING & \texttt{Mehrkomponentensysteme (2-/4-Lagen)} & Mehrkomponentensysteme (2-/4-Lagen) \\ \hline
KOMPRESSION\_GT\_MAPPING & \texttt{Kurzzugbinden} & Kurzzugbinden \\ \hline
KOMPRESSION\_GT\_MAPPING & \texttt{Adaptive Kompressionsbandagen (Wrap)} & Adaptive Kompressionsbandagen (Wrap) \\ \hline
KOMPRESSION\_GT\_MAPPING & \texttt{Medizinische Kompressionsstrümpfe (MKS)} & Medizinische Kompressionsstrümpfe (MKS) \\ \hline
KOMPRESSION\_GT\_MAPPING & \texttt{medizinische Kompressionssysteme} & Medizinische Kompressionsstrümpfe (MKS) \\ \hline
LOKALISATION\_GT\_MAPPING & \texttt{...stark exsudierendes, nekrotisch belegtes Ulcus mit grünlicher Verfärbung typisch für Pseudomonas aeruginosa.} & Enthaltung / keine Angabe \\ \hline
LOKALISATION\_GT\_MAPPING & \texttt{?} & Enthaltung / keine Angabe \\ \hline
LOKALISATION\_GT\_MAPPING & \texttt{???} & Enthaltung / keine Angabe \\ \hline
LOKALISATION\_GT\_MAPPING & \texttt{Abdomen} & Abdomen \\ \hline
LOKALISATION\_GT\_MAPPING & \texttt{Abdomen, links} & Abdomen \\ \hline
LOKALISATION\_GT\_MAPPING & \texttt{Arm} & Arm / Hand \\ \hline
LOKALISATION\_GT\_MAPPING & \texttt{Arm oder Bein??} & Enthaltung / keine Angabe \\ \hline
LOKALISATION\_GT\_MAPPING & \texttt{Aussenseitig linke Fußsohle} & Fuß \\ \hline
LOKALISATION\_GT\_MAPPING & \texttt{Außenknöchel links} & Fuß \\ \hline
LOKALISATION\_GT\_MAPPING & \texttt{Bein} & Bein \\ \hline
LOKALISATION\_GT\_MAPPING & \texttt{Bein nicht genau definierbar} & Bein \\ \hline
LOKALISATION\_GT\_MAPPING & \texttt{Bein wo genau fraglich} & Bein \\ \hline
LOKALISATION\_GT\_MAPPING & \texttt{Bein, Unterschenkel Richtung Fuß} & Bein \\ \hline
LOKALISATION\_GT\_MAPPING & \texttt{Bein, evtl. links, Unterschenkel} & Bein \\ \hline
LOKALISATION\_GT\_MAPPING & \texttt{Bein, evtl. rechts, Unterschenkel} & Bein \\ \hline
LOKALISATION\_GT\_MAPPING & \texttt{Bein, genaue Definition nicht möglich} & Bein \\ \hline
LOKALISATION\_GT\_MAPPING & \texttt{Bein, genauere Lokalisation nicht genau definierbar} & Bein \\ \hline
LOKALISATION\_GT\_MAPPING & \texttt{Bein, wo genau nicht definierbar} & Bein \\ \hline
LOKALISATION\_GT\_MAPPING & \texttt{Das Bild zeigt ein großflächiges venöses Ulcus cruris bei ausgeprägter chronisch venöser Stauung und Adipositas-assoziierter Belastung. Möglich, Oberschenkel links, Unterschenkel} & Bein \\ \hline
LOKALISATION\_GT\_MAPPING & \texttt{Ferse} & Fuß \\ \hline
LOKALISATION\_GT\_MAPPING & \texttt{Ferse, Achillessehne} & Fuß \\ \hline
LOKALISATION\_GT\_MAPPING & \texttt{Ferse, linker Fuß} & Fuß \\ \hline
LOKALISATION\_GT\_MAPPING & \texttt{Ferse, rechter Fuß} & Fuß \\ \hline
LOKALISATION\_GT\_MAPPING & \texttt{Ferse, scheint links zu sein:
gelblich-fibrinösem Belag
mäßiger Exsudation
gerötetem Wundrand} & Fuß \\ \hline
LOKALISATION\_GT\_MAPPING & \texttt{Fuß lateral links} & Fuß \\ \hline
LOKALISATION\_GT\_MAPPING & \texttt{Fuß links plantar} & Fuß \\ \hline
LOKALISATION\_GT\_MAPPING & \texttt{Fuß links plantar  lateral} & Fuß \\ \hline
LOKALISATION\_GT\_MAPPING & \texttt{Fuß rechts} & Fuß \\ \hline
LOKALISATION\_GT\_MAPPING & \texttt{Fuß, Ferse links} & Fuß \\ \hline
LOKALISATION\_GT\_MAPPING & \texttt{Fußrücken} & Fuß \\ \hline
LOKALISATION\_GT\_MAPPING & \texttt{Fußrücken linker Fuß} & Fuß \\ \hline
LOKALISATION\_GT\_MAPPING & \texttt{Fußrücken rechter Fuß} & Fuß \\ \hline
LOKALISATION\_GT\_MAPPING & \texttt{Fußsohle Ferse} & Fuß \\ \hline
LOKALISATION\_GT\_MAPPING & \texttt{Gesäß-/Sakralregion} & Gesäß / Sakral \\ \hline
LOKALISATION\_GT\_MAPPING & \texttt{Gesäß-/perianale Region eines Säuglings} & Gesäß / Sakral \\ \hline
LOKALISATION\_GT\_MAPPING & \texttt{Knöchelinnenseite} & Fuß \\ \hline
LOKALISATION\_GT\_MAPPING & \texttt{Lokalisation: Fußrücken, lateraler Vorfuß sowie Zehenbereiche} & Fuß \\ \hline
LOKALISATION\_GT\_MAPPING & \texttt{Lokalisation: Unterschenkelbereich} & Bein \\ \hline
LOKALISATION\_GT\_MAPPING & \texttt{Lokalisation: Unterschenkelregion} & Bein \\ \hline
LOKALISATION\_GT\_MAPPING & \texttt{Lokalisation: lateraler Mittelfuß/Knöchelbereich} & Fuß \\ \hline
LOKALISATION\_GT\_MAPPING & \texttt{Lokalisation: lateraler Vor-/Mittelfuß} & Fuß \\ \hline
LOKALISATION\_GT\_MAPPING & \texttt{Lokalisation: sakrogluteale/perianale Region beidseits} & Gesäß / Sakral \\ \hline
LOKALISATION\_GT\_MAPPING & \texttt{Lokalisation: vermutlich Sakral-/Gesäßregion} & Gesäß / Sakral \\ \hline
LOKALISATION\_GT\_MAPPING & \texttt{Oberarb Muskulus bizeps} & Arm / Hand \\ \hline
LOKALISATION\_GT\_MAPPING & \texttt{Os Sacrum} & Gesäß / Sakral \\ \hline
LOKALISATION\_GT\_MAPPING & \texttt{Rima Ani} & Gesäß / Sakral \\ \hline
LOKALISATION\_GT\_MAPPING & \texttt{Sakralbereich} & Gesäß / Sakral \\ \hline
LOKALISATION\_GT\_MAPPING & \texttt{Sakralbereich linke hälfte} & Gesäß / Sakral \\ \hline
LOKALISATION\_GT\_MAPPING & \texttt{Sakraler Bereich} & Gesäß / Sakral \\ \hline
LOKALISATION\_GT\_MAPPING & \texttt{Steiß} & Gesäß / Sakral \\ \hline
LOKALISATION\_GT\_MAPPING & \texttt{Unterarm} & Arm / Hand \\ \hline
LOKALISATION\_GT\_MAPPING & \texttt{Unterschenkel} & Bein \\ \hline
LOKALISATION\_GT\_MAPPING & \texttt{Unterschenkel  medial} & Bein \\ \hline
LOKALISATION\_GT\_MAPPING & \texttt{Unterschenkel warsch. gamaschenartig / bereits zirkulierend} & Bein \\ \hline
LOKALISATION\_GT\_MAPPING & \texttt{Unterschenkel weit umfassend} & Bein \\ \hline
LOKALISATION\_GT\_MAPPING & \texttt{Unterschenkel, Außenseite auf Knie Höhe} & Bein \\ \hline
LOKALISATION\_GT\_MAPPING & \texttt{Unterschenkel, rechts, Innenseite} & Bein \\ \hline
LOKALISATION\_GT\_MAPPING & \texttt{Unterschenkel/Knöchelregion, zirkumferenziell ausgedehnt} & Bein \\ \hline
LOKALISATION\_GT\_MAPPING & \texttt{Unterschenkel?} & Bein \\ \hline
LOKALISATION\_GT\_MAPPING & \texttt{chronisches Ulkus im Bereich der Achillessehne
dorsaler Unterschenkel / Achillessehnenregion} & Bein \\ \hline
LOKALISATION\_GT\_MAPPING & \texttt{ferse, fußsohle} & Fuß \\ \hline
LOKALISATION\_GT\_MAPPING & \texttt{fibrinösem gelblich-weißem Belag,
mäßiger Exsudation,
gerötetem Wundrand,
--> vermutlich entzündlicher bzw. kritisch kolonisierter Situation,} & Enthaltung / keine Angabe \\ \hline
LOKALISATION\_GT\_MAPPING & \texttt{größeres, flächiges Ulcus am Unterschenkel (Wade) bei ausgeprägtem Ödem und Adipositas-assoziierter Hautveränderung.} & Bein \\ \hline
LOKALISATION\_GT\_MAPPING & \texttt{keine Angabe möglich} & Enthaltung / keine Angabe \\ \hline
LOKALISATION\_GT\_MAPPING & \texttt{könnte Innenseite linker Fuß sein} & Fuß \\ \hline
LOKALISATION\_GT\_MAPPING & \texttt{könnte Unterschenkel oder Handgelenk sein} & Enthaltung / keine Angabe \\ \hline
LOKALISATION\_GT\_MAPPING & \texttt{könnte am linken Fuß sein... Knöchelbereich} & Fuß \\ \hline
LOKALISATION\_GT\_MAPPING & \texttt{könnte der Unterschenkel sein, rechts} & Bein \\ \hline
LOKALISATION\_GT\_MAPPING & \texttt{könnte oberhalb des os sacrum liegen} & Gesäß / Sakral \\ \hline
LOKALISATION\_GT\_MAPPING & \texttt{könnte sich im Bereich des Steißes befinden} & Gesäß / Sakral \\ \hline
LOKALISATION\_GT\_MAPPING & \texttt{lateral linker Fuß} & Fuß \\ \hline
LOKALISATION\_GT\_MAPPING & \texttt{li lateraler Malleolus} & Fuß \\ \hline
LOKALISATION\_GT\_MAPPING & \texttt{li. Fußrücken} & Fuß \\ \hline
LOKALISATION\_GT\_MAPPING & \texttt{linke Ferse} & Fuß \\ \hline
LOKALISATION\_GT\_MAPPING & \texttt{linker Fuss} & Fuß \\ \hline
LOKALISATION\_GT\_MAPPING & \texttt{linker Fussrücken} & Fuß \\ \hline
LOKALISATION\_GT\_MAPPING & \texttt{linker Fuß} & Fuß \\ \hline
LOKALISATION\_GT\_MAPPING & \texttt{linker Fuß
Lokalisation: dorsaler und lateraler Vorfuß mit Beteiligung der Zehenbasis} & Fuß \\ \hline
LOKALISATION\_GT\_MAPPING & \texttt{linker Fuß - dorsaler Vorfuß/Fußrücken} & Fuß \\ \hline
LOKALISATION\_GT\_MAPPING & \texttt{linker Fuß außen} & Fuß \\ \hline
LOKALISATION\_GT\_MAPPING & \texttt{linker Fuß lateratl} & Fuß \\ \hline
LOKALISATION\_GT\_MAPPING & \texttt{linker Fuß medial, malleolär} & Fuß \\ \hline
LOKALISATION\_GT\_MAPPING & \texttt{linker Fuß medial-lateral 
kleine Nekrose an der Ferse außen plantar} & Fuß \\ \hline
LOKALISATION\_GT\_MAPPING & \texttt{linker Fuß, Außenseite} & Fuß \\ \hline
LOKALISATION\_GT\_MAPPING & \texttt{linker Fuß, Ferse} & Fuß \\ \hline
LOKALISATION\_GT\_MAPPING & \texttt{linker Fuß, Innenseite} & Fuß \\ \hline
LOKALISATION\_GT\_MAPPING & \texttt{linker Fuß, Mittelfußsohle, Ferse} & Fuß \\ \hline
LOKALISATION\_GT\_MAPPING & \texttt{linker Fuß, Spann, Fußrücken} & Fuß \\ \hline
LOKALISATION\_GT\_MAPPING & \texttt{linker Fuß, außenseite} & Fuß \\ \hline
LOKALISATION\_GT\_MAPPING & \texttt{linker Fuß, seite} & Fuß \\ \hline
LOKALISATION\_GT\_MAPPING & \texttt{linker Fußrücken} & Fuß \\ \hline
LOKALISATION\_GT\_MAPPING & \texttt{linker Oberarm, Innenseite} & Arm / Hand \\ \hline
LOKALISATION\_GT\_MAPPING & \texttt{linker Oberarm/ Innenseitig} & Arm / Hand \\ \hline
LOKALISATION\_GT\_MAPPING & \texttt{linker proximaler Unterschenkel lateral} & Bein \\ \hline
LOKALISATION\_GT\_MAPPING & \texttt{linkes Bein unterhalb Knie lateral} & Bein \\ \hline
LOKALISATION\_GT\_MAPPING & \texttt{lokalisation nicht genau zu definieren} & Enthaltung / keine Angabe \\ \hline
LOKALISATION\_GT\_MAPPING & \texttt{mediale Fußseite linker Fuß} & Fuß \\ \hline
LOKALISATION\_GT\_MAPPING & \texttt{mehrere kleine bis mittelgroße Ulzerationen am Unterschenkel (Wade, rechts?) mit entzündlich wirkender Umgebungshaut.} & Bein \\ \hline
LOKALISATION\_GT\_MAPPING & \texttt{multiple Schädigung gesamter Gesäßbereich} & Gesäß / Sakral \\ \hline
LOKALISATION\_GT\_MAPPING & \texttt{nicht beurteilbar} & Enthaltung / keine Angabe \\ \hline
LOKALISATION\_GT\_MAPPING & \texttt{nicht genau definierbar} & Enthaltung / keine Angabe \\ \hline
LOKALISATION\_GT\_MAPPING & \texttt{nicht genau definierbar könnte Sakralbereich sein} & Gesäß / Sakral \\ \hline
LOKALISATION\_GT\_MAPPING & \texttt{nicht genau definierbar sieht für mich aus wie am Vorfuß} & Fuß \\ \hline
LOKALISATION\_GT\_MAPPING & \texttt{nicht genau definierbar, Vermutung Unterschenkel} & Bein \\ \hline
LOKALISATION\_GT\_MAPPING & \texttt{nicht genau definierbar, meist am Unterschenkel} & Bein \\ \hline
LOKALISATION\_GT\_MAPPING & \texttt{oberhalb der Achillessehne} & Fuß \\ \hline
LOKALISATION\_GT\_MAPPING & \texttt{plantarseitig rechter Fuß} & Fuß \\ \hline
LOKALISATION\_GT\_MAPPING & \texttt{re Fuß} & Fuß \\ \hline
LOKALISATION\_GT\_MAPPING & \texttt{rechte Ferse} & Fuß \\ \hline
LOKALISATION\_GT\_MAPPING & \texttt{rechte Hand} & Arm / Hand \\ \hline
LOKALISATION\_GT\_MAPPING & \texttt{rechte Hand und Arm} & Arm / Hand \\ \hline
LOKALISATION\_GT\_MAPPING & \texttt{rechter Fuß} & Fuß \\ \hline
LOKALISATION\_GT\_MAPPING & \texttt{rechter Fuß - lateraler und plantarer Fußbereich/Ferse} & Fuß \\ \hline
LOKALISATION\_GT\_MAPPING & \texttt{rechter Fuß Fußrücken und lateral} & Fuß \\ \hline
LOKALISATION\_GT\_MAPPING & \texttt{rechter Fuß medial, Richtung plantar.} & Fuß \\ \hline
LOKALISATION\_GT\_MAPPING & \texttt{rechter Fuß plantar} & Fuß \\ \hline
LOKALISATION\_GT\_MAPPING & \texttt{rechter Fuß, Ferse} & Fuß \\ \hline
LOKALISATION\_GT\_MAPPING & \texttt{rechter Fuß, Ferse, Mittelfuß - Fußsohle} & Fuß \\ \hline
LOKALISATION\_GT\_MAPPING & \texttt{rechter Fuß, Fußsohle} & Fuß \\ \hline
LOKALISATION\_GT\_MAPPING & \texttt{rechter Fuß, Innenseite, Knöchel} & Fuß \\ \hline
LOKALISATION\_GT\_MAPPING & \texttt{rechter Fuß, Spann} & Fuß \\ \hline
LOKALISATION\_GT\_MAPPING & \texttt{rechter Fuß, in Höhe der großen Zehe, unterhalb der medialen Malleolen} & Fuß \\ \hline
LOKALISATION\_GT\_MAPPING & \texttt{rechter Unterarm, rechtes Handgelenk, rechte Hand} & Arm / Hand \\ \hline
LOKALISATION\_GT\_MAPPING & \texttt{rechter Unterschenkel} & Bein \\ \hline
LOKALISATION\_GT\_MAPPING & \texttt{rechtes Bein, Kniebereich, Unterschenkel} & Bein \\ \hline
LOKALISATION\_GT\_MAPPING & \texttt{rechtes Bein, höhe der Knöchel} & Bein \\ \hline
LOKALISATION\_GT\_MAPPING & \texttt{sacralbereich} & Gesäß / Sakral \\ \hline
LOKALISATION\_GT\_MAPPING & \texttt{sakrales Dekubitalulkus} & Gesäß / Sakral \\ \hline
LOKALISATION\_GT\_MAPPING & \texttt{schwer zu beurteilen, Vermutung am Trochanter major} & Bein \\ \hline
LOKALISATION\_GT\_MAPPING & \texttt{schwer zu beurteilen, keine genaue Angabe möglich,
eventuell Rücken Richtung Oberarm /} & Enthaltung / keine Angabe \\ \hline
LOKALISATION\_GT\_MAPPING & \texttt{schwierig zu beantworten... könnte die Wade sein, evtl. rechts} & Bein \\ \hline
LOKALISATION\_GT\_MAPPING & \texttt{unterschenkel} & Bein \\ \hline
LOKALISATION\_GT\_MAPPING & \texttt{vermutlich am Außenknöchel} & Fuß \\ \hline
LOKALISATION\_GT\_MAPPING & \texttt{vermutlich am Unterschenkel} & Bein \\ \hline
LOKALISATION\_KEYWORDS & \texttt{Fuß} & fuß, fus, ferse, knöchel, zehe, spann, sohle, plantar, vorfuß, mittelfuß, fußrücken, sprunggelenk, malleolar, malleol, dorsolateral, dorsalseitig, rückfuß, großzehe, calcaneus, außenknöchel \\ \hline
LOKALISATION\_KEYWORDS & \texttt{Bein} & bein, schenkel, wade, knie, achilles, knöchelregion, unterschenkel, gaiter \\ \hline
LOKALISATION\_KEYWORDS & \texttt{Arm} & arm, oberarm, unterarm \\ \hline
LOKALISATION\_KEYWORDS & \texttt{Hand} & hand, handgelenk \\ \hline
LOKALISATION\_KEYWORDS & \texttt{Abdomen} & abdomen, bauch, bauchdecke, abdominal \\ \hline
LOKALISATION\_KEYWORDS & \texttt{Gesäß/Steiß} & gesäß, gesaess, sakral, steiß, steiss, perianal, sakrogluteal, gluteal, intergluteal, leiste, leistenregion, inguinal \\ \hline
PRODUKT\_GT\_MAPPING & \texttt{Schaumstoffverbände (Foam)} & Schaumstoffverbände (Foam) \\ \hline
PRODUKT\_GT\_MAPPING & \texttt{Schaumstoffverbände} & Schaumstoffverbände (Foam) \\ \hline
PRODUKT\_GT\_MAPPING & \texttt{Hydrofaser / Hydrofiber} & Hydrofaser / Hydrofiber \\ \hline
PRODUKT\_GT\_MAPPING & \texttt{Alginate} & Alginate \\ \hline
PRODUKT\_GT\_MAPPING & \texttt{Hydrogele (Kompresse)} & Hydrogele (Kompresse) \\ \hline
PRODUKT\_GT\_MAPPING & \texttt{Hydrokolloide} & Hydrokolloide \\ \hline
PRODUKT\_GT\_MAPPING & \texttt{Superabsorber} & Superabsorber \\ \hline
PRODUKT\_GT\_MAPPING & \texttt{Wundkontaktschichten (Silikon/Paraffin)} & Wundkontaktschichten (Silikon/Paraffin) \\ \hline
PRODUKT\_GT\_MAPPING & \texttt{Lipidokolloid-Auflagen} & Wundkontaktschichten (Silikon/Paraffin) \\ \hline
PRODUKT\_GT\_MAPPING & \texttt{Hydropolymerverbände} & Schaumstoffverbände (Foam) \\ \hline
PRODUKT\_GT\_MAPPING & \texttt{Silberhydrofaser} & Hydrofaser / Hydrofiber \\ \hline
PRODUKT\_GT\_MAPPING & \texttt{Silberalginat} & Alginate \\ \hline
PRODUKT\_GT\_MAPPING & \texttt{Silber-Schaumverbände} & Schaumstoffverbände (Foam) \\ \hline
PRODUKT\_GT\_MAPPING & \texttt{antimikrobielle Alginate} & Alginate \\ \hline
PRODUKT\_GT\_MAPPING & \texttt{Hydrofaser / Hydrofiber, silberhaltige Varianten bei bakterieller Belastung} & Hydrofaser / Hydrofiber \\ \hline
PRODUKT\_GT\_MAPPING & \texttt{Hydrofaser / Hydrofiber, Wundkontaktschichten (Silikon/Paraffin), Bei bakterieller Belastung: silberhaltige Varianten kurzfristig} & Hydrofaser / Hydrofiber, Wundkontaktschichten (Silikon/Paraffin) \\ \hline
PRODUKT\_GT\_MAPPING & \texttt{Alginate, Hydrofaser / Hydrofiber, bei kritischer Kolonisation silberhaltige Varianten erwägen} & Hydrofaser / Hydrofiber, Alginate \\ \hline
PRODUKT\_GT\_MAPPING & \texttt{Hydrofaser / Hydrofiber, Schaumstoffverbände (Foam), Silber-Schaumverbände} & Hydrofaser / Hydrofiber, Schaumstoffverbände (Foam) \\ \hline
PRODUKT\_GT\_MAPPING & \texttt{Hydrofaser / Hydrofiber, Bei Infektionsverdacht silberhaltige Varianten, Schaumstoffverbände (Foam)} & Hydrofaser / Hydrofiber, Schaumstoffverbände (Foam) \\ \hline
PRODUKT\_GT\_MAPPING & \texttt{Schaumstoffverbände (Foam), silberhaltige Verbrennungsauflagen} & Schaumstoffverbände (Foam) \\ \hline
PRODUKT\_GT\_MAPPING & \texttt{antiseptische Wundauflagen --> Silber PHMB-haltige Produkte} & PHMB-haltige Wundauflagen \\ \hline
PRODUKT\_GT\_MAPPING & \texttt{Silberhydrofaser oder Silberalginat} & Hydrofaser / Hydrofiber, Alginate \\ \hline
PRODUKT\_GT\_MAPPING & \texttt{Silberalginate, Silberhydrofaser} & Hydrofaser / Hydrofiber, Alginate \\ \hline
PRODUKT\_GT\_MAPPING & \texttt{Silberhydrofaser, Silberalginat} & Hydrofaser / Hydrofiber, Alginate \\ \hline
PRODUKT\_GT\_MAPPING & \texttt{Silberhydrofaser, antimikrobielle Alginate} & Hydrofaser / Hydrofiber, Alginate \\ \hline
PRODUKT\_GT\_MAPPING & \texttt{silberhydrofaser/Silberalginat bei kritischer Kolonisation} & Hydrofaser / Hydrofiber, Alginate \\ \hline
PRODUKT\_GT\_MAPPING & \texttt{Silberalginat, antimikrobielle Alginate ggf. PHMB-haltige Wundauflagen} & Alginate, PHMB-haltige Wundauflagen \\ \hline
PRODUKT\_GT\_MAPPING & \texttt{Bei feuchten/infizierten Bereichen: Silberhydrofaser Silberalginat antimikrobielle Wundauflagen.} & Hydrofaser / Hydrofiber, Alginate \\ \hline
PRODUKT\_GT\_MAPPING & \texttt{Bei feuchten/infizierten Bereichen: Silberhydrofaser Silberalginat antimikrobielle Wundauflagen} & Hydrofaser / Hydrofiber, Alginate \\ \hline
PRODUKT\_GT\_MAPPING & \texttt{Bei kritischer Kolonisation: Silberhydrofaser kurzfristig, Hydrogele (Kompresse)} & Hydrofaser / Hydrofiber, Hydrogele (Kompresse) \\ \hline
PRODUKT\_GT\_MAPPING & \texttt{Bei kritischer Kolonisation: Silberhydrofaser kurzfristig} & Hydrofaser / Hydrofiber \\ \hline
PRODUKT\_GT\_MAPPING & \texttt{Alginate bei stärker exsudierend} & Alginate \\ \hline
PRODUKT\_GT\_MAPPING & \texttt{ggf. Hydrogel bei trockeneren Belägen} & Hydrogele (Kompresse) \\ \hline
PRODUKT\_GT\_MAPPING & \texttt{Alginate, Hydrofaser / Hydrofiber, Superabsorber, ggf. Unterdrucktherapie (NPWT/VAC) nach Nekrosensanierung} & Alginate, Hydrofaser / Hydrofiber, Superabsorber \\ \hline
PRODUKT\_GT\_MAPPING & \texttt{PHMB-haltige Auflagen} & PHMB-haltige Wundauflagen \\ \hline
PRODUKT\_GT\_MAPPING & \texttt{PHMB-haltige Wundauflagen} & PHMB-haltige Wundauflagen \\ \hline
PRODUKT\_GT\_MAPPING & \texttt{antimikrobielle Alginate ggf. PHMB-haltige Wundauflagen} & Alginate, PHMB-haltige Wundauflagen \\ \hline
PRODUKT\_GT\_MAPPING & \texttt{trockene sterile Schutzverbände, atraumatische Kontaktlagen, keine aggressive Befeuchtung der stabilen Nekrose} & trockene sterile Schutzverbände, Wundkontaktschichten (Silikon/Paraffin) \\ \hline
PRODUKT\_GT\_MAPPING & \texttt{Bei trockener Nekrose: trockener Schutzverband atraumatische Kontaktlagen.} & trockene sterile Schutzverbände, Wundkontaktschichten (Silikon/Paraffin) \\ \hline
PRODUKT\_GT\_MAPPING & \texttt{Bei trockenen Nekrosen: atraumatische trockene Schutzverbände} & trockene sterile Schutzverbände \\ \hline
PRODUKT\_GT\_MAPPING & \texttt{trockene sterile Abdeckung, atraumatische Schutzverbände, keine aggressive Befeuchtung} & trockene sterile Schutzverbände \\ \hline
PRODUKT\_GT\_MAPPING & \texttt{Bei trockenen Nekrosen: atraumatische trockene Schutzverbände, Bei feuchten/infizierten Bereichen: Silberhydrofaser Silberalginat antimikrobielle Wundauflagen} & trockene sterile Schutzverbände, Hydrofaser / Hydrofiber, Alginate \\ \hline
PRODUKT\_GT\_MAPPING & \texttt{Hydrogel bei trockenen Nekroseanteilen, Hydrofaser oder Alginat bei Exsudat} & Hydrogele (Kompresse), Hydrofaser / Hydrofiber, Alginate \\ \hline
SEKUNDAERVERBAND\_GT\_MAPPING & \texttt{Schaumstoffverbände (Foam)} & Schaumstoffverbände (Foam) \\ \hline
SEKUNDAERVERBAND\_GT\_MAPPING & \texttt{Schaumstoffverbände} & Schaumstoffverbände (Foam) \\ \hline
SEKUNDAERVERBAND\_GT\_MAPPING & \texttt{Schaumverband} & Schaumstoffverbände (Foam) \\ \hline
SEKUNDAERVERBAND\_GT\_MAPPING & \texttt{PU-Schaumverband} & Schaumstoffverbände (Foam) \\ \hline
SEKUNDAERVERBAND\_GT\_MAPPING & \texttt{dünner PU-Schaumverband} & Schaumstoffverbände (Foam) \\ \hline
SEKUNDAERVERBAND\_GT\_MAPPING & \texttt{atraumatische Schaumverbände} & Schaumstoffverbände (Foam) \\ \hline
SEKUNDAERVERBAND\_GT\_MAPPING & \texttt{saugfähiger Schaumverband} & Schaumstoffverbände (Foam) \\ \hline
SEKUNDAERVERBAND\_GT\_MAPPING & \texttt{atraumatischer Schaumverband} & Schaumstoffverbände (Foam) \\ \hline
SEKUNDAERVERBAND\_GT\_MAPPING & \texttt{absorbierender Schaumverband} & Schaumstoffverbände (Foam) \\ \hline
SEKUNDAERVERBAND\_GT\_MAPPING & \texttt{Superabsorber} & Superabsorber \\ \hline
SEKUNDAERVERBAND\_GT\_MAPPING & \texttt{Superabsorber-Verband} & Superabsorber \\ \hline
SEKUNDAERVERBAND\_GT\_MAPPING & \texttt{stark saugfähige Superabsorber} & Superabsorber \\ \hline
SEKUNDAERVERBAND\_GT\_MAPPING & \texttt{Superabsorber bei stärkerer Sekretion} & Superabsorber \\ \hline
SEKUNDAERVERBAND\_GT\_MAPPING & \texttt{Fixiervlies / -pflaster} & Fixiervlies / -pflaster \\ \hline
SEKUNDAERVERBAND\_GT\_MAPPING & \texttt{Elastische Fixierbinden} & Elastische Fixierbinden \\ \hline
SEKUNDAERVERBAND\_GT\_MAPPING & \texttt{Fixierbinden / kohäsive Binden} & Fixierbinden / kohäsive Binden \\ \hline
SEKUNDAERVERBAND\_GT\_MAPPING & \texttt{Nicht erforderlich (selbsthaftender Primärverband)} & Kein Sekundärverband erforderlich \\ \hline
SEKUNDAERVERBAND\_GT\_MAPPING & \texttt{saugfähiger Schaumverband, atraumatische Fixierung} & Schaumstoffverbände (Foam) \\ \hline
SEKUNDAERVERBAND\_GT\_MAPPING & \texttt{absorbierender Schaumverband, atraumatische Fixierung} & Schaumstoffverbände (Foam) \\ \hline
SEKUNDAERVERBAND\_GT\_MAPPING & \texttt{atraumatische Schaumverbände, ggf. Unterdrucktherapie (VAC) nach Débridement} & Schaumstoffverbände (Foam) \\ \hline
SEKUNDAERVERBAND\_GT\_MAPPING & \texttt{PU-Schaumverband, Atraumatische Fixierung} & Schaumstoffverbände (Foam) \\ \hline
SEKUNDAERVERBAND\_GT\_MAPPING & \texttt{weicher PU-Schaumverband, weicher PU-Schaumverband} & Schaumstoffverbände (Foam) \\ \hline
SEKUNDAERVERBAND\_GT\_MAPPING & \texttt{Elastische Fixierbinden, dünner PU-Schaumverband} & Schaumstoffverbände (Foam), Elastische Fixierbinden \\ \hline
SEKUNDAERVERBAND\_GT\_MAPPING & \texttt{Superabsorber, hochabsorbierender PU-Schaumverband} & Superabsorber, Schaumstoffverbände (Foam) \\ \hline
SEKUNDAERVERBAND\_GT\_MAPPING & \texttt{Superabsorber oder PU-Schaumverband} & Superabsorber, Schaumstoffverbände (Foam) \\ \hline
SEKUNDAERVERBAND\_GT\_MAPPING & \texttt{Superabsorber weicher PU-Schaumverband} & Superabsorber, Schaumstoffverbände (Foam) \\ \hline
SEKUNDAERVERBAND\_GT\_MAPPING & \texttt{Superabsorber, hochabsorbierende Schaumverbände} & Superabsorber, Schaumstoffverbände (Foam) \\ \hline
SEKUNDAERVERBAND\_GT\_MAPPING & \texttt{PU-Schaumverband oder Superabsorber} & Superabsorber, Schaumstoffverbände (Foam) \\ \hline
SEKUNDAERVERBAND\_GT\_MAPPING & \texttt{absorbierender Schaumverband, Superabsorber bei stärkerer Sekretion} & Schaumstoffverbände (Foam), Superabsorber \\ \hline
SEKUNDAERVERBAND\_GT\_MAPPING & \texttt{weicher PU-Schaumverband Superabsorber bei stärkerem Exsudat} & Schaumstoffverbände (Foam), Superabsorber \\ \hline
SEKUNDAERVERBAND\_GT\_MAPPING & \texttt{PU-Schaumverband oder Superabsorber bei stärkerer Exsudation} & Schaumstoffverbände (Foam), Superabsorber \\ \hline
SEKUNDAERVERBAND\_GT\_MAPPING & \texttt{PU-Schaumverband, bei stärkerem Exsudat: Superabsorber} & Schaumstoffverbände (Foam), Superabsorber \\ \hline
SEKUNDAERVERBAND\_GT\_MAPPING & \texttt{PU-Schaumverband oder kleiner Superabsorber bei stärkerer Sekretion} & Schaumstoffverbände (Foam), Superabsorber \\ \hline
SEKUNDAERVERBAND\_GT\_MAPPING & \texttt{PU-Schaumverband, Superabsorber bei stärkerer Sekretion} & Schaumstoffverbände (Foam), Superabsorber \\ \hline
SEKUNDAERVERBAND\_GT\_MAPPING & \texttt{kleiner PU-Schaumverband oder weicher Superabsorber bei stärkerer Exsudation} & Schaumstoffverbände (Foam), Superabsorber \\ \hline
SEKUNDAERVERBAND\_GT\_MAPPING & \texttt{Superabsorber ergänzen} & Superabsorber \\ \hline
SEKUNDAERVERBAND\_GT\_MAPPING & \texttt{Sterile Saugkompressen / Polsterung} & Sterile Saugkompressen / Polsterung \\ \hline
SEKUNDAERVERBAND\_GT\_MAPPING & \texttt{weicher PU-Schaumverband,absorbierende sterile Sekundärauflagen} & Schaumstoffverbände (Foam), Sterile Saugkompressen / Polsterung \\ \hline
SEKUNDAERVERBAND\_GT\_MAPPING & \texttt{silikonisierter Schaumverband, absorbierende sterile Auflagen} & Schaumstoffverbände (Foam), Sterile Saugkompressen / Polsterung \\ \hline
SEKUNDAERVERBAND\_GT\_MAPPING & \texttt{weicher absorbierender Schaumverband sterile Saugkompressen} & Schaumstoffverbände (Foam), Sterile Saugkompressen / Polsterung \\ \hline
SEKUNDAERVERBAND\_GT\_MAPPING & \texttt{weicher PU-Schaumverband, sterile Polsterung} & Schaumstoffverbände (Foam), Sterile Saugkompressen / Polsterung \\ \hline
SEKUNDAERVERBAND\_GT\_MAPPING & \texttt{silikonisierter Schaumverband, absorbierender PU-Schaumverband} & Schaumstoffverbände (Foam) \\ \hline
SEKUNDAERVERBAND\_GT\_MAPPING & \texttt{weicher PU-Schaumverband, druckentlastende Polsterung} & Schaumstoffverbände (Foam) \\ \hline
SEKUNDAERVERBAND\_GT\_MAPPING & \texttt{weicher PU-Schaumverband druckreduzierende Polsterung} & Schaumstoffverbände (Foam) \\ \hline
SEKUNDAERVERBAND\_GT\_MAPPING & \texttt{weicher PU-Schaumverband, druckreduzierende Polsterung, ggf. Superabsorber} & Schaumstoffverbände (Foam), Superabsorber \\ \hline
SEKUNDAERVERBAND\_GT\_MAPPING & \texttt{Sekundär: Superabsorber + Fixier-/Kompressionssystem} & Superabsorber \\ \hline
SEKUNDAERVERBAND\_GT\_MAPPING & \texttt{abhängig von der Kompressionstherapie} & abhängig von Kompressionstherapie \\ \hline
SEKUNDAERVERBAND\_GT\_MAPPING & \texttt{abhängig von Kompressionstherapie} & abhängig von Kompressionstherapie \\ \hline
SEKUNDAERVERBAND\_GT\_MAPPING & \texttt{über Kompression} & abhängig von Kompressionstherapie \\ \hline
SEKUNDAERVERBAND\_GT\_MAPPING & \texttt{Kompressionstherapie} & abhängig von Kompressionstherapie \\ \hline
SEKUNDAERVERBAND\_GT\_MAPPING & \texttt{kommt auf die Kompression an} & abhängig von Kompressionstherapie \\ \hline
SEKUNDAERVERBAND\_GT\_MAPPING & \texttt{abhängig von Druckentlastung} & abhängig von Druckentlastung \\ \hline
SEKUNDAERVERBAND\_GT\_MAPPING & \texttt{abhängig von der Druckentlastung} & abhängig von Druckentlastung \\ \hline
SPELLING\_MAPPING & \texttt{ulcus} & ulkus \\ \hline
SPELLING\_MAPPING & \texttt{decubitus} & dekubitus \\ \hline
SPUELLOESUNG\_GT\_MAPPING & \texttt{Neutrale Spüllösung (NaCl 0,9 \% / Ringer-Lösung)} & Neutrale Spüllösung (NaCl 0,9 \% / Ringer-Lösung) \\ \hline
SPUELLOESUNG\_GT\_MAPPING & \texttt{Antimikrobielle Spüllösung (PHMB / Octenidin / Hypochlorit)} & Antimikrobielle Spüllösung (PHMB / Octenidin / Hypochlorit) \\ \hline
SPUELLOESUNG\_GT\_MAPPING & \texttt{Neutrale Spüllösung (NaCl, Ringer)} & Neutrale Spüllösung (NaCl 0,9 \% / Ringer-Lösung) \\ \hline
SPUELLOESUNG\_GT\_MAPPING & \texttt{Antimikrobielle Spüllösung (PHMB, Octenisept)} & Antimikrobielle Spüllösung (PHMB / Octenidin / Hypochlorit) \\ \hline
SPUELLOESUNG\_GT\_MAPPING & \texttt{PHMB oder Octenidin zeitlich begrenzt} & Antimikrobielle Spüllösung (PHMB / Octenidin / Hypochlorit) \\ \hline
SPUELLOESUNG\_GT\_MAPPING & \texttt{Bei kritischer Kolonisation: PHMB oder Octenidin zeitlich begrenzt} & Antimikrobielle Spüllösung (PHMB / Octenidin / Hypochlorit) \\ \hline
SPUELLOESUNG\_GT\_MAPPING & \texttt{bei kritischer Kolonisation: PHMB oder Octenidin zeitlich begrenzt} & Antimikrobielle Spüllösung (PHMB / Octenidin / Hypochlorit) \\ \hline
SPUELLOESUNG\_GT\_MAPPING & \texttt{bei kritischer Kolonisation: PHMB/Octenidin} & Antimikrobielle Spüllösung (PHMB / Octenidin / Hypochlorit) \\ \hline
SPUELLOESUNG\_GT\_MAPPING & \texttt{bei kritischer Kolonisation: mit PHMB/Octenidin} & Antimikrobielle Spüllösung (PHMB / Octenidin / Hypochlorit) \\ \hline
SPUELLOESUNG\_GT\_MAPPING & \texttt{Antiseptika nur gezielt bei: Infektionszeichen, kritischer Kolonisation} & Antimikrobielle Spüllösung (PHMB / Octenidin / Hypochlorit) \\ \hline
WUNDRAND\_GT\_MAPPING & \texttt{["Epibolie (eingerollter Wundrand)","Gerötet / entzündlich","Mazeriert"]} & Epibolie (eingerollter Wundrand), Gerötet / entzündlich, Mazeriert \\ \hline
WUNDRAND\_GT\_MAPPING & \texttt{["Epibolie (eingerollter Wundrand)","Gerötet / entzündlich"]} & Epibolie (eingerollter Wundrand), Gerötet / entzündlich \\ \hline
WUNDRAND\_GT\_MAPPING & \texttt{["Epibolie (eingerollter Wundrand)","Mazeriert"]} & Epibolie (eingerollter Wundrand), Mazeriert \\ \hline
WUNDRAND\_GT\_MAPPING & \texttt{["Epibolie (eingerollter Wundrand)"]} & Epibolie (eingerollter Wundrand) \\ \hline
WUNDRAND\_GT\_MAPPING & \texttt{["Gelbliche fibrinöse Randbeläge"]} & nan \\ \hline
WUNDRAND\_GT\_MAPPING & \texttt{["Gerötet / entzündlich","Epibolie (eingerollter Wundrand)"]} & Epibolie (eingerollter Wundrand), Gerötet / entzündlich \\ \hline
WUNDRAND\_GT\_MAPPING & \texttt{["Gerötet / entzündlich","Mazeriert","Unterminiert (Wundtaschen)"]} & Gerötet / entzündlich, Mazeriert, Unterminiert (Wundtaschen) \\ \hline
WUNDRAND\_GT\_MAPPING & \texttt{["Gerötet / entzündlich","Mazeriert"]} & Gerötet / entzündlich, Mazeriert \\ \hline
WUNDRAND\_GT\_MAPPING & \texttt{["Gerötet / entzündlich","Taschenbildung möglich"]} & Gerötet / entzündlich, Unterminiert (Wundtaschen) \\ \hline
WUNDRAND\_GT\_MAPPING & \texttt{["Gerötet / entzündlich","Unterminiert (Wundtaschen)"]} & Gerötet / entzündlich, Unterminiert (Wundtaschen) \\ \hline
WUNDRAND\_GT\_MAPPING & \texttt{["Gerötet / entzündlich","glänzende, ödematöse Haut"]} & Gerötet / entzündlich \\ \hline
WUNDRAND\_GT\_MAPPING & \texttt{["Gerötet / entzündlich","livid-erythematöse Hautveränderungen"]} & Gerötet / entzündlich \\ \hline
WUNDRAND\_GT\_MAPPING & \texttt{["Gerötet / entzündlich","relativ scharf begrenzte Wundränder"]} & Gerötet / entzündlich \\ \hline
WUNDRAND\_GT\_MAPPING & \texttt{["Gerötet / entzündlich","unregelmäßige Wundränder"]} & Gerötet / entzündlich \\ \hline
WUNDRAND\_GT\_MAPPING & \texttt{["Gerötet / entzündlich","unregelmäßigen Rändern"]} & Gerötet / entzündlich \\ \hline
WUNDRAND\_GT\_MAPPING & \texttt{["Gerötet / entzündlich"]} & Gerötet / entzündlich \\ \hline
WUNDRAND\_GT\_MAPPING & \texttt{["Hyperkeratotisch","Gerötet / entzündlich"]} & Gerötet / entzündlich, Hyperkeratotisch \\ \hline
WUNDRAND\_GT\_MAPPING & \texttt{["Hyperkeratotisch"]} & Hyperkeratotisch \\ \hline
WUNDRAND\_GT\_MAPPING & \texttt{["Mazeriert","Epibolie (eingerollter Wundrand)","Gerötet / entzündlich"]} & Epibolie (eingerollter Wundrand), Gerötet / entzündlich, Mazeriert \\ \hline
WUNDRAND\_GT\_MAPPING & \texttt{["Mazeriert","Epibolie (eingerollter Wundrand)"]} & Epibolie (eingerollter Wundrand), Mazeriert \\ \hline
WUNDRAND\_GT\_MAPPING & \texttt{["Mazeriert","Gerötet / entzündlich","Epibolie (eingerollter Wundrand)"]} & Epibolie (eingerollter Wundrand), Gerötet / entzündlich, Mazeriert \\ \hline
WUNDRAND\_GT\_MAPPING & \texttt{["Mazeriert","Gerötet / entzündlich","Unterminiert (Wundtaschen)"]} & Gerötet / entzündlich, Mazeriert, Unterminiert (Wundtaschen) \\ \hline
WUNDRAND\_GT\_MAPPING & \texttt{["Mazeriert","Gerötet / entzündlich"]} & Gerötet / entzündlich, Mazeriert \\ \hline
WUNDRAND\_GT\_MAPPING & \texttt{["Mazeriert","Unterminiert (Wundtaschen)","Gerötet / entzündlich"]} & Gerötet / entzündlich, Mazeriert, Unterminiert (Wundtaschen) \\ \hline
WUNDRAND\_GT\_MAPPING & \texttt{["Mazeriert","Unterminiert (Wundtaschen)"]} & Mazeriert, Unterminiert (Wundtaschen) \\ \hline
WUNDRAND\_GT\_MAPPING & \texttt{["Mazeriert","teilweise aufgequollen"]} & Mazeriert \\ \hline
WUNDRAND\_GT\_MAPPING & \texttt{["Mazeriert"]} & Mazeriert \\ \hline
WUNDRAND\_GT\_MAPPING & \texttt{["Reizlos / unauffällig"]} & Reizlos / unauffällig \\ \hline
WUNDRAND\_GT\_MAPPING & \texttt{["Unterminiert (Wundtaschen)","Epibolie (eingerollter Wundrand)"]} & Epibolie (eingerollter Wundrand), Unterminiert (Wundtaschen) \\ \hline
WUNDRAND\_GT\_MAPPING & \texttt{["Unterminiert (Wundtaschen)","Gerötet / entzündlich"]} & Gerötet / entzündlich, Unterminiert (Wundtaschen) \\ \hline
WUNDRAND\_GT\_MAPPING & \texttt{["Unterminiert (Wundtaschen)","Mazeriert"]} & Mazeriert, Unterminiert (Wundtaschen) \\ \hline
WUNDRAND\_GT\_MAPPING & \texttt{["Unterminiert (Wundtaschen)"]} & Unterminiert (Wundtaschen) \\ \hline
WUNDRAND\_GT\_MAPPING & \texttt{["scharf begrenzt, teils livide/verfärbt, ischämietypisch"]} & Reizlos / unauffällig \\ \hline
WUNDTYP\_GT\_MAPPING & \texttt{Ausgedehnte feuchte Nekrose} & Ausgedehnte feuchte Nekrose \\ \hline
WUNDTYP\_GT\_MAPPING & \texttt{Ausgetrtenes Blut oder Lymphe aus den entsprechenden Gefäßen mit Verteilung an der betroffenen Extremität} & Ulkus (Ätiologie unspezifisch) \\ \hline
WUNDTYP\_GT\_MAPPING & \texttt{Dekubitus} & Dekubitus \\ \hline
WUNDTYP\_GT\_MAPPING & \texttt{Dekubitus Grad 2 oder 3 nach EPUAP.
Nicht genau zu definieren, ob sich Unterminierungen vorhanden sind.} & Dekubitus \\ \hline
WUNDTYP\_GT\_MAPPING & \texttt{Dekubitus an der Ferse} & Dekubitus \\ \hline
WUNDTYP\_GT\_MAPPING & \texttt{Dekubitus epuap Stadium 3 / 4} & Dekubitus \\ \hline
WUNDTYP\_GT\_MAPPING & \texttt{Dekubitus mit Nekrose} & Dekubitus \\ \hline
WUNDTYP\_GT\_MAPPING & \texttt{Dekubitus mit Taschenbildung} & Dekubitus \\ \hline
WUNDTYP\_GT\_MAPPING & \texttt{Dekubitus sacralbereich} & Dekubitus \\ \hline
WUNDTYP\_GT\_MAPPING & \texttt{Dekubitus, vermutlich am Außenknöchel} & Dekubitus \\ \hline
WUNDTYP\_GT\_MAPPING & \texttt{Dekubtis im Sakralbereich, großflächig} & Dekubitus \\ \hline
WUNDTYP\_GT\_MAPPING & \texttt{Fersendekubitus} & Dekubitus \\ \hline
WUNDTYP\_GT\_MAPPING & \texttt{Fersenulcus / eventuell Mischulcus} & Ulcus cruris mixtum \\ \hline
WUNDTYP\_GT\_MAPPING & \texttt{Fibrinbelegtes Ulcus} & Ulkus (Ätiologie unspezifisch) \\ \hline
WUNDTYP\_GT\_MAPPING & \texttt{Fibringbelegtes Ulcus} & Ulkus (Ätiologie unspezifisch) \\ \hline
WUNDTYP\_GT\_MAPPING & \texttt{Hypergranulation nach Verbrennung} & Verbrennungswunde \\ \hline
WUNDTYP\_GT\_MAPPING & \texttt{Hämangiom ist ein Spezialgebiet. In der Regel ist dieses innerhalb des ersten Lebensjahres rückläufig.
Deshalb nehme ich hier keine Beurteilung vor.} & Hämangiom ist ein Spezialgebiet. In der Regel ist dieses innerhalb des ersten Lebensjahres rückläufig.
Deshalb nehme ich hier keine Beurteilung vor. \\ \hline
WUNDTYP\_GT\_MAPPING & \texttt{Infizierte Wunden am linken Oberarm / Innenseitig} & Infizierte Wunden am linken Oberarm / Innenseitig \\ \hline
WUNDTYP\_GT\_MAPPING & \texttt{Infiziertes Ulcus mit freiliegender Sehne} & Ulkus (Ätiologie unspezifisch) \\ \hline
WUNDTYP\_GT\_MAPPING & \texttt{Oberflächliches Ulcus} & Ulkus (Ätiologie unspezifisch) \\ \hline
WUNDTYP\_GT\_MAPPING & \texttt{Oberflächliches Ulcus (diffuse Defekte)} & Ulkus (Ätiologie unspezifisch) \\ \hline
WUNDTYP\_GT\_MAPPING & \texttt{Plantares Ulcus linker Fuß, Ursachen Klärung notwendig, bei Diabetiker Prüfung ob Neuropathisch oder Angiopathische Ursache} & Plantares Ulcus linker Fuß, Ursachen Klärung notwendig, bei Diabetiker Prüfung ob Neuropathisch oder Angiopathische Ursache \\ \hline
WUNDTYP\_GT\_MAPPING & \texttt{Platzbauch nach möglichem chirurgischen Eingriff} & Postoperative Wunde / Dehiszenz \\ \hline
WUNDTYP\_GT\_MAPPING & \texttt{Postoperative Wunde abdominal} & Postoperative Wunde / Dehiszenz \\ \hline
WUNDTYP\_GT\_MAPPING & \texttt{Taschenbildendes Ulcus am linken Fußrücken} & Ulkus (Ätiologie unspezifisch) \\ \hline
WUNDTYP\_GT\_MAPPING & \texttt{Teils nekrotisches Ulcus an der linken Ferse.
plantares Ulcus belegt mit Infektionszeichen} & Ulkus (Ätiologie unspezifisch) \\ \hline
WUNDTYP\_GT\_MAPPING & \texttt{Ulcera rechter Fuß. könnte ursächlich diabetischer Fuß sein ( Neuropathie, Angiopathie) oder arterielles ulcus schäden auf Grund fehlender Durchblutung} & Ulkus (Ätiologie unspezifisch) \\ \hline
WUNDTYP\_GT\_MAPPING & \texttt{Ulcerationen bedingt durch Vaskulitis- Form der rheumatischen Erkrankung ( Autoimmunerkrankung)} & Ulkus (Ätiologie unspezifisch) \\ \hline
WUNDTYP\_GT\_MAPPING & \texttt{Ulcus} & Ulkus (Ätiologie unspezifisch) \\ \hline
WUNDTYP\_GT\_MAPPING & \texttt{Ulcus / warsch. Ulcus cruris, Ursache nicht klar, muss bestimmt werden} & Ulkus (Ätiologie unspezifisch) \\ \hline
WUNDTYP\_GT\_MAPPING & \texttt{Ulcus Dekubitus} & Dekubitus \\ \hline
WUNDTYP\_GT\_MAPPING & \texttt{Ulcus Dekubitus Bereich OS Sacrum} & Dekubitus \\ \hline
WUNDTYP\_GT\_MAPPING & \texttt{Ulcus Dekubitus Epuap Grad 3 / 4} & Dekubitus \\ \hline
WUNDTYP\_GT\_MAPPING & \texttt{Ulcus Dekubitus nach Epuap Grad 2 bis 3} & Dekubitus \\ \hline
WUNDTYP\_GT\_MAPPING & \texttt{Ulcus Ursache nicht genau definierbar} & Ulkus (Ätiologie unspezifisch) \\ \hline
WUNDTYP\_GT\_MAPPING & \texttt{Ulcus am Fussrücken} & Ulkus (Ätiologie unspezifisch) \\ \hline
WUNDTYP\_GT\_MAPPING & \texttt{Ulcus am Unterschenkel medial} & Ulkus (Ätiologie unspezifisch) \\ \hline
WUNDTYP\_GT\_MAPPING & \texttt{Ulcus an der Achillessehne} & Ulkus (Ätiologie unspezifisch) \\ \hline
WUNDTYP\_GT\_MAPPING & \texttt{Ulcus auf Grund von Mangeldurchblutung. Könnte diabetischer Fuß sein / angiopathisch bedingt.} & Ulkus (Ätiologie unspezifisch) \\ \hline
WUNDTYP\_GT\_MAPPING & \texttt{Ulcus auf Grund von Mangelernährung des Gewebes, Druckschäden in Kombination mit Gefäßschäden} & Ulkus (Ätiologie unspezifisch) \\ \hline
WUNDTYP\_GT\_MAPPING & \texttt{Ulcus cruris} & Ulkus (Ätiologie unspezifisch) \\ \hline
WUNDTYP\_GT\_MAPPING & \texttt{Ulcus cruris  unklarer Ursache} & Ulkus (Ätiologie unspezifisch) \\ \hline
WUNDTYP\_GT\_MAPPING & \texttt{Ulcus cruris /  Gravitationsulcus: Ulcus cruris venosum} & Ulcus cruris venosum \\ \hline
WUNDTYP\_GT\_MAPPING & \texttt{Ulcus cruris li lateraler Malleolus} & Ulkus (Ätiologie unspezifisch) \\ \hline
WUNDTYP\_GT\_MAPPING & \texttt{Ulcus cruris möglich arteriell oder mixtum} & Ulkus (Ätiologie unspezifisch) \\ \hline
WUNDTYP\_GT\_MAPPING & \texttt{Ulcus cruris venosum} & Ulcus cruris venosum \\ \hline
WUNDTYP\_GT\_MAPPING & \texttt{Ulcus decubitus} & Dekubitus \\ \hline
WUNDTYP\_GT\_MAPPING & \texttt{Ulcus decubitus Nach Epuap Grad 2 bis 3} & Dekubitus \\ \hline
WUNDTYP\_GT\_MAPPING & \texttt{Ulcus dekubitus} & Dekubitus \\ \hline
WUNDTYP\_GT\_MAPPING & \texttt{Ulcus dekubitus Ferse,  Stadium 3 / 4} & Dekubitus \\ \hline
WUNDTYP\_GT\_MAPPING & \texttt{Ulcus lokal abgegrenzt. belegt könnte ebenso nekrotisch sein} & Ulkus (Ätiologie unspezifisch) \\ \hline
WUNDTYP\_GT\_MAPPING & \texttt{Ulcus mit Stauung} & Ulkus (Ätiologie unspezifisch) \\ \hline
WUNDTYP\_GT\_MAPPING & \texttt{Ulcus mit scheinender Hypergranulation oberhalb der Achillessehne} & Ulkus (Ätiologie unspezifisch) \\ \hline
WUNDTYP\_GT\_MAPPING & \texttt{Ulcus nach Meshgraft (zumindest lässt die Struktur darauf schließen), eventuell ist die Struktur aber einer Unterdrucktherapie geschuldet} & Postoperative Wunde / Dehiszenz \\ \hline
WUNDTYP\_GT\_MAPPING & \texttt{Ulcus nach Stich} & Traumatische Wunde \\ \hline
WUNDTYP\_GT\_MAPPING & \texttt{Ulcus unbekannter Herkunft, könnte traumatologisch bedingte Wunde sein, durch die Wunde bei unsachgemäßer Versorgung entstehen ulcerationen ebenfalls, nicht immer Krankheitsbedingt} & Ulkus (Ätiologie unspezifisch) \\ \hline
WUNDTYP\_GT\_MAPPING & \texttt{Ulcus unklarer Genese} & Ulkus (Ätiologie unspezifisch) \\ \hline
WUNDTYP\_GT\_MAPPING & \texttt{Ulcus unklarer Genese am Bein genaue Lokalisation nicht beurteilbar} & Ulkus (Ätiologie unspezifisch) \\ \hline
WUNDTYP\_GT\_MAPPING & \texttt{Ulcus unklarer Genese, könnte Venös oder arteriell oder mixum sein.
Beachte kleine Nekrose an der Ferse plantar} & Ulkus (Ätiologie unspezifisch) \\ \hline
WUNDTYP\_GT\_MAPPING & \texttt{Ulcus unklarer Genese, könnte diabet. Ulcus sein, möglich ebenso durch traumatologische Ursachen} & Ulkus (Ätiologie unspezifisch) \\ \hline
WUNDTYP\_GT\_MAPPING & \texttt{Ulcus, Lokalisation am Bein nicht genau definierbar} & Ulkus (Ätiologie unspezifisch) \\ \hline
WUNDTYP\_GT\_MAPPING & \texttt{Ulcus, bei vermutetem Diabetes} & Ulkus (Ätiologie unspezifisch) \\ \hline
WUNDTYP\_GT\_MAPPING & \texttt{Ulcus, durch die Mangelernährung ( keine Durchblutung) im Wundgebiet  könnte es sich um einen Dekubitus oder ulcus cruris handeln.} & Ulkus (Ätiologie unspezifisch) \\ \hline
WUNDTYP\_GT\_MAPPING & \texttt{Ulcus. Ursache klären sollte es ein diabet. Fuß sein dann Prüfung ob angiopathisch oder neuropathisch,
Druckentlastung und Ursachenbeseitigung.} & Ulkus (Ätiologie unspezifisch) \\ \hline
WUNDTYP\_GT\_MAPPING & \texttt{Ulkus Dekubitus am Os sacrum  nach Epuap Grad 2 / 3} & Dekubitus \\ \hline
WUNDTYP\_GT\_MAPPING & \texttt{Ulkus an der Ferse, warsch. Dekubitus nach EPUAP  Stadium 3} & Dekubitus \\ \hline
WUNDTYP\_GT\_MAPPING & \texttt{Verbrennung 2 und 3 Grades nach 9 er Regel etwa 9 \% Körperoberfläche} & Verbrennungswunde \\ \hline
WUNDTYP\_GT\_MAPPING & \texttt{Verbrennung 2. Grades am Bein} & Verbrennungswunde \\ \hline
WUNDTYP\_GT\_MAPPING & \texttt{Verbrennung 2. und teilweise 3. Grades Oberfläche ca 9 \%} & Verbrennungswunde \\ \hline
WUNDTYP\_GT\_MAPPING & \texttt{Verbrennung Grad 1 bis 2 am rechten Fuß} & Verbrennungswunde \\ \hline
WUNDTYP\_GT\_MAPPING & \texttt{Verbrennungswunde 1. bis 2. Grades} & Verbrennungswunde \\ \hline
WUNDTYP\_GT\_MAPPING & \texttt{Verbrennungswunde am Arm Grad 2} & Verbrennungswunde \\ \hline
WUNDTYP\_GT\_MAPPING & \texttt{Wunde am li Fußrücken mit Sehnenfreilegung} & Wunde am li Fußrücken mit Sehnenfreilegung \\ \hline
WUNDTYP\_GT\_MAPPING & \texttt{Wunde aufgrund einer Infusion, welche ins umliegende Gewebe ausgetreten ist} & Wunde aufgrund einer Infusion, welche ins umliegende Gewebe ausgetreten ist \\ \hline
WUNDTYP\_GT\_MAPPING & \texttt{["Dekubitus","Diabetisches Fußulkus"]} & Ulkus (Ätiologie unspezifisch) \\ \hline
WUNDTYP\_GT\_MAPPING & \texttt{["Dekubitus","Ulcus"]} & Ulkus (Ätiologie unspezifisch) \\ \hline
WUNDTYP\_GT\_MAPPING & \texttt{["Dekubitus"]} & Dekubitus \\ \hline
WUNDTYP\_GT\_MAPPING & \texttt{["Diabetisches Fußulkus","hochgradig destruierende Fußwunde"]} & Diabetisches Fußsyndrom (DFS) \\ \hline
WUNDTYP\_GT\_MAPPING & \texttt{["Diabetisches Fußulkus","ischämisches Ulkus"]} & Diabetisches Fußsyndrom (DFS) \\ \hline
WUNDTYP\_GT\_MAPPING & \texttt{["Diabetisches Fußulkus","mit deutlichen Zeichen einer lokalen Infektion und chronisch entzündlicher Hautveränderungen"]} & Diabetisches Fußsyndrom (DFS) \\ \hline
WUNDTYP\_GT\_MAPPING & \texttt{["Diabetisches Fußulkus"]} & Diabetisches Fußsyndrom (DFS) \\ \hline
WUNDTYP\_GT\_MAPPING & \texttt{["Extravasationsverletzung"]} & Extravasationsverletzung \\ \hline
WUNDTYP\_GT\_MAPPING & \texttt{["Gravitationsulzera/Druckulzera"]} & Ulkus (Ätiologie unspezifisch) \\ \hline
WUNDTYP\_GT\_MAPPING & \texttt{["Postoperative Wunde"]} & Postoperative Wunde / Dehiszenz \\ \hline
WUNDTYP\_GT\_MAPPING & \texttt{["Traumatische Wunde","durch Insektenstich"]} & Traumatische Wunde \\ \hline
WUNDTYP\_GT\_MAPPING & \texttt{["Ulcus Fußrücken mit deutlicher Gewebedestruktion und freiliegender Sehnenstruktur"]} & Ulkus (Ätiologie unspezifisch) \\ \hline
WUNDTYP\_GT\_MAPPING & \texttt{["Ulcus cruris venosum","Adipositas-assoziiert"]} & Ulcus cruris venosum \\ \hline
WUNDTYP\_GT\_MAPPING & \texttt{["Ulcus cruris venosum","venös-lymphatischem Ulcus"]} & Ulcus cruris venosum \\ \hline
WUNDTYP\_GT\_MAPPING & \texttt{["Ulcus cruris venosum"]} & Ulcus cruris venosum \\ \hline
WUNDTYP\_GT\_MAPPING & \texttt{["Ulcus"]} & Ulkus (Ätiologie unspezifisch) \\ \hline
WUNDTYP\_GT\_MAPPING & \texttt{["Ulkus cruris"]} & Ulkus (Ätiologie unspezifisch) \\ \hline
WUNDTYP\_GT\_MAPPING & \texttt{["Ulkus"]} & Ulkus (Ätiologie unspezifisch) \\ \hline
WUNDTYP\_GT\_MAPPING & \texttt{["Verbrennungswunde"]} & Verbrennungswunde \\ \hline
WUNDTYP\_GT\_MAPPING & \texttt{["ausgedehnte, tiefe, nekrotische Ulzeration"]} & Ulkus (Ätiologie unspezifisch) \\ \hline
WUNDTYP\_GT\_MAPPING & \texttt{["ausgedehntes, tiefes chronisches Ulcus"]} & Ulkus (Ätiologie unspezifisch) \\ \hline
WUNDTYP\_GT\_MAPPING & \texttt{["ausgeprägtes infiziertes Ulcus cruris","Ulcus cruris venosum"]} & Ulkus (Ätiologie unspezifisch) \\ \hline
WUNDTYP\_GT\_MAPPING & \texttt{["chronische Fußwunde / Ulkus"]} & Ulkus (Ätiologie unspezifisch) \\ \hline
WUNDTYP\_GT\_MAPPING & \texttt{["chronisches Ulkus cruris"]} & Ulkus (Ätiologie unspezifisch) \\ \hline
WUNDTYP\_GT\_MAPPING & \texttt{["chronisches Ulkus"]} & Ulkus (Ätiologie unspezifisch) \\ \hline
WUNDTYP\_GT\_MAPPING & \texttt{["großflächiges, flaches chronisches Ulkus"]} & Ulkus (Ätiologie unspezifisch) \\ \hline
WUNDTYP\_GT\_MAPPING & \texttt{["kleines, oberflächliches Ulcus"]} & Ulkus (Ätiologie unspezifisch) \\ \hline
WUNDTYP\_GT\_MAPPING & \texttt{["multiple infizierte chronische Fußulzera"]} & Ulkus (Ätiologie unspezifisch) \\ \hline
WUNDTYP\_GT\_MAPPING & \texttt{["multiple kleine Ulzera"]} & Ulkus (Ätiologie unspezifisch) \\ \hline
WUNDTYP\_GT\_MAPPING & \texttt{["nekrotische Fersenwunde","Dekubitus"]} & Ulkus (Ätiologie unspezifisch) \\ \hline
WUNDTYP\_GT\_MAPPING & \texttt{["nekrotische ischämische Fußwunde / Gangrän"]} & Ulcus cruris arteriosum / Ischämisches Ulkus \\ \hline
WUNDTYP\_GT\_MAPPING & \texttt{["neuroischämisches bzw. arterielles Fußulkus"]} & Ulcus cruris arteriosum / Ischämisches Ulkus \\ \hline
WUNDTYP\_GT\_MAPPING & \texttt{["rundlich-ovales Ulcus"]} & Ulkus (Ätiologie unspezifisch) \\ \hline
WUNDTYP\_GT\_MAPPING & \texttt{["tiefer reichende Ulzerationen am Unterschenkel mit chronisch entzündlicher Umgebung"]} & Ulkus (Ätiologie unspezifisch) \\ \hline
WUNDTYP\_GT\_MAPPING & \texttt{["ulzeriertes infantiles Hämangiom mit zentraler Nekrose"]} & ulzeriertes infantiles Hämangiom mit zentraler Nekrose \\ \hline
WUNDTYP\_GT\_MAPPING & \texttt{["vaskulitischen Ulzera","möglicher kutaner Vaskulitis"]} & Ulkus (Ätiologie unspezifisch) \\ \hline
WUNDTYP\_GT\_MAPPING & \texttt{arterielles Ulcus} & Ulcus cruris arteriosum / Ischämisches Ulkus \\ \hline
WUNDTYP\_GT\_MAPPING & \texttt{belegter Ulcus} & Ulkus (Ätiologie unspezifisch) \\ \hline
WUNDTYP\_GT\_MAPPING & \texttt{diffuse fibrinbelegte und teils infizierte Ulzera am rechten Fuß} & Ulkus (Ätiologie unspezifisch) \\ \hline
WUNDTYP\_GT\_MAPPING & \texttt{fibrinbelegtes Ulcus} & Ulkus (Ätiologie unspezifisch) \\ \hline
WUNDTYP\_GT\_MAPPING & \texttt{gleiche Beurteilung wie Wunde 53} & Dekubitus \\ \hline
WUNDTYP\_GT\_MAPPING & \texttt{infantiles Hämangiom,} & infantiles Hämangiom, \\ \hline
WUNDTYP\_GT\_MAPPING & \texttt{könnte arterieller Ulcus oder diabetischer Ulcus ( Angiopathie) auf Grund der Mangeldurchblutung  sein.} & Ulkus (Ätiologie unspezifisch) \\ \hline
WUNDTYP\_GT\_MAPPING & \texttt{könnte diabetischer ulcus sein auf grund einer neuropathie} & Ulkus (Ätiologie unspezifisch) \\ \hline
WUNDTYP\_GT\_MAPPING & \texttt{könnte ein diabetischer Fuß sein, Ulcus  nicht möglich,
traumatologische Indikation ebenso nicht ausgeschlossen} & Ulkus (Ätiologie unspezifisch) \\ \hline
WUNDTYP\_GT\_MAPPING & \texttt{laterales Ulcus linker Fuß} & Ulkus (Ätiologie unspezifisch) \\ \hline
WUNDTYP\_GT\_MAPPING & \texttt{massives Ulcus dekubitus} & Dekubitus \\ \hline
WUNDTYP\_GT\_MAPPING & \texttt{mumifizierte Zehen (D4 und D5), gleichzeitig Nekrose an der rechten Ferse} & Ulcus cruris arteriosum / Ischämisches Ulkus \\ \hline
WUNDTYP\_GT\_MAPPING & \texttt{möglicherweise ein feuchtes Gangrän} & Ulcus cruris arteriosum / Ischämisches Ulkus \\ \hline
WUNDTYP\_GT\_MAPPING & \texttt{nekrotischen Wunden am Fuß} & nekrotischen Wunden am Fuß \\ \hline
WUNDTYP\_GT\_MAPPING & \texttt{nekrotischer Defekt / Dekubitus} & Dekubitus \\ \hline
WUNDTYP\_GT\_MAPPING & \texttt{nekrotischer Dekubitus li Ferse} & Dekubitus \\ \hline
WUNDTYP\_GT\_MAPPING & \texttt{oberflächlicher Dekubitus ( 2 Stück)} & Dekubitus \\ \hline
WUNDTYP\_GT\_MAPPING & \texttt{oberflächliches Ulcus} & Ulkus (Ätiologie unspezifisch) \\ \hline
WUNDTYP\_GT\_MAPPING & \texttt{offene ulcera} & Ulkus (Ätiologie unspezifisch) \\ \hline
WUNDTYP\_GT\_MAPPING & \texttt{plantares ulcus} & Ulkus (Ätiologie unspezifisch) \\ \hline
WUNDTYP\_GT\_MAPPING & \texttt{pseudomas besiedeltes teils nekrotisches Ulcus.
Semizirkulär} & Ulkus (Ätiologie unspezifisch) \\ \hline
WUNDTYP\_GT\_MAPPING & \texttt{rechter Fuß lateral/ malleolär} & rechter Fuß lateral/ malleolär \\ \hline
WUNDTYP\_GT\_MAPPING & \texttt{schwere Nekrose am Fuß rechts lateral und plantar, offene Geschwüre über den Malleolen lateral
deutet auf arterielles Ulcus hin} & Ulcus cruris arteriosum / Ischämisches Ulkus \\ \hline
WUNDTYP\_GT\_MAPPING & \texttt{schwere Ulzerationen am Fuß- mehrere offene Stellen, freiliegende Sehen, teilweise nekrotisch, belegt} & Ulkus (Ätiologie unspezifisch) \\ \hline
WUNDTYP\_GT\_MAPPING & \texttt{semizirkuläres Ulcus} & Ulkus (Ätiologie unspezifisch) \\ \hline
WUNDTYP\_GT\_MAPPING & \texttt{semizirkuläres Ulcus re Fuß} & Ulkus (Ätiologie unspezifisch) \\ \hline
WUNDTYP\_GT\_MAPPING & \texttt{superinfizierte Ulcration am rechen Fuß offene Ulcera} & Ulkus (Ätiologie unspezifisch) \\ \hline
WUNDTYP\_GT\_MAPPING & \texttt{tiefer Dekubitus} & Dekubitus \\ \hline
WUNDTYP\_GT\_MAPPING & \texttt{tiefes Ulcus im Fersenbereich} & Ulkus (Ätiologie unspezifisch) \\ \hline
WUNDTYP\_GT\_MAPPING & \texttt{traumatologische Wunde ( thermische Schädigung / Verbrennung 3. und 4. Grades} & Verbrennungswunde \\ \hline
WUNDTYP\_GT\_MAPPING & \texttt{ulcus cruris} & Ulkus (Ätiologie unspezifisch) \\ \hline
WUNDTYP\_GT\_MAPPING & \texttt{vermutlich Mischulcus am rechten Bein, höhe der Knöchel.
Semizirkulär} & Ulcus cruris mixtum \\ \hline
WUNDTYP\_GT\_MAPPING & \texttt{zirkuläres Ulcus Unterschenkel (vermutlich Mischulcus)} & Ulcus cruris mixtum \\ \hline
WUNDUMGEBUNG\_GT\_MAPPING & \texttt{["CVI-typische Hautveränderungen (Hyperpigmentierung, Atrophie blanche, Lipodermatosklerose)","Erythem / Rötung"]} & CVI-typische Hautveränderungen (Hyperpigmentierung, Atrophie blanche, Lipodermatosklerose), Erythem / Rötung \\ \hline
WUNDUMGEBUNG\_GT\_MAPPING & \texttt{["CVI-typische Hautveränderungen (Hyperpigmentierung, Atrophie blanche, Lipodermatosklerose)"]} & CVI-typische Hautveränderungen (Hyperpigmentierung, Atrophie blanche, Lipodermatosklerose) \\ \hline
WUNDUMGEBUNG\_GT\_MAPPING & \texttt{["Ekzem / Dermatitis","Erythem / Rötung"]} & Ekzem / Dermatitis, Erythem / Rötung \\ \hline
WUNDUMGEBUNG\_GT\_MAPPING & \texttt{["Ekzem / Dermatitis"]} & Ekzem / Dermatitis \\ \hline
WUNDUMGEBUNG\_GT\_MAPPING & \texttt{["Erythem / Rötung","CVI-typische Hautveränderungen (Hyperpigmentierung, Atrophie blanche, Lipodermatosklerose)"]} & CVI-typische Hautveränderungen (Hyperpigmentierung, Atrophie blanche, Lipodermatosklerose), Erythem / Rötung \\ \hline
WUNDUMGEBUNG\_GT\_MAPPING & \texttt{["Erythem / Rötung","Ekzem / Dermatitis","CVI-typische Hautveränderungen (Hyperpigmentierung, Atrophie blanche, Lipodermatosklerose)"]} & CVI-typische Hautveränderungen (Hyperpigmentierung, Atrophie blanche, Lipodermatosklerose), Ekzem / Dermatitis, Erythem / Rötung \\ \hline
WUNDUMGEBUNG\_GT\_MAPPING & \texttt{["Erythem / Rötung","Ekzem / Dermatitis"]} & Ekzem / Dermatitis, Erythem / Rötung \\ \hline
WUNDUMGEBUNG\_GT\_MAPPING & \texttt{["Erythem / Rötung","Mazeration","CVI-typische Hautveränderungen (Hyperpigmentierung, Atrophie blanche, Lipodermatosklerose)"]} & CVI-typische Hautveränderungen (Hyperpigmentierung, Atrophie blanche, Lipodermatosklerose), Erythem / Rötung, Mazeration \\ \hline
WUNDUMGEBUNG\_GT\_MAPPING & \texttt{["Erythem / Rötung","Mazeration"]} & Erythem / Rötung, Mazeration \\ \hline
WUNDUMGEBUNG\_GT\_MAPPING & \texttt{["Erythem / Rötung","Weichteilreizung"]} & Erythem / Rötung \\ \hline
WUNDUMGEBUNG\_GT\_MAPPING & \texttt{["Erythem / Rötung","deutliche entzündliche Rötung der Umgebung"]} & Erythem / Rötung \\ \hline
WUNDUMGEBUNG\_GT\_MAPPING & \texttt{["Erythem / Rötung","eher trocken wirkende Umgebung"]} & Erythem / Rötung \\ \hline
WUNDUMGEBUNG\_GT\_MAPPING & \texttt{["Erythem / Rötung","entzündliche Umgebungshaut"]} & Erythem / Rötung \\ \hline
WUNDUMGEBUNG\_GT\_MAPPING & \texttt{["Erythem / Rötung","leicht gerötete, glänzende Umgebungshaut"]} & Erythem / Rötung \\ \hline
WUNDUMGEBUNG\_GT\_MAPPING & \texttt{["Erythem / Rötung","Ödem","CVI-typische Hautveränderungen (Hyperpigmentierung, Atrophie blanche, Lipodermatosklerose)"]} & CVI-typische Hautveränderungen (Hyperpigmentierung, Atrophie blanche, Lipodermatosklerose), Erythem / Rötung, Ödem \\ \hline
WUNDUMGEBUNG\_GT\_MAPPING & \texttt{["Erythem / Rötung","Ödem"]} & Erythem / Rötung, Ödem \\ \hline
WUNDUMGEBUNG\_GT\_MAPPING & \texttt{["Erythem / Rötung"]} & Erythem / Rötung \\ \hline
WUNDUMGEBUNG\_GT\_MAPPING & \texttt{["Keratosen"]} & Sonstiges \\ \hline
WUNDUMGEBUNG\_GT\_MAPPING & \texttt{["Mazeration","Erythem / Rötung","trockener wirkende Areale mit verminderter Durchblutung"]} & Erythem / Rötung, Mazeration \\ \hline
WUNDUMGEBUNG\_GT\_MAPPING & \texttt{["Mazeration","Erythem / Rötung"]} & Erythem / Rötung, Mazeration \\ \hline
WUNDUMGEBUNG\_GT\_MAPPING & \texttt{["Mazeration","glänzend"]} & Mazeration \\ \hline
WUNDUMGEBUNG\_GT\_MAPPING & \texttt{["Mazeration","glänzende, gespannte Umgebungshaut"]} & Mazeration \\ \hline
WUNDUMGEBUNG\_GT\_MAPPING & \texttt{["Mazeration"]} & Mazeration \\ \hline
WUNDUMGEBUNG\_GT\_MAPPING & \texttt{["Reizlos / intakt","Erythem / Rötung"]} & Erythem / Rötung, Reizlos / intakt \\ \hline
WUNDUMGEBUNG\_GT\_MAPPING & \texttt{["Reizlos / intakt"]} & Reizlos / intakt \\ \hline
WUNDUMGEBUNG\_GT\_MAPPING & \texttt{["atrophisch-trockene Haut mit trophischen Störungen","CVI-typische Hautveränderungen (Hyperpigmentierung, Atrophie blanche, Lipodermatosklerose)"]} & CVI-typische Hautveränderungen (Hyperpigmentierung, Atrophie blanche, Lipodermatosklerose) \\ \hline
WUNDUMGEBUNG\_GT\_MAPPING & \texttt{["empfindliche Säuglingshaut mit hoher Feuchtigkeits- und Reibungsbelastung"]} & Sonstiges \\ \hline
WUNDUMGEBUNG\_GT\_MAPPING & \texttt{["fragile, atrophe Haut"]} & Sonstiges \\ \hline
WUNDUMGEBUNG\_GT\_MAPPING & \texttt{["gerötet-glänzende Haut, Hinweis auf chronische venöse Stauung/Ödemneigung"]} & Erythem / Rötung, Ödem \\ \hline
WUNDUMGEBUNG\_GT\_MAPPING & \texttt{["gerötet-glänzende, ödematöse Haut","Erythem / Rötung"]} & Erythem / Rötung, Ödem \\ \hline
WUNDUMGEBUNG\_GT\_MAPPING & \texttt{["gut durchbluteter, überwiegend roter Wundgrund"]} & Reizlos / intakt \\ \hline
WUNDUMGEBUNG\_GT\_MAPPING & \texttt{["rocken, teils schuppig, gespannt"]} & Sonstiges \\ \hline
WUNDUMGEBUNG\_GT\_MAPPING & \texttt{["zahlreiche erythematöse Makulae/Papeln, livid-entzündliche Hautveränderungen, Hinweis auf entzündlich-vaskulären Prozess"]} & Ekzem / Dermatitis, Erythem / Rötung \\ \hline
WUNDUMGEBUNG\_GT\_MAPPING & \texttt{["Ödem","CVI-typische Hautveränderungen (Hyperpigmentierung, Atrophie blanche, Lipodermatosklerose)"]} & CVI-typische Hautveränderungen (Hyperpigmentierung, Atrophie blanche, Lipodermatosklerose), Ödem \\ \hline
WUNDUMGEBUNG\_GT\_MAPPING & \texttt{["Ödem","Erythem / Rötung"]} & Erythem / Rötung, Ödem \\ \hline
WUNDUMGEBUNG\_GT\_MAPPING & \texttt{["Ödem","gespannte, trophisch veränderte Haut"]} & Ödem \\ \hline
WUNDUMGEBUNG\_GT\_MAPPING & \texttt{["Ödem"]} & Ödem \\ \hline
WUNDUMGEBUNG\_GT\_MAPPING & \texttt{[]} & nan \\ \hline
DEBRIDEMENT\_GT\_MAPPING & \texttt{Autolytisches Debridement} & Autolytisches Debridement \\ \hline
DEBRIDEMENT\_GT\_MAPPING & \texttt{Chirurgisches Debridement} & Chirurgisches Debridement \\ \hline
DEBRIDEMENT\_GT\_MAPPING & \texttt{Debrisoft Duo} & Debrisoft Duo \\ \hline
DEBRIDEMENT\_GT\_MAPPING & \texttt{Debrisoft Lolly} & Debrisoft Lolly \\ \hline
DEBRIDEMENT\_GT\_MAPPING & \texttt{Debrisoft Pad} & Debrisoft Pad \\ \hline
DEBRIDEMENT\_GT\_MAPPING & \texttt{Ultraschall-assistiertes Debridement (UAW)} & Ultraschall-assistiertes Debridement (UAW) \\ \hline
EXSUDAT\_GT\_MAPPING & \texttt{Annahme stark} & Stark \\ \hline
EXSUDAT\_GT\_MAPPING & \texttt{Große Wunde: stark, da Mazerationen am Wundrand
Kleine Nekrose: trocken} & Stark \\ \hline
EXSUDAT\_GT\_MAPPING & \texttt{Keine} & Keine \\ \hline
EXSUDAT\_GT\_MAPPING & \texttt{Leicht} & Leicht \\ \hline
EXSUDAT\_GT\_MAPPING & \texttt{Mässig} & Mäßig \\ \hline
EXSUDAT\_GT\_MAPPING & \texttt{Mäßig} & Mäßig \\ \hline
EXSUDAT\_GT\_MAPPING & \texttt{Nekrosen: kein Exsudat. Unterschenkel mäßig} & Mäßig \\ \hline
EXSUDAT\_GT\_MAPPING & \texttt{Stark} & Stark \\ \hline
EXSUDAT\_GT\_MAPPING & \texttt{blutig bis serös teilweise vorhanden dann eher mäßig} & Mäßig \\ \hline
EXSUDAT\_GT\_MAPPING & \texttt{eher schwach b is mäßig} & Leicht, Mäßig \\ \hline
EXSUDAT\_GT\_MAPPING & \texttt{eher wenig} & Leicht \\ \hline
EXSUDAT\_GT\_MAPPING & \texttt{eher weniger} & Leicht \\ \hline
EXSUDAT\_GT\_MAPPING & \texttt{gering bis gar nicht} & Keine, Leicht \\ \hline
EXSUDAT\_GT\_MAPPING & \texttt{kein} & Keine \\ \hline
EXSUDAT\_GT\_MAPPING & \texttt{keine} & Keine \\ \hline
EXSUDAT\_GT\_MAPPING & \texttt{keine Angabe möglich} & Enthaltung / keine Angabe \\ \hline
EXSUDAT\_GT\_MAPPING & \texttt{keine Angabe möglich / eher wenig} & Leicht \\ \hline
EXSUDAT\_GT\_MAPPING & \texttt{keine Angabe möglich / vermutlich mässig} & Mäßig \\ \hline
EXSUDAT\_GT\_MAPPING & \texttt{keine Angabe möglich / vermutlich wenig} & Leicht \\ \hline
EXSUDAT\_GT\_MAPPING & \texttt{keine Angabe möglich, vermutlich gering} & Leicht \\ \hline
EXSUDAT\_GT\_MAPPING & \texttt{keine Angabe möglich, vermutlich trocken} & Keine \\ \hline
EXSUDAT\_GT\_MAPPING & \texttt{keine Einschätzung möglich} & Enthaltung / keine Angabe \\ \hline
EXSUDAT\_GT\_MAPPING & \texttt{keine genaue Angabe möglich, sieht eher trocken bis leicht exsudierend aus} & Keine, Leicht \\ \hline
EXSUDAT\_GT\_MAPPING & \texttt{könnte eher schwach sein} & Leicht \\ \hline
EXSUDAT\_GT\_MAPPING & \texttt{leicht} & Leicht \\ \hline
EXSUDAT\_GT\_MAPPING & \texttt{leicht bis mittel} & Leicht, Mäßig \\ \hline
EXSUDAT\_GT\_MAPPING & \texttt{leicht bis mäßig} & Leicht, Mäßig \\ \hline
EXSUDAT\_GT\_MAPPING & \texttt{leichte bis mäßige exsudation} & Leicht, Mäßig \\ \hline
EXSUDAT\_GT\_MAPPING & \texttt{mittel} & Mäßig \\ \hline
EXSUDAT\_GT\_MAPPING & \texttt{mittel bis stark} & Mäßig, Stark \\ \hline
EXSUDAT\_GT\_MAPPING & \texttt{mittelstark} & Mäßig \\ \hline
EXSUDAT\_GT\_MAPPING & \texttt{mäfig} & Mäßig \\ \hline
EXSUDAT\_GT\_MAPPING & \texttt{mäßig} & Mäßig \\ \hline
EXSUDAT\_GT\_MAPPING & \texttt{mäßig bis stark} & Mäßig, Stark \\ \hline
EXSUDAT\_GT\_MAPPING & \texttt{mäßig bis stark, sicher ist eine Geruchsbildung süßlich aromatisch teilweise beschrieben als traubenartig} & Mäßig, Stark \\ \hline
EXSUDAT\_GT\_MAPPING & \texttt{nicht beurteilbar} & Enthaltung / keine Angabe \\ \hline
EXSUDAT\_GT\_MAPPING & \texttt{nicht beurteilbar, vermutlich stark} & Stark \\ \hline
EXSUDAT\_GT\_MAPPING & \texttt{nicht genau definierbar, aus der Erfahrung heraus mäßige Exsudation} & Mäßig \\ \hline
EXSUDAT\_GT\_MAPPING & \texttt{nicht klar, eher trübe Wunde könnte Geruch abgeben} & Enthaltung / keine Angabe \\ \hline
EXSUDAT\_GT\_MAPPING & \texttt{nicht zu beschreiben} & Enthaltung / keine Angabe \\ \hline
EXSUDAT\_GT\_MAPPING & \texttt{offenen stelle mäßig, geschlossene stelle kein bis wenig exsudat} & Keine, Leicht \\ \hline
EXSUDAT\_GT\_MAPPING & \texttt{scheint mäßig zu sein} & Mäßig \\ \hline
EXSUDAT\_GT\_MAPPING & \texttt{schwach} & Leicht \\ \hline
EXSUDAT\_GT\_MAPPING & \texttt{schwach bin mäßig} & Leicht, Mäßig \\ \hline
EXSUDAT\_GT\_MAPPING & \texttt{schwach bis mäßig} & Leicht, Mäßig \\ \hline
EXSUDAT\_GT\_MAPPING & \texttt{schwach bis nicht vorhanden} & Keine, Leicht \\ \hline
EXSUDAT\_GT\_MAPPING & \texttt{sehr gering} & Leicht \\ \hline
EXSUDAT\_GT\_MAPPING & \texttt{stark} & Stark \\ \hline
EXSUDAT\_GT\_MAPPING & \texttt{vermutlich gering} & Leicht \\ \hline
EXSUDAT\_GT\_MAPPING & \texttt{vermutlich hoch} & Stark \\ \hline
EXSUDAT\_GT\_MAPPING & \texttt{vermutlich hohe Exsudation} & Stark \\ \hline
EXSUDAT\_GT\_MAPPING & \texttt{vermutlich kaum} & Keine \\ \hline
EXSUDAT\_GT\_MAPPING & \texttt{vermutlich mittel} & Mäßig \\ \hline
EXSUDAT\_GT\_MAPPING & \texttt{vermutlich mittelmässig} & Mäßig \\ \hline
EXSUDAT\_GT\_MAPPING & \texttt{vermutlich mittelmäßig} & Mäßig \\ \hline
EXSUDAT\_GT\_MAPPING & \texttt{vermutlich mässig} & Mäßig \\ \hline
EXSUDAT\_GT\_MAPPING & \texttt{vermutlich mäßig} & Mäßig \\ \hline
EXSUDAT\_GT\_MAPPING & \texttt{vermutlich starke Exsudation} & Stark \\ \hline
EXSUDAT\_GT\_MAPPING & \texttt{vermutlich vorhanden, in welcher Menge ist nicht zu beurteilen} & Enthaltung / keine Angabe \\ \hline
EXSUDAT\_GT\_MAPPING & \texttt{vermutlich wenig} & Leicht \\ \hline
EXSUDAT\_GT\_MAPPING & \texttt{wahrscheinlich mittelmäßig vorhanden} & Mäßig \\ \hline
EXSUDAT\_GT\_MAPPING & \texttt{warsch hoch bis sehr hoch,  klare exsudation} & Stark \\ \hline
LOKALISATION\_GT\_MAPPING & \texttt{...stark exsudierendes, nekrotisch belegtes Ulcus mit grünlicher Verfärbung typisch für Pseudomonas aeruginosa.} & Enthaltung / keine Angabe \\ \hline
LOKALISATION\_GT\_MAPPING & \texttt{?} & Enthaltung / keine Angabe \\ \hline
LOKALISATION\_GT\_MAPPING & \texttt{???} & Enthaltung / keine Angabe \\ \hline
LOKALISATION\_GT\_MAPPING & \texttt{Abdomen} & Abdomen \\ \hline
LOKALISATION\_GT\_MAPPING & \texttt{Abdomen, links} & Abdomen \\ \hline
LOKALISATION\_GT\_MAPPING & \texttt{Arm} & Arm / Hand \\ \hline
LOKALISATION\_GT\_MAPPING & \texttt{Arm oder Bein??} & Enthaltung / keine Angabe \\ \hline
LOKALISATION\_GT\_MAPPING & \texttt{Aussenseitig linke Fußsohle} & Fuß \\ \hline
LOKALISATION\_GT\_MAPPING & \texttt{Außenknöchel links} & Fuß \\ \hline
LOKALISATION\_GT\_MAPPING & \texttt{Bein} & Bein \\ \hline
LOKALISATION\_GT\_MAPPING & \texttt{Bein nicht genau definierbar} & Bein \\ \hline
LOKALISATION\_GT\_MAPPING & \texttt{Bein wo genau fraglich} & Bein \\ \hline
LOKALISATION\_GT\_MAPPING & \texttt{Bein, Unterschenkel Richtung Fuß} & Bein \\ \hline
LOKALISATION\_GT\_MAPPING & \texttt{Bein, evtl. links, Unterschenkel} & Bein \\ \hline
LOKALISATION\_GT\_MAPPING & \texttt{Bein, evtl. rechts, Unterschenkel} & Bein \\ \hline
LOKALISATION\_GT\_MAPPING & \texttt{Bein, genaue Definition nicht möglich} & Bein \\ \hline
LOKALISATION\_GT\_MAPPING & \texttt{Bein, genauere Lokalisation nicht genau definierbar} & Bein \\ \hline
LOKALISATION\_GT\_MAPPING & \texttt{Bein, wo genau nicht definierbar} & Bein \\ \hline
LOKALISATION\_GT\_MAPPING & \texttt{Das Bild zeigt ein großflächiges venöses Ulcus cruris bei ausgeprägter chronisch venöser Stauung und Adipositas-assoziierter Belastung. Möglich, Oberschenkel links, Unterschenkel} & Bein \\ \hline
LOKALISATION\_GT\_MAPPING & \texttt{Ferse} & Fuß \\ \hline
LOKALISATION\_GT\_MAPPING & \texttt{Ferse, Achillessehne} & Fuß \\ \hline
LOKALISATION\_GT\_MAPPING & \texttt{Ferse, linker Fuß} & Fuß \\ \hline
LOKALISATION\_GT\_MAPPING & \texttt{Ferse, rechter Fuß} & Fuß \\ \hline
LOKALISATION\_GT\_MAPPING & \texttt{Ferse, scheint links zu sein:
gelblich-fibrinösem Belag
mäßiger Exsudation
gerötetem Wundrand} & Fuß \\ \hline
LOKALISATION\_GT\_MAPPING & \texttt{Fuß lateral links} & Fuß \\ \hline
LOKALISATION\_GT\_MAPPING & \texttt{Fuß links plantar} & Fuß \\ \hline
LOKALISATION\_GT\_MAPPING & \texttt{Fuß links plantar  lateral} & Fuß \\ \hline
LOKALISATION\_GT\_MAPPING & \texttt{Fuß rechts} & Fuß \\ \hline
LOKALISATION\_GT\_MAPPING & \texttt{Fuß, Ferse links} & Fuß \\ \hline
LOKALISATION\_GT\_MAPPING & \texttt{Fußrücken} & Fuß \\ \hline
LOKALISATION\_GT\_MAPPING & \texttt{Fußrücken linker Fuß} & Fuß \\ \hline
LOKALISATION\_GT\_MAPPING & \texttt{Fußrücken rechter Fuß} & Fuß \\ \hline
LOKALISATION\_GT\_MAPPING & \texttt{Fußsohle Ferse} & Fuß \\ \hline
LOKALISATION\_GT\_MAPPING & \texttt{Gesäß-/Sakralregion} & Gesäß / Sakral \\ \hline
LOKALISATION\_GT\_MAPPING & \texttt{Gesäß-/perianale Region eines Säuglings} & Gesäß / Sakral \\ \hline
LOKALISATION\_GT\_MAPPING & \texttt{Knöchelinnenseite} & Fuß \\ \hline
LOKALISATION\_GT\_MAPPING & \texttt{Lokalisation: Fußrücken, lateraler Vorfuß sowie Zehenbereiche} & Fuß \\ \hline
LOKALISATION\_GT\_MAPPING & \texttt{Lokalisation: Unterschenkelbereich} & Bein \\ \hline
LOKALISATION\_GT\_MAPPING & \texttt{Lokalisation: Unterschenkelregion} & Bein \\ \hline
LOKALISATION\_GT\_MAPPING & \texttt{Lokalisation: lateraler Mittelfuß/Knöchelbereich} & Fuß \\ \hline
LOKALISATION\_GT\_MAPPING & \texttt{Lokalisation: lateraler Vor-/Mittelfuß} & Fuß \\ \hline
LOKALISATION\_GT\_MAPPING & \texttt{Lokalisation: sakrogluteale/perianale Region beidseits} & Gesäß / Sakral \\ \hline
LOKALISATION\_GT\_MAPPING & \texttt{Lokalisation: vermutlich Sakral-/Gesäßregion} & Gesäß / Sakral \\ \hline
LOKALISATION\_GT\_MAPPING & \texttt{Oberarb Muskulus bizeps} & Arm / Hand \\ \hline
LOKALISATION\_GT\_MAPPING & \texttt{Os Sacrum} & Gesäß / Sakral \\ \hline
LOKALISATION\_GT\_MAPPING & \texttt{Rima Ani} & Gesäß / Sakral \\ \hline
LOKALISATION\_GT\_MAPPING & \texttt{Sakralbereich} & Gesäß / Sakral \\ \hline
LOKALISATION\_GT\_MAPPING & \texttt{Sakralbereich linke hälfte} & Gesäß / Sakral \\ \hline
LOKALISATION\_GT\_MAPPING & \texttt{Sakraler Bereich} & Gesäß / Sakral \\ \hline
LOKALISATION\_GT\_MAPPING & \texttt{Steiß} & Gesäß / Sakral \\ \hline
LOKALISATION\_GT\_MAPPING & \texttt{Unterarm} & Arm / Hand \\ \hline
LOKALISATION\_GT\_MAPPING & \texttt{Unterschenkel} & Bein \\ \hline
LOKALISATION\_GT\_MAPPING & \texttt{Unterschenkel  medial} & Bein \\ \hline
LOKALISATION\_GT\_MAPPING & \texttt{Unterschenkel warsch. gamaschenartig / bereits zirkulierend} & Bein \\ \hline
LOKALISATION\_GT\_MAPPING & \texttt{Unterschenkel weit umfassend} & Bein \\ \hline
LOKALISATION\_GT\_MAPPING & \texttt{Unterschenkel, Außenseite auf Knie Höhe} & Bein \\ \hline
LOKALISATION\_GT\_MAPPING & \texttt{Unterschenkel, rechts, Innenseite} & Bein \\ \hline
LOKALISATION\_GT\_MAPPING & \texttt{Unterschenkel/Knöchelregion, zirkumferenziell ausgedehnt} & Bein \\ \hline
LOKALISATION\_GT\_MAPPING & \texttt{Unterschenkel?} & Bein \\ \hline
LOKALISATION\_GT\_MAPPING & \texttt{chronisches Ulkus im Bereich der Achillessehne
dorsaler Unterschenkel / Achillessehnenregion} & Bein \\ \hline
LOKALISATION\_GT\_MAPPING & \texttt{ferse, fußsohle} & Fuß \\ \hline
LOKALISATION\_GT\_MAPPING & \texttt{fibrinösem gelblich-weißem Belag,
mäßiger Exsudation,
gerötetem Wundrand,
--> vermutlich entzündlicher bzw. kritisch kolonisierter Situation,} & Enthaltung / keine Angabe \\ \hline
LOKALISATION\_GT\_MAPPING & \texttt{größeres, flächiges Ulcus am Unterschenkel (Wade) bei ausgeprägtem Ödem und Adipositas-assoziierter Hautveränderung.} & Bein \\ \hline
LOKALISATION\_GT\_MAPPING & \texttt{keine Angabe möglich} & Enthaltung / keine Angabe \\ \hline
LOKALISATION\_GT\_MAPPING & \texttt{könnte Innenseite linker Fuß sein} & Fuß \\ \hline
LOKALISATION\_GT\_MAPPING & \texttt{könnte Unterschenkel oder Handgelenk sein} & Enthaltung / keine Angabe \\ \hline
LOKALISATION\_GT\_MAPPING & \texttt{könnte am linken Fuß sein... Knöchelbereich} & Fuß \\ \hline
LOKALISATION\_GT\_MAPPING & \texttt{könnte der Unterschenkel sein, rechts} & Bein \\ \hline
LOKALISATION\_GT\_MAPPING & \texttt{könnte oberhalb des os sacrum liegen} & Gesäß / Sakral \\ \hline
LOKALISATION\_GT\_MAPPING & \texttt{könnte sich im Bereich des Steißes befinden} & Gesäß / Sakral \\ \hline
LOKALISATION\_GT\_MAPPING & \texttt{lateral linker Fuß} & Fuß \\ \hline
LOKALISATION\_GT\_MAPPING & \texttt{li lateraler Malleolus} & Fuß \\ \hline
LOKALISATION\_GT\_MAPPING & \texttt{li. Fußrücken} & Fuß \\ \hline
LOKALISATION\_GT\_MAPPING & \texttt{linke Ferse} & Fuß \\ \hline
LOKALISATION\_GT\_MAPPING & \texttt{linker Fuss} & Fuß \\ \hline
LOKALISATION\_GT\_MAPPING & \texttt{linker Fussrücken} & Fuß \\ \hline
LOKALISATION\_GT\_MAPPING & \texttt{linker Fuß} & Fuß \\ \hline
LOKALISATION\_GT\_MAPPING & \texttt{linker Fuß
Lokalisation: dorsaler und lateraler Vorfuß mit Beteiligung der Zehenbasis} & Fuß \\ \hline
LOKALISATION\_GT\_MAPPING & \texttt{linker Fuß - dorsaler Vorfuß/Fußrücken} & Fuß \\ \hline
LOKALISATION\_GT\_MAPPING & \texttt{linker Fuß außen} & Fuß \\ \hline
LOKALISATION\_GT\_MAPPING & \texttt{linker Fuß lateratl} & Fuß \\ \hline
LOKALISATION\_GT\_MAPPING & \texttt{linker Fuß medial, malleolär} & Fuß \\ \hline
LOKALISATION\_GT\_MAPPING & \texttt{linker Fuß medial-lateral 
kleine Nekrose an der Ferse außen plantar} & Fuß \\ \hline
LOKALISATION\_GT\_MAPPING & \texttt{linker Fuß, Außenseite} & Fuß \\ \hline
LOKALISATION\_GT\_MAPPING & \texttt{linker Fuß, Ferse} & Fuß \\ \hline
LOKALISATION\_GT\_MAPPING & \texttt{linker Fuß, Innenseite} & Fuß \\ \hline
LOKALISATION\_GT\_MAPPING & \texttt{linker Fuß, Mittelfußsohle, Ferse} & Fuß \\ \hline
LOKALISATION\_GT\_MAPPING & \texttt{linker Fuß, Spann, Fußrücken} & Fuß \\ \hline
LOKALISATION\_GT\_MAPPING & \texttt{linker Fuß, außenseite} & Fuß \\ \hline
LOKALISATION\_GT\_MAPPING & \texttt{linker Fuß, seite} & Fuß \\ \hline
LOKALISATION\_GT\_MAPPING & \texttt{linker Fußrücken} & Fuß \\ \hline
LOKALISATION\_GT\_MAPPING & \texttt{linker Oberarm, Innenseite} & Arm / Hand \\ \hline
LOKALISATION\_GT\_MAPPING & \texttt{linker Oberarm/ Innenseitig} & Arm / Hand \\ \hline
LOKALISATION\_GT\_MAPPING & \texttt{linker proximaler Unterschenkel lateral} & Bein \\ \hline
LOKALISATION\_GT\_MAPPING & \texttt{linkes Bein unterhalb Knie lateral} & Bein \\ \hline
LOKALISATION\_GT\_MAPPING & \texttt{lokalisation nicht genau zu definieren} & Enthaltung / keine Angabe \\ \hline
LOKALISATION\_GT\_MAPPING & \texttt{mediale Fußseite linker Fuß} & Fuß \\ \hline
LOKALISATION\_GT\_MAPPING & \texttt{mehrere kleine bis mittelgroße Ulzerationen am Unterschenkel (Wade, rechts?) mit entzündlich wirkender Umgebungshaut.} & Bein \\ \hline
LOKALISATION\_GT\_MAPPING & \texttt{multiple Schädigung gesamter Gesäßbereich} & Gesäß / Sakral \\ \hline
LOKALISATION\_GT\_MAPPING & \texttt{nicht beurteilbar} & Enthaltung / keine Angabe \\ \hline
LOKALISATION\_GT\_MAPPING & \texttt{nicht genau definierbar} & Enthaltung / keine Angabe \\ \hline
LOKALISATION\_GT\_MAPPING & \texttt{nicht genau definierbar könnte Sakralbereich sein} & Gesäß / Sakral \\ \hline
LOKALISATION\_GT\_MAPPING & \texttt{nicht genau definierbar sieht für mich aus wie am Vorfuß} & Fuß \\ \hline
LOKALISATION\_GT\_MAPPING & \texttt{nicht genau definierbar, Vermutung Unterschenkel} & Bein \\ \hline
LOKALISATION\_GT\_MAPPING & \texttt{nicht genau definierbar, meist am Unterschenkel} & Bein \\ \hline
LOKALISATION\_GT\_MAPPING & \texttt{oberhalb der Achillessehne} & Fuß \\ \hline
LOKALISATION\_GT\_MAPPING & \texttt{plantarseitig rechter Fuß} & Fuß \\ \hline
LOKALISATION\_GT\_MAPPING & \texttt{re Fuß} & Fuß \\ \hline
LOKALISATION\_GT\_MAPPING & \texttt{rechte Ferse} & Fuß \\ \hline
LOKALISATION\_GT\_MAPPING & \texttt{rechte Hand} & Arm / Hand \\ \hline
LOKALISATION\_GT\_MAPPING & \texttt{rechte Hand und Arm} & Arm / Hand \\ \hline
LOKALISATION\_GT\_MAPPING & \texttt{rechter Fuß} & Fuß \\ \hline
LOKALISATION\_GT\_MAPPING & \texttt{rechter Fuß - lateraler und plantarer Fußbereich/Ferse} & Fuß \\ \hline
LOKALISATION\_GT\_MAPPING & \texttt{rechter Fuß Fußrücken und lateral} & Fuß \\ \hline
LOKALISATION\_GT\_MAPPING & \texttt{rechter Fuß medial, Richtung plantar.} & Fuß \\ \hline
LOKALISATION\_GT\_MAPPING & \texttt{rechter Fuß plantar} & Fuß \\ \hline
LOKALISATION\_GT\_MAPPING & \texttt{rechter Fuß, Ferse} & Fuß \\ \hline
LOKALISATION\_GT\_MAPPING & \texttt{rechter Fuß, Ferse, Mittelfuß - Fußsohle} & Fuß \\ \hline
LOKALISATION\_GT\_MAPPING & \texttt{rechter Fuß, Fußsohle} & Fuß \\ \hline
LOKALISATION\_GT\_MAPPING & \texttt{rechter Fuß, Innenseite, Knöchel} & Fuß \\ \hline
LOKALISATION\_GT\_MAPPING & \texttt{rechter Fuß, Spann} & Fuß \\ \hline
LOKALISATION\_GT\_MAPPING & \texttt{rechter Fuß, in Höhe der großen Zehe, unterhalb der medialen Malleolen} & Fuß \\ \hline
LOKALISATION\_GT\_MAPPING & \texttt{rechter Unterarm, rechtes Handgelenk, rechte Hand} & Arm / Hand \\ \hline
LOKALISATION\_GT\_MAPPING & \texttt{rechter Unterschenkel} & Bein \\ \hline
LOKALISATION\_GT\_MAPPING & \texttt{rechtes Bein, Kniebereich, Unterschenkel} & Bein \\ \hline
LOKALISATION\_GT\_MAPPING & \texttt{rechtes Bein, höhe der Knöchel} & Bein \\ \hline
LOKALISATION\_GT\_MAPPING & \texttt{sacralbereich} & Gesäß / Sakral \\ \hline
LOKALISATION\_GT\_MAPPING & \texttt{sakrales Dekubitalulkus} & Gesäß / Sakral \\ \hline
LOKALISATION\_GT\_MAPPING & \texttt{schwer zu beurteilen, Vermutung am Trochanter major} & Bein \\ \hline
LOKALISATION\_GT\_MAPPING & \texttt{schwer zu beurteilen, keine genaue Angabe möglich,
eventuell Rücken Richtung Oberarm /} & Enthaltung / keine Angabe \\ \hline
LOKALISATION\_GT\_MAPPING & \texttt{schwierig zu beantworten... könnte die Wade sein, evtl. rechts} & Bein \\ \hline
LOKALISATION\_GT\_MAPPING & \texttt{unterschenkel} & Bein \\ \hline
LOKALISATION\_GT\_MAPPING & \texttt{vermutlich am Außenknöchel} & Fuß \\ \hline
LOKALISATION\_GT\_MAPPING & \texttt{vermutlich am Unterschenkel} & Bein \\ \hline
LOKALISATION\_KEYWORDS & \texttt{Abdomen} & abdomen, bauch, peristom, stoma \\ \hline
LOKALISATION\_KEYWORDS & \texttt{Fuß} & fuß, fuss, zehe, ferse, fußsohle, fußrücken, plantar, vorfuß, rückfuß, malleol, achillessehne, calcaneus, hallux, metatarsal \\ \hline
LOKALISATION\_KEYWORDS & \texttt{Bein} & bein, knöchel, unterschenkel, oberschenkel, knie, gaiter \\ \hline
LOKALISATION\_KEYWORDS & \texttt{Hand} & \bhand\b, \bhände\b, finger, handgelenk, handrücken \\ \hline
LOKALISATION\_KEYWORDS & \texttt{Arm} & arm, ellenbogen, oberarm, unterarm, bizeps, axilla, achsel \\ \hline
LOKALISATION\_KEYWORDS & \texttt{Gesäß} & gesäß, sakral, kreuzbein, steißbein, os sacrum, sacral, rima ani, steiß, trochanter, gluteal, paraglutäal \\ \hline
LOKALISATION\_KEYWORDS & \texttt{Rücken} & (?<!fuß)(?<!fuss)(?<!hand)(?<!händ)rücken \\ \hline
LOKALISATION\_KEYWORDS & \texttt{Flanke} & flanke \\ \hline
LOKALISATION\_KEYWORDS & \texttt{Kopf} & \bkopf\b, hals, nacken \\ \hline
LOKALISATION\_KEYWORDS & \texttt{Brust} & brust, thorax \\ \hline
PRODUKT\_GT\_MAPPING & \texttt{Actico UlcerSys System} & Actico UlcerSys System \\ \hline
PRODUKT\_GT\_MAPPING & \texttt{ReadyWrap Untere Extremität} & ReadyWrap Untere Extremität \\ \hline
PRODUKT\_GT\_MAPPING & \texttt{Rosidal TCS} & Rosidal TCS \\ \hline
PRODUKT\_GT\_MAPPING & \texttt{Rosidal soft} & Rosidal soft \\ \hline
PRODUKT\_GT\_MAPPING & \texttt{Suprasorb P SensiFlex multisite border lite} & Suprasorb P Sensiflex \\ \hline
PRODUKT\_GT\_MAPPING & \texttt{Suprasorb P selbstklebend} & Suprasorb P \\ \hline
PRODUKT\_GT\_MAPPING & \texttt{Suprasorb P sensitive (finger/toe)} & Suprasorb P Sensitive \\ \hline
PRODUKT\_GT\_MAPPING & \texttt{Suprasorb P sensitive (sacrum)} & Suprasorb P Sensitive \\ \hline
PRODUKT\_GT\_MAPPING & \texttt{Suprasorb X + PHMB Tamponade} & Suprasorb X + PHMB \\ \hline
PRODUKT\_GT\_MAPPING & \texttt{Vliwasorb sensitive border sacrum} & Vliwasorb sensitive border \\ \hline
PRODUKT\_GT\_MAPPING & \texttt{tg Fertigverbände} & tg Fertigverbände \\ \hline
PRODUKT\_GT\_MAPPING & \texttt{tg Fertigverbände (Hand-/Fußverband)} & tg Fertigverbände \\ \hline
PRODUKT\_GT\_MAPPING & \texttt{tg Fertigverbände Hand-/Fußverband} & tg Fertigverbände \\ \hline
PRODUKT\_GT\_MAPPING & \texttt{Suprasorb A + Ag} & Suprasorb A + Ag \\ \hline
PRODUKT\_GT\_MAPPING & \texttt{Suprasorb A + Ag (Kompresse)} & Suprasorb A + Ag \\ \hline
PRODUKT\_GT\_MAPPING & \texttt{Suprasorb A + Ag Kompresse} & Suprasorb A + Ag \\ \hline
PRODUKT\_GT\_MAPPING & \texttt{Suprasorb A + Ag Tamponade} & Suprasorb A + Ag \\ \hline
PRODUKT\_GT\_MAPPING & \texttt{Suprasorb A Pro} & Suprasorb A Pro \\ \hline
PRODUKT\_GT\_MAPPING & \texttt{Suprasorb A Pro (Kompresse)} & Suprasorb A Pro \\ \hline
PRODUKT\_GT\_MAPPING & \texttt{Suprasorb A Pro (Tamponade)} & Suprasorb A Pro \\ \hline
PRODUKT\_GT\_MAPPING & \texttt{Suprasorb A Pro Kompresse} & Suprasorb A Pro \\ \hline
PRODUKT\_GT\_MAPPING & \texttt{Suprasorb A Pro Tamponade} & Suprasorb A Pro \\ \hline
PRODUKT\_GT\_MAPPING & \texttt{Suprasorb Liquacel Pro} & Suprasorb Liquacel Pro \\ \hline
PRODUKT\_GT\_MAPPING & \texttt{Suprasorb Liquacel Pro (Kompresse)} & Suprasorb Liquacel Pro \\ \hline
PRODUKT\_GT\_MAPPING & \texttt{Suprasorb Liquacel Pro (Tamponade)} & Suprasorb Liquacel Pro \\ \hline
PRODUKT\_GT\_MAPPING & \texttt{Suprasorb Liquacel Pro Kompresse} & Suprasorb Liquacel Pro \\ \hline
PRODUKT\_GT\_MAPPING & \texttt{Suprasorb Liquacel Pro Tamponade} & Suprasorb Liquacel Pro \\ \hline
PRODUKT\_GT\_MAPPING & \texttt{Suprasorb P} & Suprasorb P \\ \hline
PRODUKT\_GT\_MAPPING & \texttt{Suprasorb P (nicht klebend)} & Suprasorb P \\ \hline
PRODUKT\_GT\_MAPPING & \texttt{Suprasorb P (selbstklebend)} & Suprasorb P \\ \hline
PRODUKT\_GT\_MAPPING & \texttt{Suprasorb P (self-adhesive, sacrum)} & Suprasorb P \\ \hline
PRODUKT\_GT\_MAPPING & \texttt{Suprasorb P nicht klebend} & Suprasorb P \\ \hline
PRODUKT\_GT\_MAPPING & \texttt{Suprasorb P heel (selbstklebend)} & Suprasorb P \\ \hline
PRODUKT\_GT\_MAPPING & \texttt{Suprasorb P + PHMB} & Suprasorb P + PHMB \\ \hline
PRODUKT\_GT\_MAPPING & \texttt{Suprasorb P SensiFlex (border rechteckig)} & Suprasorb P Sensiflex \\ \hline
PRODUKT\_GT\_MAPPING & \texttt{Suprasorb P SensiFlex border} & Suprasorb P Sensiflex \\ \hline
PRODUKT\_GT\_MAPPING & \texttt{Suprasorb P SensiFlex border rechteckig} & Suprasorb P Sensiflex \\ \hline
PRODUKT\_GT\_MAPPING & \texttt{Suprasorb P SensiFlex multisite border} & Suprasorb P Sensiflex \\ \hline
PRODUKT\_GT\_MAPPING & \texttt{Suprasorb P Sensiflex} & Suprasorb P Sensiflex \\ \hline
PRODUKT\_GT\_MAPPING & \texttt{Suprasorb P sensitive} & Suprasorb P Sensitive \\ \hline
PRODUKT\_GT\_MAPPING & \texttt{Suprasorb P sensitive (heel)} & Suprasorb P Sensitive \\ \hline
PRODUKT\_GT\_MAPPING & \texttt{Suprasorb P sensitive (nicht klebend)} & Suprasorb P Sensitive \\ \hline
PRODUKT\_GT\_MAPPING & \texttt{Suprasorb P sensitive (selbstklebend)} & Suprasorb P Sensitive \\ \hline
PRODUKT\_GT\_MAPPING & \texttt{Suprasorb P sensitive (selbstklebend) als abdeckender Verband} & Suprasorb P Sensitive \\ \hline
PRODUKT\_GT\_MAPPING & \texttt{Suprasorb P sensitive heel} & Suprasorb P Sensitive \\ \hline
PRODUKT\_GT\_MAPPING & \texttt{Suprasorb P sensitive heel (selbstklebend) – als polsternde, absorbierende Sekundärlage über der Tamponade} & Suprasorb P Sensitive \\ \hline
PRODUKT\_GT\_MAPPING & \texttt{Suprasorb P sensitive sacrum} & Suprasorb P Sensitive \\ \hline
PRODUKT\_GT\_MAPPING & \texttt{Suprasorb P sensitive selbstklebend} & Suprasorb P Sensitive \\ \hline
PRODUKT\_GT\_MAPPING & \texttt{Suprasorb X} & Suprasorb X \\ \hline
PRODUKT\_GT\_MAPPING & \texttt{Suprasorb X Kompresse} & Suprasorb X \\ \hline
PRODUKT\_GT\_MAPPING & \texttt{Suprasorb X Pro} & Suprasorb X Pro \\ \hline
PRODUKT\_GT\_MAPPING & \texttt{Suprasorb X Pro (Kompresse)} & Suprasorb X Pro \\ \hline
PRODUKT\_GT\_MAPPING & \texttt{Suprasorb X Pro Kompresse} & Suprasorb X Pro \\ \hline
PRODUKT\_GT\_MAPPING & \texttt{Suprasorb X + PHMB} & Suprasorb X + PHMB \\ \hline
PRODUKT\_GT\_MAPPING & \texttt{Suprasorb X + PHMB (Kompresse)} & Suprasorb X + PHMB \\ \hline
PRODUKT\_GT\_MAPPING & \texttt{Suprasorb X + PHMB Kompresse} & Suprasorb X + PHMB \\ \hline
PRODUKT\_GT\_MAPPING & \texttt{Suprasorb G Gel-Kompresse} & Suprasorb G Gel-Kompresse \\ \hline
PRODUKT\_GT\_MAPPING & \texttt{Suprasorb H} & Suprasorb H \\ \hline
PRODUKT\_GT\_MAPPING & \texttt{Suprasorb F} & Suprasorb F \\ \hline
PRODUKT\_GT\_MAPPING & \texttt{Suprasorb F Protect} & Suprasorb F Protect \\ \hline
PRODUKT\_GT\_MAPPING & \texttt{Lomatuell Pro} & Lomatuell Pro \\ \hline
PRODUKT\_GT\_MAPPING & \texttt{Lomatuell H} & Lomatuell H \\ \hline
PRODUKT\_GT\_MAPPING & \texttt{Vliwasorb Pro} & Vliwasorb Pro \\ \hline
PRODUKT\_GT\_MAPPING & \texttt{Vliwasorb sensitive} & Vliwasorb sensitive \\ \hline
PRODUKT\_GT\_MAPPING & \texttt{Vliwazell} & Vliwazell \\ \hline
PRODUKT\_GT\_MAPPING & \texttt{Vliwazell Pro} & Vliwazell Pro \\ \hline
PRODUKT\_GT\_MAPPING & \texttt{Vliwazell Pro als abdeckender Verband} & Vliwazell Pro \\ \hline
PRODUKT\_GT\_MAPPING & \texttt{Vliwazell Pro – als hochabsorbierende Sekundärlage über der Tamponade} & Vliwazell Pro \\ \hline
PRODUKT\_GT\_MAPPING & \texttt{Vliwaktiv Ag} & Vliwaktiv Ag \\ \hline
PRODUKT\_GT\_MAPPING & \texttt{Vliwaktiv Ag Saugkompresse} & Vliwaktiv Ag \\ \hline
PRODUKT\_GT\_MAPPING & \texttt{Vliwaktiv Ag Tamponade} & Vliwaktiv Ag \\ \hline
PRODUKT\_GT\_MAPPING & \texttt{Solvaline N} & Solvaline N \\ \hline
PRODUKT\_GT\_MAPPING & \texttt{Metalline Kompresse} & Metalline Kompresse \\ \hline
PRODUKT\_GT\_MAPPING & \texttt{Curafix H} & Curafix H \\ \hline
PRODUKT\_GT\_MAPPING & \texttt{Porofix} & Porofix \\ \hline
PRODUKT\_GT\_MAPPING & \texttt{Silkafix} & Silkafix \\ \hline
PRODUKT\_GT\_MAPPING & \texttt{Mollelast} & Mollelast \\ \hline
PRODUKT\_GT\_MAPPING & \texttt{Mollelast (Elastische Fixierbinde)} & Mollelast \\ \hline
PRODUKT\_GT\_MAPPING & \texttt{Mollelast haft latexfrei} & Mollelast haft latexfrei \\ \hline
PRODUKT\_GT\_MAPPING & \texttt{Haftelast latexfrei} & Haftelast latexfrei \\ \hline
PRODUKT\_GT\_MAPPING & \texttt{Haftelast latexfrei (kohäsive Fixierbinde)} & Haftelast latexfrei \\ \hline
PRODUKT\_GT\_MAPPING & \texttt{ActiFast} & ActiFast \\ \hline
PRODUKT\_GT\_MAPPING & \texttt{ActiFast Schlauchverband} & ActiFast \\ \hline
PRODUKT\_GT\_MAPPING & \texttt{tg Schlauchverband} & tg Schlauchverband \\ \hline
PRODUKT\_GT\_MAPPING & \texttt{Curapor} & Curapor \\ \hline
PRODUKT\_GT\_MAPPING & \texttt{Curapor transparent} & Curapor transparent \\ \hline
PRODUKT\_GT\_MAPPING & \texttt{Fixierbinde} & Fixierbinde \\ \hline
PRODUKT\_GT\_MAPPING & \texttt{Fixierbinden} & Fixierbinde \\ \hline
PRODUKT\_GT\_MAPPING & \texttt{Universalbinde} & Universalbinde \\ \hline
PRODUKT\_GT\_MAPPING & \texttt{Universalbinden} & Universalbinde \\ \hline
PRODUKT\_GT\_MAPPING & \texttt{Vliwasorb sensitive border} & Vliwasorb sensitive border \\ \hline
PRODUKT\_GT\_MAPPING & \texttt{Vliwasorb sensitive border (sacrum)} & Vliwasorb sensitive border \\ \hline
PRODUKT\_GT\_MAPPING & \texttt{Suprasorb P sensitive (Heel, selbstklebend)} & Suprasorb P Sensitive \\ \hline
PRODUKT\_GT\_MAPPING & \texttt{Suprasorb P sensitive (selbstklebend, sacrum)} & Suprasorb P Sensitive \\ \hline
PRODUKT\_GT\_MAPPING & \texttt{Suprasorb P (Heel)} & Suprasorb P \\ \hline
WUNDGRUND\_GT\_MAPPING & \texttt{nekrotisch} & Nekrose \\ \hline
WUNDGRUND\_GT\_MAPPING & \texttt{Fibrinbelegt} & Fibrinbelag \\ \hline
WUNDGRUND\_GT\_MAPPING & \texttt{Sauber, bis auf die aufgetragene Salbe} & Sauber \\ \hline
WUNDGRUND\_GT\_MAPPING & \texttt{Fibrinbelegte Wunde} & Fibrinbelag \\ \hline
WUNDGRUND\_GT\_MAPPING & \texttt{nekrotisch belegt} & Nekrose, Fibrinbelag \\ \hline
WUNDGRUND\_GT\_MAPPING & \texttt{Fibringbelegt mit durschscheinendem Granulationsgewebe} & Fibrinbelag, Granulation \\ \hline
WUNDGRUND\_GT\_MAPPING & \texttt{Mischuntergrund. Fibrinbelag, nekrotisch, teils mit Granulationsinseln} & Mischuntergrund, Fibrinbelag, Nekrose, Granulation \\ \hline
WUNDGRUND\_GT\_MAPPING & \texttt{teilweise fibrinbelegt mit einzelnen Granulationsinseln} & Fibrinbelag, Granulation \\ \hline
WUNDGRUND\_GT\_MAPPING & \texttt{sauberer Wundgrund mit Inseln von Granulation} & Sauber, Granulation \\ \hline
WUNDGRUND\_GT\_MAPPING & \texttt{Fibrinbelegt mit einzelnen Granulationsinseln} & Fibrinbelag, Granulation \\ \hline
WUNDGRUND\_GT\_MAPPING & \texttt{Fibrinbelegtes tiefes Ulcus} & Fibrinbelag, Ulzeration \\ \hline
WUNDGRUND\_GT\_MAPPING & \texttt{Fibrinbelegt bei Wundrandepithelisierung} & Fibrinbelag, Wundrandepithelisierung \\ \hline
WUNDGRUND\_GT\_MAPPING & \texttt{belegt mit Hypergranulation} & Hypergranulation \\ \hline
WUNDGRUND\_GT\_MAPPING & \texttt{Fibrinbelegte Wundverhältnisse} & Fibrinbelag \\ \hline
WUNDGRUND\_GT\_MAPPING & \texttt{Fibrinbelegt, zerklüftet} & Fibrinbelag, Zerklüftet \\ \hline
WUNDGRUND\_GT\_MAPPING & \texttt{Sauberer Wundgrund mit kleinen Fibrininseln} & Sauber, Fibrinbelag \\ \hline
WUNDGRUND\_GT\_MAPPING & \texttt{belegt, teilweise nekrotisch} & Belag, Nekrose \\ \hline
WUNDGRUND\_GT\_MAPPING & \texttt{stark infiziert, massiv belegt} & Infektion, Belag \\ \hline
WUNDGRUND\_GT\_MAPPING & \texttt{verbrannte  Epidermis} & Verbrennung \\ \hline
WUNDGRUND\_GT\_MAPPING & \texttt{belegt, nekrotisch} & Belag, Nekrose \\ \hline
WUNDGRUND\_GT\_MAPPING & \texttt{massive großflächige Verbrennung teilweise 4. Grades, je mehr es vom Wundzentrum nach außen geht werden Verbrennungen 3 7 2. und 1. Grades sichtbar. Nach der 9 er regel könnten hier bis zu 9 \% der Körperoberfläche geschädigt sein} & Verbrennung \\ \hline
WUNDGRUND\_GT\_MAPPING & \texttt{Verbrennungen, Blasenbildung, offene Wunden keine verbrannten Sehnen / Knochen  epidermis geschädigt} & Verbrennung, Blasenbildung \\ \hline
WUNDGRUND\_GT\_MAPPING & \texttt{stark belgtes Ulcus, teilweise einzelne Granulationsinseln,} & Belag, Granulation \\ \hline
WUNDGRUND\_GT\_MAPPING & \texttt{hämatös} & Hämatom \\ \hline
WUNDGRUND\_GT\_MAPPING & \texttt{Fußwunde: Nekrotisch  Unterschenkelwunde fibrinös belegt,} & Nekrose, Fibrinbelag \\ \hline
WUNDGRUND\_GT\_MAPPING & \texttt{sehr ausgeprägte massive Nekrose, teilweise frei liegende Strukturen, blutig seröses Exsudat} & Nekrose, Freiliegende Strukturen, Exsudat \\ \hline
WUNDGRUND\_GT\_MAPPING & \texttt{belegte Wunde, teilweise infiziert} & Belag, Infektion \\ \hline
WUNDGRUND\_GT\_MAPPING & \texttt{Nekrotische und fibrinös belegte Wunde- auf Grund der Mazeration nach plantar hin könnte mäßige Exsudation auftreten. Wundinfektion möglich} & Nekrose, Fibrinbelag, Mazeration \\ \hline
WUNDGRUND\_GT\_MAPPING & \texttt{fibrinös belegt, zum Teil mit Granulationsinseln, teilweise freiliegende Sehnen.
Nekrose plantar wandständig trocken} & Fibrinbelag, Granulation, Freiliegende Sehnen, Nekrose \\ \hline
WUNDGRUND\_GT\_MAPPING & \texttt{mehrere Stadien zu erkennen, nekrotisch, belegt, teilweise granuliert} & Nekrose, Belag, Granulation \\ \hline
WUNDGRUND\_GT\_MAPPING & \texttt{Nekrotisch und durch Infektion belegt, Wunde bildet Geruch, Frage der Keimbestimmung wichtig, Wundabstrich notwendig,} & Nekrose, Infektion \\ \hline
WUNDGRUND\_GT\_MAPPING & \texttt{teilweise nekrotisch, belegt} & Nekrose, Belag \\ \hline
WUNDGRUND\_GT\_MAPPING & \texttt{offenen stelle entzündet. belegt, nekrotisch, geschlossene stelle nekrotisch} & Entzündung, Belag, Nekrose \\ \hline
WUNDGRUND\_GT\_MAPPING & \texttt{wandständige Nekrose, Prüfung, ob Knochen / Sehnen betroffen sind. Unbedingt chirurgisch intervenieren.
Druckentlastung essentiell.} & Nekrose \\ \hline
WUNDGRUND\_GT\_MAPPING & \texttt{fibrinös belegt, unterminierte Wundränder siehe punktierte Markierung
Eventuelle freie Sehnen} & Fibrinbelag, Unterminierung, Freiliegende Sehnen \\ \hline
WUNDGRUND\_GT\_MAPPING & \texttt{freiliegende Sehnen Strang 5, freiliegende Knochen,
teilweise nekrotisch} & Freiliegende Sehnen, Freiliegende Knochen, Nekrose \\ \hline
WUNDGRUND\_GT\_MAPPING & \texttt{großes Ulcus: belegt, teilweise nekrotisch
Zehe: trockene Nekrose- bitte nicht durch autolytisches Débridement anfeuchten. Zeh fällt irgendwann ab, da die Durchblutung massiv gestört ist. Gefäßstatus prüfen und gfls. sanieren} & Belag, Nekrose \\ \hline
WUNDGRUND\_GT\_MAPPING & \texttt{Kleine Nekrose oberhalb des großen offenen Ulcus, muss mit versorgt werden
sonstig massiv belegt, Nekrosenbildung bei Nichtbehandlung wahrscheinlich. Prüfung ob knöcherne Strukturen / Sehnen betroffen sind} & Nekrose, Belag \\ \hline
WUNDGRUND\_GT\_MAPPING & \texttt{belegt, teilweise nekrotisch, teilweise überschießende Granulation} & Belag, Nekrose, Hypergranulation \\ \hline
WUNDGRUND\_GT\_MAPPING & \texttt{stark belegt, neben Fibrin könnte es sich auch um feuchte, wandständige Nekrosen handeln} & Fibrinbelag \\ \hline
WUNDGRUND\_GT\_MAPPING & \texttt{Granulation und Beläge gleichzeitig} & Granulation, Belag \\ \hline
WUNDGRUND\_GT\_MAPPING & \texttt{belegt, mäßige Exsudation Infektion möglich} & Belag, Exsudat \\ \hline
WUNDGRUND\_GT\_MAPPING & \texttt{leichte Granulation, belegt mit Fibrin, 2 Inseln Beurteilung hier schwer, da nicht genau zu erkennen sind, ob es ich um Nekrosen handelt.} & Granulation, Fibrinbelag \\ \hline
WUNDGRUND\_GT\_MAPPING & \texttt{belegt, hypergranulation} & Belag, Hypergranulation \\ \hline
WUNDGRUND\_GT\_MAPPING & \texttt{nekrotisch und fibrinös belegt} & Nekrose, Fibrinbelag \\ \hline
WUNDGRUND\_GT\_MAPPING & \texttt{belegt, teilweise kleine Nekrosen, könnte infiziert sein} & Belag, Nekrose \\ \hline
WUNDGRUND\_GT\_MAPPING & \texttt{kleinere Nekrosen, belegt} & Nekrose, Belag \\ \hline
WUNDGRUND\_GT\_MAPPING & \texttt{belegt, teilweise nekrotisch könnte infiziert sein} & Belag, Nekrose \\ \hline
WUNDGRUND\_GT\_MAPPING & \texttt{grünlich belegt, deutet auf Infektion hin, Nekrosen am  Wundgrund belegt} & Belag, Infektion, Nekrose \\ \hline
WUNDGRUND\_GT\_MAPPING & \texttt{stark belegte, super infizierte Wunde, teilweise freiliegende Strukturen,
Nekrosen.} & Belag, Infektion, Freiliegende Strukturen, Nekrose \\ \hline
WUNDGRUND\_GT\_MAPPING & \texttt{stark geschädigt, nekrotisch, fibrinös belegt, freiliegende Sehnen / Knochen} & Nekrose, Fibrinbelag, Freiliegende Sehnen, Freiliegende Knochen \\ \hline
WUNDGRUND\_GT\_MAPPING & \texttt{Granulationsgewebe, teilweise fibrinös belegt, Biofilm} & Granulation, Fibrinbelag, Biofilm \\ \hline
WUNDGRUND\_GT\_MAPPING & \texttt{belegt, fibrinös, nekrotisch} & Fibrinbelag, Nekrose \\ \hline
WUNDGRUND\_GT\_MAPPING & \texttt{Granulationsinseln, fibrinös belegt} & Granulation, Fibrinbelag \\ \hline
WUNDGRUND\_GT\_MAPPING & \texttt{belegt, teilweise kleine Nekrosen, könnte infiziert sein, belegt} & Belag, Nekrose \\ \hline
WUNDGRUND\_GT\_MAPPING & \texttt{belegt, teilweise kleine Nekrosen freiliegende Sehne sichtbar} & Belag, Nekrose, Freiliegende Sehnen \\ \hline
WUNDGRUND\_GT\_MAPPING & \texttt{tiefer Dekubitus Wundgrund belegt, Taschenbildung teilweise nekrotisch} & Belag, Taschenbildung, Nekrose \\ \hline
WUNDGRUND\_GT\_MAPPING & \texttt{feuchte Nekrose,} & Nekrose \\ \hline
WUNDGRUND\_GT\_MAPPING & \texttt{nekrotisch, wandständige Nekrosen} & Nekrose \\ \hline
WUNDGRUND\_GT\_MAPPING & \texttt{nekrotisch, feuchte nekrose} & Nekrose \\ \hline
WUNDGRUND\_GT\_MAPPING & \texttt{fleischige Wunde mit Taschenbildung nach Feststellung der Ausdehnung} & Granulation, Taschenbildung \\ \hline
WUNDGRUND\_GT\_MAPPING & \texttt{gerötet, stark geschädigt, teilweise nekrotisch belegt} & Rötung, Nekrose, Belag \\ \hline
WUNDGRUND\_GT\_MAPPING & \texttt{Ausgedehnte gelblich-fibrinöse Beläge mit schwarz-braunen Nekrosen, dazwischen rötliche granulierende Areale; feucht-glänzend.} & Fibrinbelag, Nekrose, Rötung, Granulation \\ \hline
WUNDGRUND\_GT\_MAPPING & \texttt{Ausgedehnte schwarz-braune Nekrose/Eschar mit feucht-glänzender Oberfläche, teils fibrinös belegt; keine Granulation sichtbar.} & Nekrose, Fibrinbelag \\ \hline
WUNDGRUND\_GT\_MAPPING & \texttt{Feucht glänzender, teils gelblich fibrinös belegter Wundgrund mit fragiler Granulation, oberflächlich, keine erkennbaren Taschen.} & Fibrinbelag, Granulation \\ \hline
WUNDGRUND\_GT\_MAPPING & \texttt{Feucht-glänzender Wundgrund mit flächigen, gelblich-weißen fibrinösen Belägen/Slough und vereinzelten rötlichen Granulationsinseln; keine schwarze Nekrose erkennbar; keine freiliegenden Sehnen/Knochen.} & Fibrinbelag, Granulation \\ \hline
WUNDGRUND\_GT\_MAPPING & \texttt{Feuchtes, gut durchblutetes granulierendes Wundbett mit kleinen, fest anhaftenden fibrinösen Inseln/Belägen; kein nekrotisches Gewebe sichtbar.} & Granulation, Fibrinbelag \\ \hline
WUNDGRUND\_GT\_MAPPING & \texttt{Feuchtes, teils granulierendes Wundbett mit ausgeprägten gelb-weißlichen Fibrinbelägen; keine trockene Nekrose erkennbar.} & Granulation, Fibrinbelag \\ \hline
WUNDGRUND\_GT\_MAPPING & \texttt{Fibrinös-gelblicher Belag mit Inseln rötlicher Granulation; feucht, keine trockene Nekrose sichtbar.} & Fibrinbelag, Rötung, Granulation \\ \hline
WUNDGRUND\_GT\_MAPPING & \texttt{Flaches Ulkus mit gelblich-fibrinösen Belägen und Anteilen von Granulationsgewebe; feucht.} & Ulzeration, Fibrinbelag, Granulation \\ \hline
WUNDGRUND\_GT\_MAPPING & \texttt{Flaches, irreguläres Ulkus mit gelblich-fibrinösen Belägen und Anteilen von Granulationsgewebe; feucht.} & Ulzeration, Fibrinbelag, Granulation \\ \hline
WUNDGRUND\_GT\_MAPPING & \texttt{Flaches, ovales Ulkus mit überwiegend gelblich-beigem Fibrinbelag; feucht glänzend; keine schwarzen Nekrosen und keine freiliegenden tiefen Strukturen sichtbar.} & Ulzeration, Fibrinbelag \\ \hline
WUNDGRUND\_GT\_MAPPING & \texttt{Flächendeckende, zähe gelblich-weiße Fibrinbeläge/Slough (\textasciitilde{}80\%) mit geringen granulierenden Inseln; feucht; keine schwarze Nekrose; unregelmäßige Wundtiefe.} & Fibrinbelag, Granulation \\ \hline
WUNDGRUND\_GT\_MAPPING & \texttt{Flächig erythematös-erosiv, teils mit weißlich-schuppigen Auflagerungen; kein nekrotisches Gewebe erkennbar; feucht-glänzend; mäßige Exsudation.} & Rötung, Exsudat \\ \hline
WUNDGRUND\_GT\_MAPPING & \texttt{Gelb-beiger, feuchter fibrinöser Belag/flächiger Slough, kaum Granulation sichtbar.} & Fibrinbelag \\ \hline
WUNDGRUND\_GT\_MAPPING & \texttt{Gelb-ockerfarbener, fest anhaftender Fibrin-/Nekrosebelag mit zentral dunklen (schwarz-braunen) Nekroseinseln; keine sichtbare Granulation; Oberfläche teils krustös, teils feucht glänzend.} & Fibrinbelag, Nekrose \\ \hline
WUNDGRUND\_GT\_MAPPING & \texttt{Gelb‑bräunlicher Fibrinbelag mit zentralen schwarz‑nekrotischen Arealen; kaum/keine sichtbare Granulation.} & Fibrinbelag, Nekrose \\ \hline
WUNDGRUND\_GT\_MAPPING & \texttt{Gelblich-beiger, feuchter Fibrinbelag/Slough (\textasciitilde{}70\%) mit randständiger rötlicher Granulation (\textasciitilde{}30\%); keine schwarze Nekrose sichtbar.} & Fibrinbelag, Rötung, Granulation \\ \hline
WUNDGRUND\_GT\_MAPPING & \texttt{Gelblich-fibrinös belegter, feucht-glänzender Wundgrund mit teils rötlichen Arealen; kein trockener Nekroseanteil sichtbar.} & Fibrinbelag, Rötung \\ \hline
WUNDGRUND\_GT\_MAPPING & \texttt{Gelblich-weißer, zäher fibrinöser Belag mit wenigen punktförmigen Blutungen; feucht, vereinzelte rötliche Areale.} & Fibrinbelag, Rötung \\ \hline
WUNDGRUND\_GT\_MAPPING & \texttt{Gemischt: überwiegend rote Granulation mit teils gelblich-fibrinösen Belägen; feuchter Wundgrund.} & Rötung, Granulation, Fibrinbelag \\ \hline
WUNDGRUND\_GT\_MAPPING & \texttt{Gemischter Wundgrund mit gelblichen fibrinösen Belägen/Slough und rötlichem Granulationsgewebe; feucht, mäßige Tiefe; kein freiliegender Knochen sichtbar.} & Fibrinbelag, Rötung, Granulation \\ \hline
WUNDGRUND\_GT\_MAPPING & \texttt{Gemischtes Wundbett mit schwarzen nekrotischen Arealen und gelb-grauem Fibrinbelag; zentral rötliches, feuchtes Granulationsgewebe; tiefe Ulzeration mit möglicher Unterminierung.} & Nekrose, Fibrinbelag, Rötung, Granulation, Ulzeration, Unterminierung \\ \hline
WUNDGRUND\_GT\_MAPPING & \texttt{Granulierender Wundgrund mit ausgeprägten gelblichen fibrinösen Belägen/Slough, teils adhärent; vereinzelt dunklere nekrotische Areale distal.} & Granulation, Fibrinbelag, Nekrose \\ \hline
WUNDGRUND\_GT\_MAPPING & \texttt{Großflächig gelb‑grünlicher, dicker Fibrin-/Slough-Belag, feucht und glänzend; am unteren Rand schwarz‑braune nekrotische/krustige Areale; nur wenige rötliche Granulationsinseln sichtbar; keine freiliegenden Strukturen (Knochen/Sehnen) erkennbar.} & Fibrinbelag, Nekrose, Granulation \\ \hline
WUNDGRUND\_GT\_MAPPING & \texttt{Großflächig nekrotisch mit schwarz-braunen Escharen; darüber ausgedehnter gelb-grünlicher fibrinöser Belag/Slough, feucht-glänzend; keine sichtbare Granulation oder Epithelisierung.} & Nekrose, Fibrinbelag \\ \hline
WUNDGRUND\_GT\_MAPPING & \texttt{Großflächiger gelb‑grünlicher Fibrinbelag/Slough mit einzelnen nekrotischen Arealen distal; feucht, kaum sichtbare Granulation.} & Fibrinbelag, Nekrose \\ \hline
WUNDGRUND\_GT\_MAPPING & \texttt{Großflächiger, feucht glänzender Wundgrund mit hell- bis dunkelroter Granulation; multiple gelblich‑weiße Fibrinbeläge/Slough-Inseln; unregelmäßige Oberfläche; keine trockene Nekrose sichtbar.} & Granulation, Fibrinbelag \\ \hline
WUNDGRUND\_GT\_MAPPING & \texttt{Großflächiges, feuchtes Wundbett mit dicken gelblich-grünlichen Fibrinbelägen/Slough; stellenweise rötliche Granulationsinseln; keine schwarze Trockennekrose erkennbar.} & Fibrinbelag, Granulation \\ \hline
WUNDGRUND\_GT\_MAPPING & \texttt{Großes rundes Ulkus mit überwiegend gelblich-weißem Fibrin-/Schleimhautbelag, zentral teils rötliche Granulationsinseln; randständig trockene nekrotische Krusten; feucht-glänzend; kein freiliegender Knochen oder Sehne sichtbar.} & Ulzeration, Fibrinbelag, Granulation, Nekrose \\ \hline
WUNDGRUND\_GT\_MAPPING & \texttt{Hauptläsion: flaches, ovales Ulkus mit gelblich-beigem fibrinösem Belag, randständig teils rötlich, feucht glänzend; keine schwarze Nekrose, feuchtem Wundbett. Zweite kleinere, oberflächliche Läsion lateral mit rötlichem, keine sichtbaren tieferen Strukturen.} & Ulzeration, Fibrinbelag, Rötung \\ \hline
WUNDGRUND\_GT\_MAPPING & \texttt{Kavitäre Ulzeration mit überwiegend gelblich-cremigem fibrinösem Belag/Slough; stellenweise rötliche Granulationsinseln an der Innenwand; feucht mit dickflüssigem Exsudat; keine freiliegenden Knochen oder Sehnen sichtbar.} & Ulzeration, Fibrinbelag, Granulation, Exsudat \\ \hline
WUNDGRUND\_GT\_MAPPING & \texttt{Kleines rundes Ulkus mit gelblich-fibrinösem Belag; kein freiliegendes tieferes Gewebe sichtbar.} & Ulzeration, Fibrinbelag \\ \hline
WUNDGRUND\_GT\_MAPPING & \texttt{Kleines, ovales Ulkus mit gelblich-fibrinösem Belag; feucht, keine freiliegenden Strukturen sichtbar.} & Ulzeration, Fibrinbelag \\ \hline
WUNDGRUND\_GT\_MAPPING & \texttt{Kleine punktförmige Öffnungen/Pusteln mit serösem bis seropurulentem Exsudat, umgebend entzündlich gerötet; kein sichtbarer Nekrosenbelag.} & Rötung, Exsudat \\ \hline
WUNDGRUND\_GT\_MAPPING & \texttt{Kleine rundliche Ulzeration mit gelblich-fibrinösem Belag, flach bis mäßig tief; kein sichtbares nekrotisches Gewebe.} & Ulzeration, Fibrinbelag \\ \hline
WUNDGRUND\_GT\_MAPPING & \texttt{Kraterförmige Ulzeration mit überwiegend gelblich-weißem fibrinösem Belag (Slough) und Anteilen rötlichen, feuchten Granulationsgewebes; kein schwarz nekrotisches Gewebe erkennbar.} & Ulzeration, Fibrinbelag, Granulation \\ \hline
WUNDGRUND\_GT\_MAPPING & \texttt{Kraterförmiges Ulkus mit gelblichen fibrinösen Belägen, teilweise rötlich-granulierende Areale; mäßig exsudierend.} & Ulzeration, Fibrinbelag, Rötung, Granulation \\ \hline
WUNDGRUND\_GT\_MAPPING & \texttt{Kräftig rot, homogen granulierend, feucht-glänzend ohne sichtbare Nekrosen oder dicke Fibrinbeläge.} & Rötung, Granulation \\ \hline
WUNDGRUND\_GT\_MAPPING & \texttt{Livide bis schwärzliche Areale mit partiell intakten/rupturierten Blasen; oberflächliche Hautnekrosen, kein sichtbares Granulationsgewebe; flächige Beteiligung ohne erkennbare Taschen.} & Blasenbildung, Nekrose \\ \hline
WUNDGRUND\_GT\_MAPPING & \texttt{Mehrere flache bis mäßig vertiefte, teilweise konfluierende Ulzera mit ausgedehntem gelb-grünlichem fibrinösem Belag/Slough; stellenweise rötliche, feinkörnige Granulation sichtbar; feucht glänzend; keine schwarze Nekrose.} & Ulzeration, Fibrinbelag, Granulation \\ \hline
WUNDGRUND\_GT\_MAPPING & \texttt{Mehrere flache, rundliche Ulzera mit gelblichen Fibrinbelägen, teils rötliche Granulationsinseln; feucht glänzend.} & Ulzeration, Fibrinbelag, Rötung, Granulation \\ \hline
WUNDGRUND\_GT\_MAPPING & \texttt{Mehrere oberflächliche Läsionen (Erosionen/Ulzera) mit gelb-beigen Fibrinbelägen und rötlicher Granulation; feucht-glänzend; keine schwarze Nekrose sichtbar.} & Ulzeration, Fibrinbelag, Granulation \\ \hline
WUNDGRUND\_GT\_MAPPING & \texttt{Mehrere oberflächliche Ulzerationen mit gelblich-grünlichem fibrinösem Belag, feucht; teils rote Granulationsinseln.} & Ulzeration, Fibrinbelag, Granulation, Rötung \\ \hline
WUNDGRUND\_GT\_MAPPING & \texttt{Mehrere sakrale Ulzera, teils tiefer liegend; gemischter Wundgrund mit rötlichem Granulationsgewebe und gelblichen Fibrinbelägen; feucht.} & Ulzeration, Granulation, Fibrinbelag, Rötung, Feucht \\ \hline
WUNDGRUND\_GT\_MAPPING & \texttt{Mehrere tiefe Ulzera mit gelblich-weißem Fibrinbelag und schwarzer Nekrose, feucht, teils mit Taschenbildung.} & Ulzeration, Fibrinbelag, Nekrose, Taschenbildung \\ \hline
WUNDGRUND\_GT\_MAPPING & \texttt{Mehrere ulzerierende Läsionen, teils mit schwarzer Nekrose/Eschar, teils gelb-fibrinös belegt; randständig teils rötlich, vereinzelt feuchte Areale mit beginnender Granulation.} & Ulzeration, Nekrose, Fibrinbelag, Rötung, Granulation \\ \hline
WUNDGRUND\_GT\_MAPPING & \texttt{Mehrere zirkuläre bis ovale, flache bis mäßig tiefe Ulzera mit gelblich-beigen Fibrinbelägen/Slough; zentral/peripher rötliche Granulationsinseln; feucht glänzend; keine trockene Schwarznekrose sichtbar.} & Ulzeration, Fibrinbelag, Granulation \\ \hline
WUNDGRUND\_GT\_MAPPING & \texttt{Oberflächlich bis mäßig tiefes Ulkus mit gelblich-fibrinösen Belägen und Anteilen rötlicher Granulation.} & Ulzeration, Fibrinbelag, Rötung, Granulation \\ \hline
WUNDGRUND\_GT\_MAPPING & \texttt{Oberflächliche Ulzeration mit gelb-grünlichem Fibrin-/Slough-Belag, stellenweise rötliche Granulation sichtbar; feucht, unregelmäßige Kontur, glänzende Oberfläche.} & Ulzeration, Fibrinbelag, Granulation \\ \hline
WUNDGRUND\_GT\_MAPPING & \texttt{Oberflächliche Ulzerationen mit gelblichen Fibrinbelägen und Anteilen von feuchtem rotem Granulationsgewebe; keine schwarze Nekrose sichtbar.} & Ulzeration, Fibrinbelag, Rötung, Granulation \\ \hline
WUNDGRUND\_GT\_MAPPING & \texttt{Oberflächliche, feuchte, rötliche Dermis mit erosiven Arealen; multiple intakte und rupturierte seröse Blasen; kein nekrotisches Gewebe sichtbar.} & Rötung, Blasenbildung \\ \hline
WUNDGRUND\_GT\_MAPPING & \texttt{Oberflächliche, teils nässende Erosionen mit gerötetem Wundgrund; keine tiefe Nekrose sichtbar.} & Rötung \\ \hline
WUNDGRUND\_GT\_MAPPING & \texttt{Oberflächlicher, größerer Defekt mit gelblich-fibrinösen Belägen und teils rötlicher Granulation; feucht glänzend, kein schwarzes Nekrosegewebe sichtbar.} & Fibrinbelag, Rötung, Granulation \\ \hline
WUNDGRUND\_GT\_MAPPING & \texttt{Oberflächliches Ulkus mit ausgedehnten gelblichen Fibrinbelägen, teils rötlich granulierend; feucht glänzend.} & Ulzeration, Fibrinbelag, Rötung, Granulation \\ \hline
WUNDGRUND\_GT\_MAPPING & \texttt{Oberflächliches Ulkus mit gelb-grünlichem Fibrin-/Belag, teils nekrotisch; feucht, unregelmäßiger Wundgrund.} & Ulzeration, Fibrinbelag, Nekrose \\ \hline
WUNDGRUND\_GT\_MAPPING & \texttt{Oberflächliches Ulkus mit gelblich-fibrinösen Belägen, randständig teils granulierend; kleine zweite Läsion lateral.} & Ulzeration, Fibrinbelag, Granulation \\ \hline
WUNDGRUND\_GT\_MAPPING & \texttt{Ovaler Defekt mit zentralem gelblich-beigem Fibrin/Slough, peripher rötliches granulierendes Gewebe; feucht-glänzende Oberfläche; keine sichtbare trockene Nekrose oder freiliegende tiefe Strukturen.} & Ulzeration, Fibrinbelag, Granulation \\ \hline
WUNDGRUND\_GT\_MAPPING & \texttt{Ovale, mäßig tiefe Ulzeration mit überwiegend rötlichem, feuchtem Granulationsgewebe; randständig teils adhärente gelblich-fibrinöse Beläge; keine schwarze Nekrose sichtbar.} & Ulzeration, Granulation, Fibrinbelag \\ \hline
WUNDGRUND\_GT\_MAPPING & \texttt{Ovale, flache Ulzeration mit überwiegend gelblich-fibrinösem Belag (ca. 70-80\%) und dazwischen rötlicher, feucht glänzender Granulation (ca. 20-30\%); keine schwarze Nekrose, keine sichtbare Sehnen- oder Knochenexposition.} & Ulzeration, Fibrinbelag, Granulation \\ \hline
WUNDGRUND\_GT\_MAPPING & \texttt{Punktförmige, schlitzartige Fistelöffnung(en); Wundgrund nicht einsehbar. Eine Öffnung mit klar-gelblich seröser Flüssigkeit; kein sichtbarer Fibrin- oder Nekrosebelag.} & Ulzeration, Exsudat \\ \hline
WUNDGRUND\_GT\_MAPPING & \texttt{Runde Ulzeration mit zentral gelblich-fibrinösem Belag, mäßig feucht; peripher schmaler erythematöser Saum, kein freiliegender Knochen/Sehne sichtbar.} & Ulzeration, Fibrinbelag, Rötung \\ \hline
WUNDGRUND\_GT\_MAPPING & \texttt{Runder, kraterförmiger Defekt mit überwiegend gelblich-fibrinösem Belag; feucht glänzend; keine schwarze Nekrose sichtbar; kaum bis keine erkennbare Granulation.} & Ulzeration, Fibrinbelag \\ \hline
WUNDGRUND\_GT\_MAPPING & \texttt{Rundliche Ulzeration mit überwiegend gelblich-fibrinösem Belag; punktuell rote, feuchte Areale/Granulationsinseln; glänzend-feucht, keine schwarze Nekrose sichtbar.} & Ulzeration, Fibrinbelag, Granulation \\ \hline
WUNDGRUND\_GT\_MAPPING & \texttt{Stark nässend; dicke gelb-grünliche Fibrin-/Beläge mit mutmaßlichem Biofilm, teils nekrotische Areale; kaum Granulation sichtbar.} & Fibrinbelag, Biofilm, Nekrose \\ \hline
WUNDGRUND\_GT\_MAPPING & \texttt{Tief klaffende, längsovale Wunde; Wundgrund überwiegend mit dicken gelb-beigen fibrinösen Belägen/Slough bedeckt, teils graugelb; feucht-glänzend; kaum/kein sichtbares Granulationsgewebe; unregelmäßige Taschen/Kavernen erkennbar.} & Fibrinbelag, Taschenbildung \\ \hline
WUNDGRUND\_GT\_MAPPING & \texttt{Tiefes Ulkus mit gelb- bis grünlichem Fibrin-/Belag, teils avital, Biofilm-verdächtig; Höhle mit möglicher Unterminierung; geringe Inseln von Granulation.} & Ulzeration, Fibrinbelag, Biofilm, Unterminierung, Granulation \\ \hline
WUNDGRUND\_GT\_MAPPING & \texttt{Tiefes, kavernöses Areal mit überwiegend gelb-braunem, dickem fibrinös-purulentem Belag; vereinzelte nekrotische Areale; kaum sichtbares vitales Granulationsgewebe; feucht glänzend.} & Fibrinbelag, Nekrose, Taschenbildung \\ \hline
WUNDGRUND\_GT\_MAPPING & \texttt{Tiefes, rundes Ulkus mit zentral gelb-weißem Fibrin/nekrotischen Anteilen; feuchte Beläge, mehere Läsionen sichtbar.} & Ulzeration, Fibrinbelag, Nekrose \\ \hline
WUNDGRUND\_GT\_MAPPING & \texttt{Tiefe moderat, kavitäre Ulzeration mit überwiegend gelblich-cremigem fibrinösem Belag/Slough; stellenweise rötliche Granulationsinseln an der Innenwand; feucht mit dickflüssigem Exsudat; keine freiliegenden Knochen oder Sehnen sichtbar.} & Ulzeration, Fibrinbelag, Granulation, Exsudat \\ \hline
WUNDGRUND\_GT\_MAPPING & \texttt{Tiefe moderat, Wundbett nahezu vollständig (ca. 90-100\%) mit gelb-beigem, faserigem Fibrin/Slough bedeckt; kein schwarzer Schorf sichtbar; kaum bis keine Granulation erkennbar; vereinzelte rötliche Areale am Boden; Wunde klar begrenzt, feucht-glänzendem, keine offen erkennbaren freiliegenden tieferen Strukturen.} & Fibrinbelag, Rötung \\ \hline
WUNDGRUND\_GT\_MAPPING & \texttt{Tiefer Defekt mit überwiegend vitalem, feuchtem Granulationsgewebe; randständig gelblich-fibrinöse Beläge.} & Granulation, Fibrinbelag \\ \hline
WUNDGRUND\_GT\_MAPPING & \texttt{Tieferes Ulkus mit dicken gelb-gräulichen Fibrinbelägen, feucht; kaum sichtbares Granulationsgewebe.} & Ulzeration, Fibrinbelag \\ \hline
WUNDGRUND\_GT\_MAPPING & \texttt{Überwiegend gelb-grauer, dicker fibrinöser/slough-Belag mit teils nekrotischem Aspekt; feucht, zäh; kaum bis kein vitales Granulationsgewebe sichtbar; Anzeichen von Tiefe/Unterminierung.} & Fibrinbelag, Nekrose, Taschenbildung \\ \hline
WUNDGRUND\_GT\_MAPPING & \texttt{Überwiegend gelb-grünlicher fibrinöser Belag; zentral schwarz-brauner Nekroseanteil (Eschar); feuchte, glänzende Oberfläche; keine sichtbare Granulation oder Epithelinseln.} & Fibrinbelag, Nekrose \\ \hline
WUNDGRUND\_GT\_MAPPING & \texttt{Überwiegend gelb-weißlicher Fibrin-/Slough-Belag mit einzelnen Arealen rötlicher Granulation; feucht-glänzende Oberfläche; keine schwarze Nekrose sichtbar; keine freiliegenden tieferen Strukturen erkennbar.} & Fibrinbelag, Granulation \\ \hline
WUNDGRUND\_GT\_MAPPING & \texttt{Überwiegend gelb-weißlicher Fibrin-/Slough-Belag mit einzelnen grau-schwarzen nekrotischen Arealen; randständig geringe rötliche Gewebeanteile; unregelmäßiger, teilweise kavitiert; feucht glänzend; keine exponierten Knochen/Sehnen sichtbar.} & Fibrinbelag, Nekrose, Rötung, Taschenbildung \\ \hline
WUNDGRUND\_GT\_MAPPING & \texttt{Überwiegend gelblich-beiger Fibrinbelag/Slough (\textasciitilde{}70\%) mit randständiger rötlicher, feucht-glänzender Granulation; keine schwarze Nekrose sichtbar; vereinzelte punktförmige Blutungen.} & Fibrinbelag, Rötung, Granulation \\ \hline
WUNDGRUND\_GT\_MAPPING & \texttt{Überwiegend gelblich-beiger fibrinöser Belag mit Anteilen rötlicher, feucht-glänzender Granulation; keine schwarze Nekrose sichtbar; vereinzelte punktförmige Blutungen.} & Fibrinbelag, Rötung, Granulation \\ \hline
WUNDGRUND\_GT\_MAPPING & \texttt{Überwiegend gelblich-weißer Fibrinbelag/Slough mit einzelnen rötlichen Granulationsinseln; feucht-glänzend; oberflächlich bis mitteltief; kein trockener schwarzer Nekroseschorf; keine freiliegenden Sehnen oder Knochen sichtbar.} & Fibrinbelag, Granulation \\ \hline
WUNDGRUND\_GT\_MAPPING & \texttt{Überwiegend gelblich-weißer Fibrinbelag/Biofilm, dazwischen kleine Inseln granulierenden, roten Gewebes; feucht glänzend.} & Fibrinbelag, Biofilm, Granulation, Rötung \\ \hline
WUNDGRUND\_GT\_MAPPING & \texttt{Überwiegend rötlich granulierendes Gewebe mit ausgedehnten gelblich-cremigen Fibrinbelägen/Slough; vereinzelt dunklere (bräunlich-schwarze) Beläge distal; feucht glänzend; keine freiliegenden Sehnen oder Knochen erkennbar.} & Fibrinbelag, Nekrose, Rötung, Granulation \\ \hline
WUNDGRUND\_GT\_MAPPING & \texttt{Überwiegend vitales, feuchtes Granulationsgewebe mit geringen fibrinösen Belägen.} & Granulation, Fibrinbelag \\ \hline
WUNDGRUND\_GT\_MAPPING & \texttt{Ulzeration mit gelblich-fibrinösem Belag am Rand, zentral teils vitales rötliches Granulationsgewebe; mäßig tiefe Kavität, keine trockene Nekrose sichtbar.} & Ulzeration, Fibrinbelag, Rötung, Granulation \\ \hline
WUNDGRUND\_GT\_MAPPING & \texttt{Unregelmäßig konfiguriertes, eher flaches Ulkus mit überwiegend gelblich-weißlichen fibrinösen Belägen und zähem Exsudat; dazwischen rötliche, feucht-körnige Granulation; keine schwarze Nekrose sichtbar.} & Ulzeration, Fibrinbelag, Granulation, Exsudat \\ \hline
WUNDGRUND\_GT\_MAPPING & \texttt{Weitgehend rötlich-granulierend mit teils fibrinösen Belägen; feucht glänzend, mehere oberflächliche Läsionen im Cluster.} & Rötung, Granulation, Fibrinbelag \\ \hline
WUNDGRUND\_GT\_MAPPING & \texttt{Zentral schwarz-braune Nekrose, umgebend gelblich-fibrinöser Belag mit feuchtem Exsudatfilm; kein vitales Granulationsgewebe erkennbar.} & Nekrose, Fibrinbelag \\ \hline
WUNDGRUND\_GT\_MAPPING & \texttt{Zentral schwarz-braune Eschar mit fest anhaftenden Belägen; randständig livid-erythematös und leicht erhaben.} & Nekrose, Rötung \\ \hline
WUNDGRUND\_GT\_MAPPING & \texttt{Zentral trockene schwarz‑braune Nekrose (Eschar) mit gelblich‑bräunlichem Fibrin/Schorf; kein sichtbares Granulations- oder Epithelgewebe.} & Nekrose, Fibrinbelag \\ \hline
WUNDGRUND\_GT\_MAPPING & \texttt{Zwei Läsionen: kaudal ausgedehnte schwarze Nekrose (Eschar), kranial gelb-grünlicher fibrinöser Belag; feuchtes Milieu, kein sichtbares Granulationsgewebe.} & Nekrose, Fibrinbelag \\ \hline
WUNDGRUND\_GT\_MAPPING & \texttt{Zwei oberflächliche Ulzerationen mit gelblichen fibrinösen Belägen, teils rötlich granulierend, feucht; keine ausgedehnte trockene Nekrose sichtbar.} & Ulzeration, Fibrinbelag, Rötung, Granulation \\ \hline
WUNDGRUND\_GT\_MAPPING & \texttt{Zwei ulzerierende Läsionen mit dickem gelb-weißlichem Fibrinbelag/Detritus und viskösem Exsudat; kein freiliegender Knochen oder Sehne sichtbar.} & Ulzeration, Fibrinbelag, Exsudat \\ \hline
WUNDGRUND\_GT\_MAPPING & \texttt{Zwei Ulzera mit gelb-grünlichen Fibrinbelägen und teils schwarzer Nekrose; feuchter Wundgrund mit randständiger Rötung.} & Ulzeration, Fibrinbelag, Nekrose, Rötung \\ \hline
WUNDGRUND\_GT\_MAPPING & \texttt{Zwei Ulzera: oberes Ulkus mit dickem gelblich-grünlichem Fibrin-/Schmierbelag, feucht glänzend; unteres, größeres Ulkus mit ausgedehnter schwarzer, teils feucht glänzender Nekrose/Eschare. Kein sichtbar vitales Granulations- oder Epithelgewebe; Beläge überwiegend adhärent.} & Ulzeration, Fibrinbelag, Nekrose \\ \hline
WUNDGRUND\_GT\_MAPPING & \texttt{Tiefere Wunde mit dicken gelblichen fibrinösen Belägen/Detritus, kaum sichtbare Granulation; möglicherweise Wundtaschen.} & Fibrinbelag, Granulation, Taschenbildung \\ \hline
WUNDGRUND\_GT\_MAPPING & \texttt{Großflächige trockene, schwarze Eschar lateral am Rückfuß/Ferse; kleinere nekrotische Läsion am lateralen Zehenbereich. Umgebung teils stark gerötet, glänzend, keine sichtbare Granulation.} & Nekrose, Rötung \\ \hline
WUNDGRUND\_GT\_MAPPING & \texttt{Weitgehend rötlich-granulierend mit teils fibrinösen Belägen; feucht glänzend, mehrere oberflächliche Läsionen im Cluster.} & Rötung, Granulation, Fibrinbelag \\ \hline
WUNDGRUND\_GT\_MAPPING & \texttt{Großflächige, trockene schwarze Eschar (nekrotisches Gewebe) am Fersenpolster, keine sichtbare Granulation oder Sickerblutung.} & Nekrose \\ \hline
WUNDGRUND\_GT\_MAPPING & \texttt{Überwiegend gelblich-weißlicher fibrinöser Belag mit Anteilen rötlicher Granulation; kein freiliegender Knochen sichtbar.} & Fibrinbelag, Rötung, Granulation \\ \hline
WUNDGRUND\_GT\_MAPPING & \texttt{Tiefes, rundes Ulkus mit zentral gelb-weißem Fibrin/nekrotischen Anteilen; feuchte Beläge, mehrere Läsionen sichtbar.} & Ulzeration, Fibrinbelag, Nekrose \\ \hline
WUNDGRUND\_GT\_MAPPING & \texttt{Schwarze Nekrosenanteile mit dicken gelb-grünlichen fibrinösen Belägen, feucht; peripher teils gerötetes/angegriffenes Gewebe.} & Nekrose, Fibrinbelag, Rötung \\ \hline
WUNDGRUND\_GT\_MAPPING & \texttt{Zentrale schwarz-gelbe Eschar mit fest anhaftenden Belägen; randständig livid-erythematös und leicht erhaben.} & Nekrose \\ \hline
WUNDGRUND\_GT\_MAPPING & \texttt{Flache, ovale Ulzeration mit zentral gelblichen Fibrinbelägen; randständig rötliche, feuchte Granulation.} & Ulzeration, Fibrinbelag, Granulation, Rötung \\ \hline
WUNDGRUND\_GT\_MAPPING & \texttt{Kleiner runder Ulkus (\textasciitilde{}1.5–2 cm) mit gelblich-weißen fibrinösen Belägen, punktuell rötlich/blasende Areale; oberflächlich und feucht.} & Ulzeration, Fibrinbelag, Rötung, Blasenbildung \\ \hline
WUNDGRUND\_GT\_MAPPING & \texttt{Ovales Ulkus mit dicken gelb-weißlichen Fibrinbelägen, zentral teils grau-schwarze nekrotische Areale; kaum sichtbare Granulation.} & Ulzeration, Fibrinbelag, Nekrose \\ \hline
WUNDGRUND\_GT\_MAPPING & \texttt{starke Granulation, überschießendes Gewebe, infektfrei} & Granulation, Hypergranulation \\ \hline
WUNDGRUND\_GT\_MAPPING & \texttt{Hämangiome sind Neoplasien ( Neubildungen) , umschriebenes Hämangiom, klare Abgrenzung. segmentale Fehlbildung.
Behandlung unbedingt notwendig.} & Hämangiom \\ \hline
WUNDGRUND\_GT\_MAPPING & \texttt{Exsudation} & Exsudat \\ \hline
WUNDGRUND\_GT\_MAPPING & \texttt{Tiefer Defekt mit dicken gelblichen Fibrin-/Schleimbelägen, teils nekrotisch; kaum Granulation sichtbar.} & Fibrinbelag, Nekrose \\ \hline
WUNDGRUND\_GT\_MAPPING & \texttt{Vitales, homogen granulierendes Gewebe, feucht glänzend, keine sichtbaren Beläge oder Nekrosen.} & Granulation \\ \hline
WUNDGRUND\_GT\_MAPPING & \texttt{Großflächige schwarze/braune Nekrosen und gelblich-fibrinöse Beläge, teils feucht; randständig kleine Areale mit frischer Blutung/Granulation.} & Fibrinbelag, Granulation, Nekrose \\ \hline
WUNDGRUND\_GT\_MAPPING & \texttt{Teilweise freiliegende, feucht-rosige Dermis mit oberflächlichen Erosionen; mehrere intakte und rupturierte seröse Blasen, vereinzelte dünne Fibrinauflagen.} & Blasenbildung, Fibrinbelag \\ \hline
WUNDGRUND\_GT\_MAPPING & \texttt{Flächig oberflächlich erosiv, feucht glänzend mit teils fibrinösen Belägen und rötlich-granulierenden Arealen.} & Fibrinbelag, Granulation \\ \hline
WUNDGRUND\_GT\_MAPPING & \texttt{Kein freiliegender Wundgrund sichtbar; intakte bzw. teils abgehobene Epidermis mit Hämatom und Blasen} & Blasenbildung, Hämatom \\ \hline
WUNDGRUND\_GT\_MAPPING & \texttt{Großflächig trockene schwarze Nekrosen (Eschar) an Ferse und lateraler Fußseite; kleinere Areale mit gelblich-fibrinösem Belag proximal; keine erkennbare Granulation.} & Fibrinbelag, Nekrose \\ \hline
WUNDGRUND\_GT\_MAPPING & \texttt{Ausgedehnte schwarz-braune Nekrosen mit gelblich-fibrinösen Belägen; zentral rötlich-granulierende Anteile; unregelmäßiger, teils tiefer Defekt.} & Fibrinbelag, Granulation, Nekrose \\ \hline
WUNDGRUND\_GT\_MAPPING & \texttt{Schwarz-braune Nekrosen und fest anhaftende Fibrinbeläge; teils blutiges, freiliegendes Gewebe am Vorfuß; mehrere oberflächliche bis mäßig tiefe Ulzera.} & Fibrinbelag, Nekrose \\ \hline
WUNDGRUND\_GT\_MAPPING & \texttt{Granulierendes Gewebe mit gelblich-fibrinösen Belägen; schräger, mäßig tiefer Defekt mit möglicher Randunterminierung; feucht glänzend.} & Fibrinbelag, Granulation, Taschenbildung \\ \hline
WUNDGRUND\_GT\_MAPPING & \texttt{Gemischt: gelblich-fibrinöse Beläge mit Anteilen von feuchter Granulation; unregelmäßige, flach bis mäßig tiefe Läsion.} & Fibrinbelag, Granulation \\ \hline
WUNDGRUND\_GT\_MAPPING & \texttt{Flacher Defekt mit überwiegender roter Granulation und teils gelblich-fibrinösen Belägen; keine trockene Nekrose sichtbar.} & Fibrinbelag, Granulation \\ \hline
WUNDGRUND\_GT\_MAPPING & \texttt{Mehrere tiefe Ulzerationen mit gelblich-fibrinösen Belägen und schwarzer Nekrose, teilweise rötlich-granulierende Areale, feucht.} & Fibrinbelag, Granulation, Nekrose \\ \hline
WUNDGRUND\_GT\_MAPPING & \texttt{vitalrotes, glänzend-feuchtes Granulationsgewebe, teils bereits beginnende Epithelisierung, keine Nekrosen sichtbar} & Granulation, Wundrandepithelisierung \\ \hline
WUNDGRUND\_GT\_MAPPING & \texttt{Mehrere flache Ulzera mit gelblich-fibrinösen Belägen, teils granuliert; einzelne Blutkrusten, kein freiliegender Knochen/Sehne sichtbar.} & Fibrinbelag, Granulation \\ \hline
WUNDGRUND\_GT\_MAPPING & \texttt{gelblich-brauner Fibrinbelag mit teils rötlichen Anteilen; keine freiliegenden Strukturen sichtbar} & Fibrinbelag \\ \hline
WUNDGRUND\_GT\_MAPPING & \texttt{Gemischter Wundgrund mit gelb-grünlichen Fibrinbelägen und teils schwarzer trockener Nekrose; teilweise gerötete, leicht feuchte Umgebung; geringe Blutungsspuren} & Fibrinbelag, Nekrose \\ \hline
WUNDGRUND\_GT\_MAPPING & \texttt{Großflächige schwarz-braune, teils feucht glänzende Nekrose/Schorf ohne sichtbares vitales Gewebe; Übergangszonen mit grauen Belägen.} & Fibrinbelag, Nekrose \\ \hline
WUNDGRUND\_GT\_MAPPING & \texttt{Großflächige, schwarz-braune, trockene Nekrose/Eschar; keine sichtbare Granulation, minimaler erythematöser Saum} & Nekrose \\ \hline
WUNDGRUND\_GT\_MAPPING & \texttt{Punktförmige Läsion mit gelb-grünlichem Belag/Eiterpfropf, Tiefe unklar.} & Fibrinbelag \\ \hline
WUNDGRUND\_GT\_MAPPING & \texttt{Zentral gelblicher fibrinöser Belag, randständig teils rötliche Granulation; kein freiliegender Knochen sichtbar.} & Fibrinbelag, Granulation \\ \hline
WUNDGRUND\_GT\_MAPPING & \texttt{Großflächiger gelblich-fibrinöser Belag, teils serös‑blutig unterlaufen; oberflächlicher bis teildermaler Substanzverlust.} & Fibrinbelag \\ \hline
WUNDGRUND\_GT\_MAPPING & \texttt{gelblich-fibrinös belegter, teils nekrotischer Wundgrund mit Ulkuskrater; zweite Läsion mit trockener Nekrose/Schorf} & Fibrinbelag, Nekrose \\ \hline
WUNDGRUND\_GT\_MAPPING & \texttt{Ausgedehnter schwarz-brauner Nekroseanteil mit gelblich-grünem Fibrin/Slough, feucht glänzend; möglicher Biofilm} & Biofilm, Fibrinbelag, Nekrose \\ \hline
WUNDGRUND\_GT\_MAPPING & \texttt{Gelb-weißliche fibrinöse Beläge mit rötlich granulierenden Arealen; feuchter Wundgrund.} & Fibrinbelag, Granulation \\ \hline
WUNDGRUND\_GT\_MAPPING & \texttt{Zentrale schwarz-braune Eschar mit gelblich-fibrinösen Arealen; verhärteter, livid-erythematöser Rand.} & Fibrinbelag, Nekrose, Rötung \\ \hline
WUNDGRUND\_GT\_MAPPING & \texttt{Irregulärer Ulkus mit gelblichen fibrinösen Belägen und Anteilen rötlicher Granulation; feucht glänzend.} & Fibrinbelag, Granulation \\ \hline
WUNDGRUND\_GT\_MAPPING & \texttt{Gelblicher, zäher Fibrinbelag; kaum Granulation; teils tieferliegendes Gewebe/Faszie sichtbar.} & Fibrinbelag \\ \hline
WUNDGRUND\_GT\_MAPPING & \texttt{Großflächig, feucht; überwiegend rote Granulation mit gelblichen fibrinösen Belägen/Slough.} & Fibrinbelag, Granulation \\ \hline
WUNDGRUND\_GT\_MAPPING & \texttt{Zentral gelblich-fibrinöser Belag, flach; geringe/keine sichtbare Granulation; kein nekrotisches Gewebe erkennbar.} & Fibrinbelag \\ \hline
WUNDGRUND\_GT\_MAPPING & \texttt{Überwiegend rot-granulierend mit wenigen gelblich-fibrinösen Inseln; feucht; flach.} & Fibrinbelag, Granulation \\ \hline
WUNDGRUND\_GT\_MAPPING & \texttt{ausgedehnter gelblich-fibrinöser Belag/Biofilm mit einzelnen rötlich-granulierenden Anteilen; feucht} & Biofilm, Fibrinbelag, Granulation \\ \hline
WUNDGRUND\_GT\_MAPPING & \texttt{oberflächliche Ulzeration mit gelblich-weißen Fibrinbelägen und vereinzelten rötlich-granulierenden Arealen; keine trockene Nekrose sichtbar} & Fibrinbelag, Granulation \\ \hline
WUNDGRUND\_GT\_MAPPING & \texttt{Gelb-bräunliche Fibrinauflagerungen mit zentralen schwarz-braunen Nekroseanteilen; kein Granulationsgewebe sichtbar.} & Fibrinbelag, Nekrose \\ \hline
WUNDGRUND\_GT\_MAPPING & \texttt{Ovales Ulkus mit gelblich-fibrinösem Belag, feucht-glänzend; punktuell gerötete Granulation sichtbar.} & Fibrinbelag, Granulation \\ \hline
WUNDGRUND\_GT\_MAPPING & \texttt{Ovales, oberflächliches Ulkus mit gelblich-fibrinösem Belag zentral, peripher rötlich granulierend.} & Fibrinbelag, Granulation \\ \hline
WUNDGRUND\_GT\_MAPPING & \texttt{Runde, oberflächliche Läsion mit gelblich-fibrinösem Belag und punktuellen Granulationen; feucht glänzend.} & Fibrinbelag, Granulation \\ \hline
WUNDGRUND\_GT\_MAPPING & \texttt{Zentral gelblich-weisslicher Fibrinbelag mit zwei rötlichen Granulationsinseln; oberflächlich, feucht.} & Fibrinbelag, Granulation \\ \hline
WUNDGRUND\_GT\_MAPPING & \texttt{Ovales Ulkus mit randständigem gelblich-fibrinösem Belag und zentralen Granulationsarealen; keine freiliegenden Strukturen sichtbar.} & Fibrinbelag, Granulation \\ \hline
WUNDGRUND\_GT\_MAPPING & \texttt{Ausgedehnte gelb-grünliche Fibrin-/Beläge mit Biofilm, teils weich-nekrotische Areale; feucht, keine freiliegenden Strukturen erkennbar.} & Biofilm, Fibrinbelag, Nekrose \\ \hline
WUNDGRUND\_GT\_MAPPING & \texttt{Ausgedehnter gelb-grünlicher Fibrin-/Detritusbelag mit einzelnen nekrotischen Arealen, kaum Granulation sichtbar; oberflächlich bis flach, unregelmäßige Ränder.} & Fibrinbelag, Nekrose \\ \hline
WUNDGRUND\_GT\_MAPPING & \texttt{Weitläufige, flache Ulzeration mit gelblich-fibrinösem Belag, feucht; vereinzelte punktförmige Blutungen; keine freiliegenden Strukturen erkennbar.} & Fibrinbelag \\ \hline
WUNDGRUND\_GT\_MAPPING & \texttt{Kleiner ovaler Defekt mit gelblich-fibrinösem Belag, kaum sichtbare Granulation, oberflächlich.} & Fibrinbelag \\ \hline
WUNDGRUND\_GT\_MAPPING & \texttt{Gelb-grünlicher Fibrinbelag mit zentraler schwarz-brauner Nekrose; feucht erscheinend.} & Fibrinbelag, Nekrose \\ \hline
WUNDGRUND\_GT\_MAPPING & \texttt{Großflächiger gelblich-bräunlicher fibrinöser Belag mit feuchtem Wundbett; teils rötliche Granulation sichtbar; oberflächlich.} & Fibrinbelag, Granulation \\ \hline
WUNDGRUND\_GT\_MAPPING & \texttt{Gelblich-grünlicher Fibrin-/Belag mit teils rötlich punktförmig granulierenden Arealen; oberflächliche, irreguläre Ulzera.} & Fibrinbelag, Granulation \\ \hline
WUNDGRUND\_GT\_MAPPING & \texttt{Gemischter Wundgrund mit gelb-weißen Fibrinbelägen und dazwischen vitalem Granulationsgewebe; feucht, oval.} & Fibrinbelag, Granulation \\ \hline
WUNDGRUND\_GT\_MAPPING & \texttt{Kleines rundes Ulkus mit gelblich-fibrinösen Belägen, eher trocken; kein deutlich granulierendes Gewebe sichtbar.} & Fibrinbelag \\ \hline
WUNDGRUND\_GT\_MAPPING & \texttt{Oberflächliches Ulkus mit gelblich-fibrinösen Belägen; randständig teils epithelisierend; zweite kleine oberflächliche Läsion lateral.} & Fibrinbelag, Wundrandepithelisierung \\ \hline
WUNDGRUND\_GT\_MAPPING & \texttt{Gelbliche fibrinöse Beläge, teils granulierend; flach-kraterförmiger Defekt.} & Fibrinbelag, Granulation \\ \hline
WUNDGRUND\_GT\_MAPPING & \texttt{Ausgedehnte schwarze trockene Nekrose am unteren Ulkus, oberes Ulkus mit gelblichen Fibrinbelägen; kein Granulationsgewebe sichtbar.} & Fibrinbelag, Nekrose \\ \hline
WUNDGRUND\_GT\_MAPPING & \texttt{Tiefer Ulkus mit gelblich-fibrinösen Belägen, teils nekrotisch, wenig Granulation; Höhle mit möglicher Randunterminierung.} & Fibrinbelag, Granulation, Nekrose, Taschenbildung \\ \hline
WUNDGRUND\_GT\_MAPPING & \texttt{Gelb-weißlicher fibrinöser Belag mit zentral dunkler nekrotischer Komponente; bislang kaum Granulation sichtbar.} & Fibrinbelag, Nekrose \\ \hline
WUNDGRUND\_GT\_MAPPING & \texttt{Gelb-grauer fibrinöser Belag mit feuchtem Slough; Tiefe/Unterminierungen nicht sicher beurteilbar.} & Fibrinbelag, Taschenbildung \\ \hline
WUNDGRUND\_GT\_MAPPING & \texttt{Kleine, serös gefüllte Blase mit möglicher punktueller Erosion; kein nekrotisches Gewebe sichtbar.} & Blasenbildung \\ \hline
WUNDGRUND\_GT\_MAPPING & \texttt{Rötlich-feuchtes, gut durchblutetes Granulationsgewebe, punktuell dünner Fibrinbelag; keine Nekrosen sichtbar.} & Fibrinbelag, Granulation \\ \hline
WUNDGRUND\_GT\_MAPPING & \texttt{Mischbild aus vitalem Granulationsgewebe mit gelblich-fibrinösen Belägen; zwei tiefer reichende Ulzera, kein freiliegender Knochen sichtbar.} & Fibrinbelag, Granulation \\ \hline
WUNDGRUND\_GT\_MAPPING & \texttt{Mehrere konfluierende, oberflächlich erodierte Areale mit feucht-glänzend rotem Gewebe; punktuell serös-sanguinöse Anteile; kein nekrotisches Gewebe erkennbar; teils weißliche, schuppige Auflagerungen an und um die Erosionen.} & nan \\ \hline
WUNDGRUND\_GT\_MAPPING & \texttt{Großflächig flaches Areal mit kräftig rotem, feucht-glänzendem, homogenem Granulationsgewebe; keine sichtbaren Nekrosen oder gelblichen Fibrinbeläge; randständige beginnende Epithelisierung.} & Granulation, Wundrandepithelisierung \\ \hline
WUNDGRUND\_GT\_MAPPING & \texttt{Zwei tiefe Ulzerationen mit dickem gelblich-weißlichem fibrinösem/purulentem Belag; randständig teils rötliche Granulation; feucht, kein schwarzes Nekrosengewebe sichtbar.} & Fibrinbelag, Granulation, Infektion \\ \hline
WUNDGRUND\_GT\_MAPPING & \texttt{Überwiegend gelblich-beiger, feucht-glänzender Fibrin-/Slough-Belag; ausgedehnte dunkelbraun-schwarze Nekrose-/Escharareale; vereinzelte rötliche Granulationsinseln, vor allem randständig; heterogene Tiefenwirkung} & Fibrinbelag, Granulation, Nekrose \\ \hline
WUNDGRUND\_GT\_MAPPING & \texttt{Mehrere oberflächliche Erosionen mit feuchtem, rosig-rotem, glänzendem Wundgrund; teils intakte seröse Blasen, teils rupturierte Blasen mit freiliegender Dermis; keine sichtbaren Nekrosen oder dicker Fibrinbelag.} & Blasenbildung \\ \hline
WUNDGRUND\_GT\_MAPPING & \texttt{Flacher, ausgedehnter Defekt mit rötlich granulierendem Gewebe und ausgedehnten gelblich-fibrinösen Belägen; stark feucht/glänzend; keine schwarze Nekrose sichtbar.} & Fibrinbelag, Granulation \\ \hline
WUNDGRUND\_GT\_MAPPING & \texttt{Ausgedehnte livid‑violette Hautverfärbung mit deutlicher Schwellung; multiple intakte und kollabierte Blasen (teils serös/serosanguinös gefüllt); oberflächliche Erosionen/Krusten in rechteckigen Arealen; keine sichtbare Granulation, kein freiliegendes tiefes Gewebe erkennbar.} & Blasenbildung \\ \hline
WUNDGRUND\_GT\_MAPPING & \texttt{Großflächige schwarz-braune, trockene, fest haftende Nekrose (Eschar) an der Ferse/lateralem Rückfuß; weitere kleinere nekrotische Läsion am lateralen Vorfuß nahe dem 5. Strahl; kein sichtbares Granulations- oder Epithelgewebe; stellenweise gelblich-fibrinöse Beläge proximal am Knöchel.} & Fibrinbelag, Nekrose \\ \hline
WUNDGRUND\_GT\_MAPPING & \texttt{Gemischter Wundgrund: großflächige schwarz-braune Nekrosen/Eschar und gelb-grauer Fibrinbelag; zentral tieferer Defekt mit rötlichem, feuchtem Granulationsgewebe; unregelmäßig, teilweise kavitiert; feucht glänzend; keine exponierten Knochen/Sehnen sichtbar.} & Fibrinbelag, Granulation, Nekrose \\ \hline
WUNDGRUND\_GT\_MAPPING & \texttt{Mehrere Ulzera: 1) rundliches Ulkus am Fußrand mit zentraler schwarz‑brauner, trockener Nekrose/Eschar; 2) irreguläres Ulkus plantar am Vorfuß/Zehenbasis mit rötlicher Granulation, gelblich-fibrinösen Belägen und teils blutigen Auflagerungen; insgesamt feucht glänzend, keine sichtbaren Sehnen oder Knochen.} & Fibrinbelag, Granulation, Nekrose \\ \hline
WUNDGRUND\_GT\_MAPPING & \texttt{Überwiegend rötliches, feuchtes Granulationsgewebe mit anteilig gelblich-fibrinösen Belägen; unregelmäßige Ulkusränder; teils kavernös wirkend mit Unterminierung; keine schwarze Nekrose sichtbar.} & Fibrinbelag, Granulation, Taschenbildung \\ \hline
WUNDGRUND\_GT\_MAPPING & \texttt{Gemischter Wundgrund mit rötlicher, feucht glänzender Granulation und ausgedehnten gelblich-fibrinösen Belägen; flach und unregelmäßig, teils zusammenfließende Ulzerationen; keine schwarze Nekrose sichtbar.} & Fibrinbelag, Granulation \\ \hline
WUNDGRUND\_GT\_MAPPING & \texttt{Großflächige Ulzeration mit überwiegend vitalem, feucht-glänzendem, hell- bis dunkelrotem Granulationsgewebe; eingestreut gelblich-weißliche fibrinöse Beläge (vor allem zentral und distal); keine schwarze Nekrose sichtbar; unregelmäßige Kontur.} & Fibrinbelag, Granulation \\ \hline
WUNDGRUND\_GT\_MAPPING & \texttt{Drei rund-ovale, teils tief reichende Defekte mit zentral schwarz-brauner Nekrose (Eschar) und gelblich-weißen fibrinösen Belägen; feucht-glänzende Oberfläche; am Rand abschnittsweise rötliches Granulationsgewebe erkennbar.} & Fibrinbelag, Granulation, Nekrose \\ \hline
WUNDGRUND\_GT\_MAPPING & \texttt{Überwiegend granulierendes, rötlich bis dunkelrotes, feuchtes Gewebe; mehrere flache, teils konfluierende Ulzera; stellenweise dünner gelblicher Fibrinbelag an den Rändern; keine schwarze Nekrose sichtbar.} & Fibrinbelag, Granulation \\ \hline
WUNDGRUND\_GT\_MAPPING & \texttt{Multiple rund-ovale, teils konfluierende Ulzera mit gelblich-beigem fibrinösem Belag; Abschnitte mit rötlicher Granulation; vereinzelt hämorrhagische Krusten/Schorfinseln; überwiegend oberflächlich bis partiell dermal, feucht-glänzend ohne sichtbare tiefe Strukturen.} & Fibrinbelag, Granulation \\ \hline
WUNDGRUND\_GT\_MAPPING & \texttt{Zwei benachbarte Ulzera: proximal überwiegend gelb-beiger fibrinöser Belag; distal rötlich-feuchtes, teils granuliertes Gewebe. Keine schwarze Nekrose sichtbar. Oberflächlich bis mäßig tief.} & Fibrinbelag, Granulation \\ \hline
WUNDGRUND\_GT\_MAPPING & \texttt{Ovales Ulkus, zentral überwiegend gelblich-bräunlicher Fibrinbelag/Slough (\textasciitilde{}60–70\%), peripher rötliche feucht wirkende Granulation (\textasciitilde{}30–40\%); keine freiliegenden Sehnen oder Knochen erkennbar; moderat tief, klar begrenzte Defektzone.} & Fibrinbelag, Granulation \\ \hline
WUNDGRUND\_GT\_MAPPING & \texttt{Zwei getrennte Ulzera. Proximal: ovaler Defekt mit überwiegend gelb-grauem Fibrin/Slough (ca. 60-70\%), restlich rötlich-feuchtes Gewebe. Distal: gemischter Wundgrund mit randständiger rötlicher Granulation und zentral schwarz-brauner trockener Nekrose/Schorf; teils serös-blutige Beläge. Tiefe bis ins subkutane Gewebe, keine freiliegenden Sehnen oder Knochen sichtbar.} & Fibrinbelag, Granulation, Nekrose \\ \hline
WUNDGRUND\_GT\_MAPPING & \texttt{Großflächige schwarz-braune, teils glänzend-feuchte Eschar/Nekrose, fest haftend; randständig gelblich-graue Beläge; kein vitales Granulations- oder Epithelgewebe sichtbar.} & Fibrinbelag, Nekrose \\ \hline
WUNDGRUND\_GT\_MAPPING & \texttt{Wundbett vollständig von trockener, schwarz-brauner, fest anhaftender Ledernekrose (Eschar) bedeckt; teils grauweißliche Anteile; keine sichtbare Granulation oder Epithelisierung.} & Nekrose \\ \hline
WUNDGRUND\_GT\_MAPPING & \texttt{Kleine runde Öffnung mit gelb-grünlichem purulentem/fibrinösem Belag (Eiterpfropf); kein sichtbares Granulations- oder Epithelgewebe.} & Fibrinbelag, Infektion \\ \hline
WUNDGRUND\_GT\_MAPPING & \texttt{Überwiegend gelblich-beiger Fibrinbelag/Slough (\textasciitilde{}70\%) mit randständiger rötlicher, feuchter Granulation (\textasciitilde{}30\%); kein schwarzes Nekrosegewebe sichtbar; Wundbett feucht-glänzend, mäßige Tiefe; keine freiliegenden Strukturen (Knochen/Sehne) erkennbar.} & Fibrinbelag, Granulation \\ \hline
WUNDGRUND\_GT\_MAPPING & \texttt{Überwiegend gelblich-beiger, dicker Fibrin-/Belag mit teils grünlicher Tönung, feucht-glänzend; randständig vereinzelte rötliche granulierende Areale und punktförmige Blutungen; keine schwarze Nekrose sichtbar.} & Fibrinbelag, Granulation \\ \hline
WUNDGRUND\_GT\_MAPPING & \texttt{Großes rundes Ulkus mit überwiegend gelblich-weißem Fibrin-/Schleimhautbelag, zentral teils rötliche Granulationsinseln; randständig trockene nekrotische Krusten; feucht-glänzend, kein freiliegender Knochen oder Sehne sichtbar.} & Fibrinbelag, Granulation, Nekrose \\ \hline
WUNDGRUND\_GT\_MAPPING & \texttt{Tiefe, kraterförmige Ulzeration mit überwiegend gelblich-weißem fibrinösem Belag (Slough) und Anteilen rötlichen, feuchten Granulationsgewebes; kein schwarz nekrotisches Gewebe erkennbar.} & Fibrinbelag, Granulation \\ \hline
WUNDGRUND\_GT\_MAPPING & \texttt{Wundbett nahezu vollständig (ca. 90-100\%) mit gelb-beigem, feucht-glänzendem, faserigem Fibrin/Slough bedeckt; kein schwarzer Schorf sichtbar; kaum bis keine Granulation erkennbar; vereinzelte rötliche Areale am Boden; Wunde klar begrenzt, Tiefe moderat, keine offen erkennbaren freiliegenden tieferen Strukturen.} & Fibrinbelag \\ \hline
WUNDGRUND\_GT\_MAPPING & \texttt{Großflächiger, oberflächlicher bis mäßig tiefer Ulkus mit überwiegend vitaler, feucht glänzender, hell- bis dunkelroter Granulation; multiple gelblich‑weiße Fibrinbeläge/Slough-Inseln; unregelmäßige Oberfläche; keine trockene Nekrose sichtbar.} & Fibrinbelag, Granulation \\ \hline
WUNDGRUND\_GT\_MAPPING & \texttt{Zentral überwiegend gelblich-weißlicher fibrinöser Belag (Slough), feucht; randständig schmaler rötlicher Saum mit beginnender Granulation; keine schwarze Nekrose sichtbar.} & Fibrinbelag, Granulation \\ \hline
WUNDGRUND\_GT\_MAPPING & \texttt{Überwiegend feucht-rote Granulation, flach; zwei zentral gelegene rundliche weißlich-gelbliche Fibrininseln/Beläge; keine schwarze Nekrose; beginnende Epithelisierung vom Rand.} & Fibrinbelag, Granulation, Wundrandepithelisierung \\ \hline
WUNDGRUND\_GT\_MAPPING & \texttt{Oberflächlicher, irregulär geformter Wundgrund mit überwiegend gelblich-grünlichem fibrinösem Belag/Slough; dazwischen rötliche, feucht-glänzende Granulation; kein trockener Nekroseanteil; flach ohne sichtbare Taschen.} & Fibrinbelag, Granulation \\ \hline
WUNDGRUND\_GT\_MAPPING & \texttt{Ovale, flache Ulzeration mit überwiegend gelblich-fibrinösem Belag (ca. 70–80\%) und dazwischen rötlicher, feucht glänzender Granulation (ca. 20–30\%); keine schwarze Nekrose, keine sichtbare Sehnen- oder Knochenexposition.} & Fibrinbelag, Granulation \\ \hline
WUNDGRUND\_GT\_MAPPING & \texttt{Überwiegend dicke gelb-grüne, schmierig-fibrinöse Beläge mit purulentem, viskösem Exsudat; dazwischen rötliche Granulationsinseln; einzelne bräunlich-dunklere Areale; Wundbett feucht und unregelmäßig, keine freiliegenden Sehnen oder Knochen sichtbar.} & Exsudat, Fibrinbelag, Granulation, Infektion \\ \hline
WUNDGRUND\_GT\_MAPPING & \texttt{Oberflächliches, ovales Ulkus mit überwiegend gelblich-weißlichem, feuchtem Fibrinbelag; stellenweise rötliche Granulation sichtbar; keine schwarze Nekrose; unregelmäßige Kontur, glänzende Oberfläche.} & Fibrinbelag, Granulation \\ \hline
WUNDGRUND\_GT\_MAPPING & \texttt{Hauptläsion: flaches, ovales Ulkus mit gelblich-beigem fibrinösem Belag, randständig teils rötlich, feucht glänzend; keine schwarze Nekrose, keine sichtbaren tieferen Strukturen. Zweite kleinere, oberflächliche Läsion lateral mit rötlichem, feuchtem Wundbett.} & Fibrinbelag, Rötung \\ \hline
WUNDGRUND\_GT\_MAPPING & \texttt{Überwiegend gelb-weißlicher fibrinöser Belag (Slough), stellenweise rötliche Granulation sichtbar; feucht, unregelmäßig-kraterförmig; keine trockene schwarze Nekrose erkennbar.} & Fibrinbelag, Granulation \\ \hline
WUNDGRUND\_GT\_MAPPING & \texttt{Tiefe, kavitäre Ulzeration mit überwiegend gelblich-cremigem fibrinösem Belag/Slough; stellenweise rötliche Granulationsinseln an der Innenwand; feucht mit dickflüssigem Exsudat; keine freiliegenden Knochen oder Sehnen sichtbar.} & Exsudat, Fibrinbelag, Granulation \\ \hline
WUNDGRUND\_GT\_MAPPING & \texttt{Überwiegend gelblich-weißlicher Fibrin-/Slough-Belag mit einzelnen grau-schwarzen nekrotischen Arealen; randständig geringe rötliche Gewebeanteile; unregelmäßiger, mäßig tiefer Defekt.} & Fibrinbelag, Nekrose \\ \hline
WUNDGRUND\_GT\_MAPPING & \texttt{Überwiegend gelb-grauer fibrinöser Slough/Belag, feucht und zäh; kaum sichtbares Granulationsgewebe; unregelmäßig-oval.} & Fibrinbelag \\ \hline
WUNDGRUND\_GT\_MAPPING & \texttt{Rötliches, feuchtes Granulationsgewebe; teils gelblich fibrinöser Belag randständig; kraterförmiger, tiefer Defekt ohne schwarze Nekrosen; glänzende Oberfläche.} & Fibrinbelag, Granulation \\ \hline
WUNDGRUND\_GT\_MAPPING & \texttt{Zwei größere kraterförmige Ulzera mit überwiegend vitaler roter Granulation und teils gelblich-beigen fibrinösen Belägen; feuchtes, glänzendes Wundbett; keine trockene Schwarznekrose sichtbar.} & Fibrinbelag, Granulation \\ \hline
WUNDRAND\_GT\_MAPPING & \texttt{["Epibolie (eingerollter Wundrand)","Gerötet / entzündlich","Mazeriert"]} & Epibolie (eingerollter Wundrand), Gerötet / entzündlich, Mazeriert \\ \hline
WUNDRAND\_GT\_MAPPING & \texttt{["Epibolie (eingerollter Wundrand)","Gerötet / entzündlich"]} & Epibolie (eingerollter Wundrand), Gerötet / entzündlich \\ \hline
WUNDRAND\_GT\_MAPPING & \texttt{["Epibolie (eingerollter Wundrand)","Mazeriert"]} & Epibolie (eingerollter Wundrand), Mazeriert \\ \hline
WUNDRAND\_GT\_MAPPING & \texttt{["Epibolie (eingerollter Wundrand)"]} & Epibolie (eingerollter Wundrand) \\ \hline
WUNDRAND\_GT\_MAPPING & \texttt{["Gelbliche fibrinöse Randbeläge"]} & nan \\ \hline
WUNDRAND\_GT\_MAPPING & \texttt{["Gerötet / entzündlich","Epibolie (eingerollter Wundrand)"]} & Epibolie (eingerollter Wundrand), Gerötet / entzündlich \\ \hline
WUNDRAND\_GT\_MAPPING & \texttt{["Gerötet / entzündlich","Mazeriert","Unterminiert (Wundtaschen)"]} & Gerötet / entzündlich, Mazeriert, Unterminiert (Wundtaschen) \\ \hline
WUNDRAND\_GT\_MAPPING & \texttt{["Gerötet / entzündlich","Mazeriert"]} & Gerötet / entzündlich, Mazeriert \\ \hline
WUNDRAND\_GT\_MAPPING & \texttt{["Gerötet / entzündlich","Taschenbildung möglich"]} & Gerötet / entzündlich, Unterminiert (Wundtaschen) \\ \hline
WUNDRAND\_GT\_MAPPING & \texttt{["Gerötet / entzündlich","Unterminiert (Wundtaschen)"]} & Gerötet / entzündlich, Unterminiert (Wundtaschen) \\ \hline
WUNDRAND\_GT\_MAPPING & \texttt{["Gerötet / entzündlich","glänzende, ödematöse Haut"]} & Gerötet / entzündlich \\ \hline
WUNDRAND\_GT\_MAPPING & \texttt{["Gerötet / entzündlich","livid-erythematöse Hautveränderungen"]} & Gerötet / entzündlich \\ \hline
WUNDRAND\_GT\_MAPPING & \texttt{["Gerötet / entzündlich","relativ scharf begrenzte Wundränder"]} & Gerötet / entzündlich \\ \hline
WUNDRAND\_GT\_MAPPING & \texttt{["Gerötet / entzündlich","unregelmäßige Wundränder"]} & Gerötet / entzündlich \\ \hline
WUNDRAND\_GT\_MAPPING & \texttt{["Gerötet / entzündlich","unregelmäßigen Rändern"]} & Gerötet / entzündlich \\ \hline
WUNDRAND\_GT\_MAPPING & \texttt{["Gerötet / entzündlich"]} & Gerötet / entzündlich \\ \hline
WUNDRAND\_GT\_MAPPING & \texttt{["Hyperkeratotisch","Gerötet / entzündlich"]} & Gerötet / entzündlich, Hyperkeratotisch \\ \hline
WUNDRAND\_GT\_MAPPING & \texttt{["Hyperkeratotisch"]} & Hyperkeratotisch \\ \hline
WUNDRAND\_GT\_MAPPING & \texttt{["Mazeriert","Epibolie (eingerollter Wundrand)","Gerötet / entzündlich"]} & Epibolie (eingerollter Wundrand), Gerötet / entzündlich, Mazeriert \\ \hline
WUNDRAND\_GT\_MAPPING & \texttt{["Mazeriert","Epibolie (eingerollter Wundrand)"]} & Epibolie (eingerollter Wundrand), Mazeriert \\ \hline
WUNDRAND\_GT\_MAPPING & \texttt{["Mazeriert","Gerötet / entzündlich","Epibolie (eingerollter Wundrand)"]} & Epibolie (eingerollter Wundrand), Gerötet / entzündlich, Mazeriert \\ \hline
WUNDRAND\_GT\_MAPPING & \texttt{["Mazeriert","Gerötet / entzündlich","Unterminiert (Wundtaschen)"]} & Gerötet / entzündlich, Mazeriert, Unterminiert (Wundtaschen) \\ \hline
WUNDRAND\_GT\_MAPPING & \texttt{["Mazeriert","Gerötet / entzündlich"]} & Gerötet / entzündlich, Mazeriert \\ \hline
WUNDRAND\_GT\_MAPPING & \texttt{["Mazeriert","Unterminiert (Wundtaschen)","Gerötet / entzündlich"]} & Gerötet / entzündlich, Mazeriert, Unterminiert (Wundtaschen) \\ \hline
WUNDRAND\_GT\_MAPPING & \texttt{["Mazeriert","Unterminiert (Wundtaschen)"]} & Mazeriert, Unterminiert (Wundtaschen) \\ \hline
WUNDRAND\_GT\_MAPPING & \texttt{["Mazeriert","teilweise aufgequollen"]} & Mazeriert \\ \hline
WUNDRAND\_GT\_MAPPING & \texttt{["Mazeriert"]} & Mazeriert \\ \hline
WUNDRAND\_GT\_MAPPING & \texttt{["Reizlos / unauffällig"]} & Reizlos / unauffällig \\ \hline
WUNDRAND\_GT\_MAPPING & \texttt{["Unterminiert (Wundtaschen)","Epibolie (eingerollter Wundrand)"]} & Epibolie (eingerollter Wundrand), Unterminiert (Wundtaschen) \\ \hline
WUNDRAND\_GT\_MAPPING & \texttt{["Unterminiert (Wundtaschen)","Gerötet / entzündlich"]} & Gerötet / entzündlich, Unterminiert (Wundtaschen) \\ \hline
WUNDRAND\_GT\_MAPPING & \texttt{["Unterminiert (Wundtaschen)","Mazeriert"]} & Mazeriert, Unterminiert (Wundtaschen) \\ \hline
WUNDRAND\_GT\_MAPPING & \texttt{["Unterminiert (Wundtaschen)"]} & Unterminiert (Wundtaschen) \\ \hline
WUNDRAND\_GT\_MAPPING & \texttt{["scharf begrenzt, teils livide/verfärbt, ischämietypisch"]} & Reizlos / unauffällig \\ \hline
WUNDTYP\_GT\_MAPPING & \texttt{Ausgedehnte feuchte Nekrose} & Ausgedehnte feuchte Nekrose \\ \hline
WUNDTYP\_GT\_MAPPING & \texttt{Ausgetrtenes Blut oder Lymphe aus den entsprechenden Gefäßen mit Verteilung an der betroffenen Extremität} & Ulkus (Ätiologie unspezifisch) \\ \hline
WUNDTYP\_GT\_MAPPING & \texttt{Dekubitus} & Dekubitus \\ \hline
WUNDTYP\_GT\_MAPPING & \texttt{Dekubitus Grad 2 oder 3 nach EPUAP.
Nicht genau zu definieren, ob sich Unterminierungen vorhanden sind.} & Dekubitus \\ \hline
WUNDTYP\_GT\_MAPPING & \texttt{Dekubitus an der Ferse} & Dekubitus \\ \hline
WUNDTYP\_GT\_MAPPING & \texttt{Dekubitus epuap Stadium 3 / 4} & Dekubitus \\ \hline
WUNDTYP\_GT\_MAPPING & \texttt{Dekubitus mit Nekrose} & Dekubitus \\ \hline
WUNDTYP\_GT\_MAPPING & \texttt{Dekubitus mit Taschenbildung} & Dekubitus \\ \hline
WUNDTYP\_GT\_MAPPING & \texttt{Dekubitus sacralbereich} & Dekubitus \\ \hline
WUNDTYP\_GT\_MAPPING & \texttt{Dekubitus, vermutlich am Außenknöchel} & Dekubitus \\ \hline
WUNDTYP\_GT\_MAPPING & \texttt{Dekubtis im Sakralbereich, großflächig} & Dekubitus \\ \hline
WUNDTYP\_GT\_MAPPING & \texttt{Fersendekubitus} & Dekubitus \\ \hline
WUNDTYP\_GT\_MAPPING & \texttt{Fersenulcus / eventuell Mischulcus} & Ulcus cruris mixtum \\ \hline
WUNDTYP\_GT\_MAPPING & \texttt{Fibrinbelegtes Ulcus} & Ulkus (Ätiologie unspezifisch) \\ \hline
WUNDTYP\_GT\_MAPPING & \texttt{Fibringbelegtes Ulcus} & Ulkus (Ätiologie unspezifisch) \\ \hline
WUNDTYP\_GT\_MAPPING & \texttt{Hypergranulation nach Verbrennung} & Verbrennungswunde \\ \hline
WUNDTYP\_GT\_MAPPING & \texttt{Hämangiom ist ein Spezialgebiet. In der Regel ist dieses innerhalb des ersten Lebensjahres rückläufig.
Deshalb nehme ich hier keine Beurteilung vor.} & Hämangiom ist ein Spezialgebiet. In der Regel ist dieses innerhalb des ersten Lebensjahres rückläufig.
Deshalb nehme ich hier keine Beurteilung vor. \\ \hline
WUNDTYP\_GT\_MAPPING & \texttt{Infizierte Wunden am linken Oberarm / Innenseitig} & Infizierte Wunden am linken Oberarm / Innenseitig \\ \hline
WUNDTYP\_GT\_MAPPING & \texttt{Infiziertes Ulcus mit freiliegender Sehne} & Ulkus (Ätiologie unspezifisch) \\ \hline
WUNDTYP\_GT\_MAPPING & \texttt{Oberflächliches Ulcus} & Ulkus (Ätiologie unspezifisch) \\ \hline
WUNDTYP\_GT\_MAPPING & \texttt{Oberflächliches Ulcus (diffuse Defekte)} & Ulkus (Ätiologie unspezifisch) \\ \hline
WUNDTYP\_GT\_MAPPING & \texttt{Plantares Ulcus linker Fuß, Ursachen Klärung notwendig, bei Diabetiker Prüfung ob Neuropathisch oder Angiopathische Ursache} & Plantares Ulcus linker Fuß, Ursachen Klärung notwendig, bei Diabetiker Prüfung ob Neuropathisch oder Angiopathische Ursache \\ \hline
WUNDTYP\_GT\_MAPPING & \texttt{Platzbauch nach möglichem chirurgischen Eingriff} & Postoperative Wunde / Dehiszenz \\ \hline
WUNDTYP\_GT\_MAPPING & \texttt{Postoperative Wunde abdominal} & Postoperative Wunde / Dehiszenz \\ \hline
WUNDTYP\_GT\_MAPPING & \texttt{Taschenbildendes Ulcus am linken Fußrücken} & Ulkus (Ätiologie unspezifisch) \\ \hline
WUNDTYP\_GT\_MAPPING & \texttt{Teils nekrotisches Ulcus an der linken Ferse.
plantares Ulcus belegt mit Infektionszeichen} & Ulkus (Ätiologie unspezifisch) \\ \hline
WUNDTYP\_GT\_MAPPING & \texttt{Ulcera rechter Fuß. könnte ursächlich diabetischer Fuß sein ( Neuropathie, Angiopathie) oder arterielles ulcus schäden auf Grund fehlender Durchblutung} & Ulkus (Ätiologie unspezifisch) \\ \hline
WUNDTYP\_GT\_MAPPING & \texttt{Ulcerationen bedingt durch Vaskulitis- Form der rheumatischen Erkrankung ( Autoimmunerkrankung)} & Ulkus (Ätiologie unspezifisch) \\ \hline
WUNDTYP\_GT\_MAPPING & \texttt{Ulcus} & Ulkus (Ätiologie unspezifisch) \\ \hline
WUNDTYP\_GT\_MAPPING & \texttt{Ulcus / warsch. Ulcus cruris, Ursache nicht klar, muss bestimmt werden} & Ulkus (Ätiologie unspezifisch) \\ \hline
WUNDTYP\_GT\_MAPPING & \texttt{Ulcus Dekubitus} & Dekubitus \\ \hline
WUNDTYP\_GT\_MAPPING & \texttt{Ulcus Dekubitus Bereich OS Sacrum} & Dekubitus \\ \hline
WUNDTYP\_GT\_MAPPING & \texttt{Ulcus Dekubitus Epuap Grad 3 / 4} & Dekubitus \\ \hline
WUNDTYP\_GT\_MAPPING & \texttt{Ulcus Dekubitus nach Epuap Grad 2 bis 3} & Dekubitus \\ \hline
WUNDTYP\_GT\_MAPPING & \texttt{Ulcus Ursache nicht genau definierbar} & Ulkus (Ätiologie unspezifisch) \\ \hline
WUNDTYP\_GT\_MAPPING & \texttt{Ulcus am Fussrücken} & Ulkus (Ätiologie unspezifisch) \\ \hline
WUNDTYP\_GT\_MAPPING & \texttt{Ulcus am Unterschenkel medial} & Ulkus (Ätiologie unspezifisch) \\ \hline
WUNDTYP\_GT\_MAPPING & \texttt{Ulcus an der Achillessehne} & Ulkus (Ätiologie unspezifisch) \\ \hline
WUNDTYP\_GT\_MAPPING & \texttt{Ulcus auf Grund von Mangeldurchblutung. Könnte diabetischer Fuß sein / angiopathisch bedingt.} & Ulkus (Ätiologie unspezifisch) \\ \hline
WUNDTYP\_GT\_MAPPING & \texttt{Ulcus auf Grund von Mangelernährung des Gewebes, Druckschäden in Kombination mit Gefäßschäden} & Ulkus (Ätiologie unspezifisch) \\ \hline
WUNDTYP\_GT\_MAPPING & \texttt{Ulcus cruris} & Ulkus (Ätiologie unspezifisch) \\ \hline
WUNDTYP\_GT\_MAPPING & \texttt{Ulcus cruris  unklarer Ursache} & Ulkus (Ätiologie unspezifisch) \\ \hline
WUNDTYP\_GT\_MAPPING & \texttt{Ulcus cruris /  Gravitationsulcus: Ulcus cruris venosum} & Ulcus cruris venosum \\ \hline
WUNDTYP\_GT\_MAPPING & \texttt{Ulcus cruris li lateraler Malleolus} & Ulkus (Ätiologie unspezifisch) \\ \hline
WUNDTYP\_GT\_MAPPING & \texttt{Ulcus cruris möglich arteriell oder mixtum} & Ulkus (Ätiologie unspezifisch) \\ \hline
WUNDTYP\_GT\_MAPPING & \texttt{Ulcus cruris venosum} & Ulcus cruris venosum \\ \hline
WUNDTYP\_GT\_MAPPING & \texttt{Ulcus decubitus} & Dekubitus \\ \hline
WUNDTYP\_GT\_MAPPING & \texttt{Ulcus decubitus Nach Epuap Grad 2 bis 3} & Dekubitus \\ \hline
WUNDTYP\_GT\_MAPPING & \texttt{Ulcus dekubitus} & Dekubitus \\ \hline
WUNDTYP\_GT\_MAPPING & \texttt{Ulcus dekubitus Ferse,  Stadium 3 / 4} & Dekubitus \\ \hline
WUNDTYP\_GT\_MAPPING & \texttt{Ulcus lokal abgegrenzt. belegt könnte ebenso nekrotisch sein} & Ulkus (Ätiologie unspezifisch) \\ \hline
WUNDTYP\_GT\_MAPPING & \texttt{Ulcus mit Stauung} & Ulkus (Ätiologie unspezifisch) \\ \hline
WUNDTYP\_GT\_MAPPING & \texttt{Ulcus mit scheinender Hypergranulation oberhalb der Achillessehne} & Ulkus (Ätiologie unspezifisch) \\ \hline
WUNDTYP\_GT\_MAPPING & \texttt{Ulcus nach Meshgraft (zumindest lässt die Struktur darauf schließen), eventuell ist die Struktur aber einer Unterdrucktherapie geschuldet} & Postoperative Wunde / Dehiszenz \\ \hline
WUNDTYP\_GT\_MAPPING & \texttt{Ulcus nach Stich} & Traumatische Wunde \\ \hline
WUNDTYP\_GT\_MAPPING & \texttt{Ulcus unbekannter Herkunft, könnte traumatologisch bedingte Wunde sein, durch die Wunde bei unsachgemäßer Versorgung entstehen ulcerationen ebenfalls, nicht immer Krankheitsbedingt} & Ulkus (Ätiologie unspezifisch) \\ \hline
WUNDTYP\_GT\_MAPPING & \texttt{Ulcus unklarer Genese} & Ulkus (Ätiologie unspezifisch) \\ \hline
WUNDTYP\_GT\_MAPPING & \texttt{Ulcus unklarer Genese am Bein genaue Lokalisation nicht beurteilbar} & Ulkus (Ätiologie unspezifisch) \\ \hline
WUNDTYP\_GT\_MAPPING & \texttt{Ulcus unklarer Genese, könnte Venös oder arteriell oder mixum sein.
Beachte kleine Nekrose an der Ferse plantar} & Ulkus (Ätiologie unspezifisch) \\ \hline
WUNDTYP\_GT\_MAPPING & \texttt{Ulcus unklarer Genese, könnte diabet. Ulcus sein, möglich ebenso durch traumatologische Ursachen} & Ulkus (Ätiologie unspezifisch) \\ \hline
WUNDTYP\_GT\_MAPPING & \texttt{Ulcus, Lokalisation am Bein nicht genau definierbar} & Ulkus (Ätiologie unspezifisch) \\ \hline
WUNDTYP\_GT\_MAPPING & \texttt{Ulcus, bei vermutetem Diabetes} & Ulkus (Ätiologie unspezifisch) \\ \hline
WUNDTYP\_GT\_MAPPING & \texttt{Ulcus, durch die Mangelernährung ( keine Durchblutung) im Wundgebiet  könnte es sich um einen Dekubitus oder ulcus cruris handeln.} & Ulkus (Ätiologie unspezifisch) \\ \hline
WUNDTYP\_GT\_MAPPING & \texttt{Ulcus. Ursache klären sollte es ein diabet. Fuß sein dann Prüfung ob angiopathisch oder neuropathisch,
Druckentlastung und Ursachenbeseitigung.} & Ulkus (Ätiologie unspezifisch) \\ \hline
WUNDTYP\_GT\_MAPPING & \texttt{Ulkus Dekubitus am Os sacrum  nach Epuap Grad 2 / 3} & Dekubitus \\ \hline
WUNDTYP\_GT\_MAPPING & \texttt{Ulkus an der Ferse, warsch. Dekubitus nach EPUAP  Stadium 3} & Dekubitus \\ \hline
WUNDTYP\_GT\_MAPPING & \texttt{Verbrennung 2 und 3 Grades nach 9 er Regel etwa 9 \% Körperoberfläche} & Verbrennungswunde \\ \hline
WUNDTYP\_GT\_MAPPING & \texttt{Verbrennung 2. Grades am Bein} & Verbrennungswunde \\ \hline
WUNDTYP\_GT\_MAPPING & \texttt{Verbrennung 2. und teilweise 3. Grades Oberfläche ca 9 \%} & Verbrennungswunde \\ \hline
WUNDTYP\_GT\_MAPPING & \texttt{Verbrennung Grad 1 bis 2 am rechten Fuß} & Verbrennungswunde \\ \hline
WUNDTYP\_GT\_MAPPING & \texttt{Verbrennungswunde 1. bis 2. Grades} & Verbrennungswunde \\ \hline
WUNDTYP\_GT\_MAPPING & \texttt{Verbrennungswunde am Arm Grad 2} & Verbrennungswunde \\ \hline
WUNDTYP\_GT\_MAPPING & \texttt{Wunde am li Fußrücken mit Sehnenfreilegung} & Wunde am li Fußrücken mit Sehnenfreilegung \\ \hline
WUNDTYP\_GT\_MAPPING & \texttt{Wunde aufgrund einer Infusion, welche ins umliegende Gewebe ausgetreten ist} & Wunde aufgrund einer Infusion, welche ins umliegende Gewebe ausgetreten ist \\ \hline
WUNDTYP\_GT\_MAPPING & \texttt{["Dekubitus","Diabetisches Fußulkus"]} & Ulkus (Ätiologie unspezifisch) \\ \hline
WUNDTYP\_GT\_MAPPING & \texttt{["Dekubitus","Ulcus"]} & Ulkus (Ätiologie unspezifisch) \\ \hline
WUNDTYP\_GT\_MAPPING & \texttt{["Dekubitus"]} & Dekubitus \\ \hline
WUNDTYP\_GT\_MAPPING & \texttt{["Diabetisches Fußulkus","hochgradig destruierende Fußwunde"]} & Diabetisches Fußsyndrom (DFS) \\ \hline
WUNDTYP\_GT\_MAPPING & \texttt{["Diabetisches Fußulkus","ischämisches Ulkus"]} & Diabetisches Fußsyndrom (DFS) \\ \hline
WUNDTYP\_GT\_MAPPING & \texttt{["Diabetisches Fußulkus","mit deutlichen Zeichen einer lokalen Infektion und chronisch entzündlicher Hautveränderungen"]} & Diabetisches Fußsyndrom (DFS) \\ \hline
WUNDTYP\_GT\_MAPPING & \texttt{["Diabetisches Fußulkus"]} & Diabetisches Fußsyndrom (DFS) \\ \hline
WUNDTYP\_GT\_MAPPING & \texttt{["Extravasationsverletzung"]} & Extravasationsverletzung \\ \hline
WUNDTYP\_GT\_MAPPING & \texttt{["Gravitationsulzera/Druckulzera"]} & Ulkus (Ätiologie unspezifisch) \\ \hline
WUNDTYP\_GT\_MAPPING & \texttt{["Postoperative Wunde"]} & Postoperative Wunde / Dehiszenz \\ \hline
WUNDTYP\_GT\_MAPPING & \texttt{["Traumatische Wunde","durch Insektenstich"]} & Traumatische Wunde \\ \hline
WUNDTYP\_GT\_MAPPING & \texttt{["Ulcus Fußrücken mit deutlicher Gewebedestruktion und freiliegender Sehnenstruktur"]} & Ulkus (Ätiologie unspezifisch) \\ \hline
WUNDTYP\_GT\_MAPPING & \texttt{["Ulcus cruris venosum","Adipositas-assoziiert"]} & Ulcus cruris venosum \\ \hline
WUNDTYP\_GT\_MAPPING & \texttt{["Ulcus cruris venosum","venös-lymphatischem Ulcus"]} & Ulcus cruris venosum \\ \hline
WUNDTYP\_GT\_MAPPING & \texttt{["Ulcus cruris venosum"]} & Ulcus cruris venosum \\ \hline
WUNDTYP\_GT\_MAPPING & \texttt{["Ulcus"]} & Ulkus (Ätiologie unspezifisch) \\ \hline
WUNDTYP\_GT\_MAPPING & \texttt{["Ulkus cruris"]} & Ulkus (Ätiologie unspezifisch) \\ \hline
WUNDTYP\_GT\_MAPPING & \texttt{["Ulkus"]} & Ulkus (Ätiologie unspezifisch) \\ \hline
WUNDTYP\_GT\_MAPPING & \texttt{["Verbrennungswunde"]} & Verbrennungswunde \\ \hline
WUNDTYP\_GT\_MAPPING & \texttt{["ausgedehnte, tiefe, nekrotische Ulzeration"]} & Ulkus (Ätiologie unspezifisch) \\ \hline
WUNDTYP\_GT\_MAPPING & \texttt{["ausgedehntes, tiefes chronisches Ulcus"]} & Ulkus (Ätiologie unspezifisch) \\ \hline
WUNDTYP\_GT\_MAPPING & \texttt{["ausgeprägtes infiziertes Ulcus cruris","Ulcus cruris venosum"]} & Ulkus (Ätiologie unspezifisch) \\ \hline
WUNDTYP\_GT\_MAPPING & \texttt{["chronische Fußwunde / Ulkus"]} & Ulkus (Ätiologie unspezifisch) \\ \hline
WUNDTYP\_GT\_MAPPING & \texttt{["chronisches Ulkus cruris"]} & Ulkus (Ätiologie unspezifisch) \\ \hline
WUNDTYP\_GT\_MAPPING & \texttt{["chronisches Ulkus"]} & Ulkus (Ätiologie unspezifisch) \\ \hline
WUNDTYP\_GT\_MAPPING & \texttt{["großflächiges, flaches chronisches Ulkus"]} & Ulkus (Ätiologie unspezifisch) \\ \hline
WUNDTYP\_GT\_MAPPING & \texttt{["kleines, oberflächliches Ulcus"]} & Ulkus (Ätiologie unspezifisch) \\ \hline
WUNDTYP\_GT\_MAPPING & \texttt{["multiple infizierte chronische Fußulzera"]} & Ulkus (Ätiologie unspezifisch) \\ \hline
WUNDTYP\_GT\_MAPPING & \texttt{["multiple kleine Ulzera"]} & Ulkus (Ätiologie unspezifisch) \\ \hline
WUNDTYP\_GT\_MAPPING & \texttt{["nekrotische Fersenwunde","Dekubitus"]} & Ulkus (Ätiologie unspezifisch) \\ \hline
WUNDTYP\_GT\_MAPPING & \texttt{["nekrotische ischämische Fußwunde / Gangrän"]} & Ulcus cruris arteriosum / Ischämisches Ulkus \\ \hline
WUNDTYP\_GT\_MAPPING & \texttt{["neuroischämisches bzw. arterielles Fußulkus"]} & Ulcus cruris arteriosum / Ischämisches Ulkus \\ \hline
WUNDTYP\_GT\_MAPPING & \texttt{["rundlich-ovales Ulcus"]} & Ulkus (Ätiologie unspezifisch) \\ \hline
WUNDTYP\_GT\_MAPPING & \texttt{["tiefer reichende Ulzerationen am Unterschenkel mit chronisch entzündlicher Umgebung"]} & Ulkus (Ätiologie unspezifisch) \\ \hline
WUNDTYP\_GT\_MAPPING & \texttt{["ulzeriertes infantiles Hämangiom mit zentraler Nekrose"]} & ulzeriertes infantiles Hämangiom mit zentraler Nekrose \\ \hline
WUNDTYP\_GT\_MAPPING & \texttt{["vaskulitischen Ulzera","möglicher kutaner Vaskulitis"]} & Ulkus (Ätiologie unspezifisch) \\ \hline
WUNDTYP\_GT\_MAPPING & \texttt{arterielles Ulcus} & Ulcus cruris arteriosum / Ischämisches Ulkus \\ \hline
WUNDTYP\_GT\_MAPPING & \texttt{belegter Ulcus} & Ulkus (Ätiologie unspezifisch) \\ \hline
WUNDTYP\_GT\_MAPPING & \texttt{diffuse fibrinbelegte und teils infizierte Ulzera am rechten Fuß} & Ulkus (Ätiologie unspezifisch) \\ \hline
WUNDTYP\_GT\_MAPPING & \texttt{fibrinbelegtes Ulcus} & Ulkus (Ätiologie unspezifisch) \\ \hline
WUNDTYP\_GT\_MAPPING & \texttt{gleiche Beurteilung wie Wunde 53} & Dekubitus \\ \hline
WUNDTYP\_GT\_MAPPING & \texttt{infantiles Hämangiom,} & infantiles Hämangiom, \\ \hline
WUNDTYP\_GT\_MAPPING & \texttt{könnte arterieller Ulcus oder diabetischer Ulcus ( Angiopathie) auf Grund der Mangeldurchblutung  sein.} & Ulkus (Ätiologie unspezifisch) \\ \hline
WUNDTYP\_GT\_MAPPING & \texttt{könnte diabetischer ulcus sein auf grund einer neuropathie} & Ulkus (Ätiologie unspezifisch) \\ \hline
WUNDTYP\_GT\_MAPPING & \texttt{könnte ein diabetischer Fuß sein, Ulcus  nicht möglich,
traumatologische Indikation ebenso nicht ausgeschlossen} & Ulkus (Ätiologie unspezifisch) \\ \hline
WUNDTYP\_GT\_MAPPING & \texttt{laterales Ulcus linker Fuß} & Ulkus (Ätiologie unspezifisch) \\ \hline
WUNDTYP\_GT\_MAPPING & \texttt{massives Ulcus dekubitus} & Dekubitus \\ \hline
WUNDTYP\_GT\_MAPPING & \texttt{mumifizierte Zehen (D4 und D5), gleichzeitig Nekrose an der rechten Ferse} & Ulcus cruris arteriosum / Ischämisches Ulkus \\ \hline
WUNDTYP\_GT\_MAPPING & \texttt{möglicherweise ein feuchtes Gangrän} & Ulcus cruris arteriosum / Ischämisches Ulkus \\ \hline
WUNDTYP\_GT\_MAPPING & \texttt{nekrotischen Wunden am Fuß} & nekrotischen Wunden am Fuß \\ \hline
WUNDTYP\_GT\_MAPPING & \texttt{nekrotischer Defekt / Dekubitus} & Dekubitus \\ \hline
WUNDTYP\_GT\_MAPPING & \texttt{nekrotischer Dekubitus li Ferse} & Dekubitus \\ \hline
WUNDTYP\_GT\_MAPPING & \texttt{oberflächlicher Dekubitus ( 2 Stück)} & Dekubitus \\ \hline
WUNDTYP\_GT\_MAPPING & \texttt{oberflächliches Ulcus} & Ulkus (Ätiologie unspezifisch) \\ \hline
WUNDTYP\_GT\_MAPPING & \texttt{offene ulcera} & Ulkus (Ätiologie unspezifisch) \\ \hline
WUNDTYP\_GT\_MAPPING & \texttt{plantares ulcus} & Ulkus (Ätiologie unspezifisch) \\ \hline
WUNDTYP\_GT\_MAPPING & \texttt{pseudomas besiedeltes teils nekrotisches Ulcus.
Semizirkulär} & Ulkus (Ätiologie unspezifisch) \\ \hline
WUNDTYP\_GT\_MAPPING & \texttt{rechter Fuß lateral/ malleolär} & rechter Fuß lateral/ malleolär \\ \hline
WUNDTYP\_GT\_MAPPING & \texttt{schwere Nekrose am Fuß rechts lateral und plantar, offene Geschwüre über den Malleolen lateral
deutet auf arterielles Ulcus hin} & Ulcus cruris arteriosum / Ischämisches Ulkus \\ \hline
WUNDTYP\_GT\_MAPPING & \texttt{schwere Ulzerationen am Fuß- mehrere offene Stellen, freiliegende Sehen, teilweise nekrotisch, belegt} & Ulkus (Ätiologie unspezifisch) \\ \hline
WUNDTYP\_GT\_MAPPING & \texttt{semizirkuläres Ulcus} & Ulkus (Ätiologie unspezifisch) \\ \hline
WUNDTYP\_GT\_MAPPING & \texttt{semizirkuläres Ulcus re Fuß} & Ulkus (Ätiologie unspezifisch) \\ \hline
WUNDTYP\_GT\_MAPPING & \texttt{superinfizierte Ulcration am rechen Fuß offene Ulcera} & Ulkus (Ätiologie unspezifisch) \\ \hline
WUNDTYP\_GT\_MAPPING & \texttt{tiefer Dekubitus} & Dekubitus \\ \hline
WUNDTYP\_GT\_MAPPING & \texttt{tiefes Ulcus im Fersenbereich} & Ulkus (Ätiologie unspezifisch) \\ \hline
WUNDTYP\_GT\_MAPPING & \texttt{traumatologische Wunde ( thermische Schädigung / Verbrennung 3. und 4. Grades} & Verbrennungswunde \\ \hline
WUNDTYP\_GT\_MAPPING & \texttt{ulcus cruris} & Ulkus (Ätiologie unspezifisch) \\ \hline
WUNDTYP\_GT\_MAPPING & \texttt{vermutlich Mischulcus am rechten Bein, höhe der Knöchel.
Semizirkulär} & Ulcus cruris mixtum \\ \hline
WUNDTYP\_GT\_MAPPING & \texttt{zirkuläres Ulcus Unterschenkel (vermutlich Mischulcus)} & Ulcus cruris mixtum \\ \hline
WUNDUMGEBUNG\_GT\_MAPPING & \texttt{["CVI-typische Hautveränderungen (Hyperpigmentierung, Atrophie blanche, Lipodermatosklerose)","Erythem / Rötung"]} & CVI-typische Hautveränderungen (Hyperpigmentierung, Atrophie blanche, Lipodermatosklerose), Erythem / Rötung \\ \hline
WUNDUMGEBUNG\_GT\_MAPPING & \texttt{["CVI-typische Hautveränderungen (Hyperpigmentierung, Atrophie blanche, Lipodermatosklerose)"]} & CVI-typische Hautveränderungen (Hyperpigmentierung, Atrophie blanche, Lipodermatosklerose) \\ \hline
WUNDUMGEBUNG\_GT\_MAPPING & \texttt{["Ekzem / Dermatitis","Erythem / Rötung"]} & Ekzem / Dermatitis, Erythem / Rötung \\ \hline
WUNDUMGEBUNG\_GT\_MAPPING & \texttt{["Ekzem / Dermatitis"]} & Ekzem / Dermatitis \\ \hline
WUNDUMGEBUNG\_GT\_MAPPING & \texttt{["Erythem / Rötung","CVI-typische Hautveränderungen (Hyperpigmentierung, Atrophie blanche, Lipodermatosklerose)"]} & CVI-typische Hautveränderungen (Hyperpigmentierung, Atrophie blanche, Lipodermatosklerose), Erythem / Rötung \\ \hline
WUNDUMGEBUNG\_GT\_MAPPING & \texttt{["Erythem / Rötung","Ekzem / Dermatitis","CVI-typische Hautveränderungen (Hyperpigmentierung, Atrophie blanche, Lipodermatosklerose)"]} & CVI-typische Hautveränderungen (Hyperpigmentierung, Atrophie blanche, Lipodermatosklerose), Ekzem / Dermatitis, Erythem / Rötung \\ \hline
WUNDUMGEBUNG\_GT\_MAPPING & \texttt{["Erythem / Rötung","Ekzem / Dermatitis"]} & Ekzem / Dermatitis, Erythem / Rötung \\ \hline
WUNDUMGEBUNG\_GT\_MAPPING & \texttt{["Erythem / Rötung","Mazeration","CVI-typische Hautveränderungen (Hyperpigmentierung, Atrophie blanche, Lipodermatosklerose)"]} & CVI-typische Hautveränderungen (Hyperpigmentierung, Atrophie blanche, Lipodermatosklerose), Erythem / Rötung, Mazeration \\ \hline
WUNDUMGEBUNG\_GT\_MAPPING & \texttt{["Erythem / Rötung","Mazeration"]} & Erythem / Rötung, Mazeration \\ \hline
WUNDUMGEBUNG\_GT\_MAPPING & \texttt{["Erythem / Rötung","Weichteilreizung"]} & Erythem / Rötung \\ \hline
WUNDUMGEBUNG\_GT\_MAPPING & \texttt{["Erythem / Rötung","deutliche entzündliche Rötung der Umgebung"]} & Erythem / Rötung \\ \hline
WUNDUMGEBUNG\_GT\_MAPPING & \texttt{["Erythem / Rötung","eher trocken wirkende Umgebung"]} & Erythem / Rötung \\ \hline
WUNDUMGEBUNG\_GT\_MAPPING & \texttt{["Erythem / Rötung","entzündliche Umgebungshaut"]} & Erythem / Rötung \\ \hline
WUNDUMGEBUNG\_GT\_MAPPING & \texttt{["Erythem / Rötung","leicht gerötete, glänzende Umgebungshaut"]} & Erythem / Rötung \\ \hline
WUNDUMGEBUNG\_GT\_MAPPING & \texttt{["Erythem / Rötung","Ödem","CVI-typische Hautveränderungen (Hyperpigmentierung, Atrophie blanche, Lipodermatosklerose)"]} & CVI-typische Hautveränderungen (Hyperpigmentierung, Atrophie blanche, Lipodermatosklerose), Erythem / Rötung, Ödem \\ \hline
WUNDUMGEBUNG\_GT\_MAPPING & \texttt{["Erythem / Rötung","Ödem"]} & Erythem / Rötung, Ödem \\ \hline
WUNDUMGEBUNG\_GT\_MAPPING & \texttt{["Erythem / Rötung"]} & Erythem / Rötung \\ \hline
WUNDUMGEBUNG\_GT\_MAPPING & \texttt{["Keratosen"]} & Sonstiges \\ \hline
WUNDUMGEBUNG\_GT\_MAPPING & \texttt{["Mazeration","Erythem / Rötung","trockener wirkende Areale mit verminderter Durchblutung"]} & Erythem / Rötung, Mazeration \\ \hline
WUNDUMGEBUNG\_GT\_MAPPING & \texttt{["Mazeration","Erythem / Rötung"]} & Erythem / Rötung, Mazeration \\ \hline
WUNDUMGEBUNG\_GT\_MAPPING & \texttt{["Mazeration","glänzend"]} & Mazeration \\ \hline
WUNDUMGEBUNG\_GT\_MAPPING & \texttt{["Mazeration","glänzende, gespannte Umgebungshaut"]} & Mazeration \\ \hline
WUNDUMGEBUNG\_GT\_MAPPING & \texttt{["Mazeration"]} & Mazeration \\ \hline
WUNDUMGEBUNG\_GT\_MAPPING & \texttt{["Reizlos / intakt","Erythem / Rötung"]} & Erythem / Rötung, Reizlos / intakt \\ \hline
WUNDUMGEBUNG\_GT\_MAPPING & \texttt{["Reizlos / intakt"]} & Reizlos / intakt \\ \hline
WUNDUMGEBUNG\_GT\_MAPPING & \texttt{["atrophisch-trockene Haut mit trophischen Störungen","CVI-typische Hautveränderungen (Hyperpigmentierung, Atrophie blanche, Lipodermatosklerose)"]} & CVI-typische Hautveränderungen (Hyperpigmentierung, Atrophie blanche, Lipodermatosklerose) \\ \hline
WUNDUMGEBUNG\_GT\_MAPPING & \texttt{["empfindliche Säuglingshaut mit hoher Feuchtigkeits- und Reibungsbelastung"]} & Sonstiges \\ \hline
WUNDUMGEBUNG\_GT\_MAPPING & \texttt{["fragile, atrophe Haut"]} & Sonstiges \\ \hline
WUNDUMGEBUNG\_GT\_MAPPING & \texttt{["gerötet-glänzende Haut, Hinweis auf chronische venöse Stauung/Ödemneigung"]} & Erythem / Rötung, Ödem \\ \hline
WUNDUMGEBUNG\_GT\_MAPPING & \texttt{["gerötet-glänzende, ödematöse Haut","Erythem / Rötung"]} & Erythem / Rötung, Ödem \\ \hline
WUNDUMGEBUNG\_GT\_MAPPING & \texttt{["gut durchbluteter, überwiegend roter Wundgrund"]} & Reizlos / intakt \\ \hline
WUNDUMGEBUNG\_GT\_MAPPING & \texttt{["rocken, teils schuppig, gespannt"]} & Sonstiges \\ \hline
WUNDUMGEBUNG\_GT\_MAPPING & \texttt{["zahlreiche erythematöse Makulae/Papeln, livid-entzündliche Hautveränderungen, Hinweis auf entzündlich-vaskulären Prozess"]} & Ekzem / Dermatitis, Erythem / Rötung \\ \hline
WUNDUMGEBUNG\_GT\_MAPPING & \texttt{["Ödem","CVI-typische Hautveränderungen (Hyperpigmentierung, Atrophie blanche, Lipodermatosklerose)"]} & CVI-typische Hautveränderungen (Hyperpigmentierung, Atrophie blanche, Lipodermatosklerose), Ödem \\ \hline
WUNDUMGEBUNG\_GT\_MAPPING & \texttt{["Ödem","Erythem / Rötung"]} & Erythem / Rötung, Ödem \\ \hline
WUNDUMGEBUNG\_GT\_MAPPING & \texttt{["Ödem","gespannte, trophisch veränderte Haut"]} & Ödem \\ \hline
WUNDUMGEBUNG\_GT\_MAPPING & \texttt{["Ödem"]} & Ödem \\ \hline
WUNDUMGEBUNG\_GT\_MAPPING & \texttt{[]} & nan \\ \hline
\end{longtable}